% DER_III-V: Grothendieck's text with the commentary set in, then the overview, inside a group that applies the chunk's own definitions
\dgrchapter{III}{Chapter III}{Hom externes dans les dérivateurs}
\dgrpart{0}{III-V}{Manuscript pages 1--12}{III-V\enspace\textperiodcentered\enspace Chapter III\enspace\textperiodcentered\enspace manuscript pages 1--12}{Chapter III \textperiodcentered\ \textup{III-V}}
\dgrchunkbegin
% DER_III-V: the definitions of the chunk's preamble, applied inside its group
\let\printedpage\relax\newcommand{\printedpage}[1]{\marginpar{\footnotesize p.~#1}}
\let\msp\relax\newcommand{\msp}[1]{\marginpar{\footnotesize [#1]}}
\let\Hot\relax\newcommand{\Hot}{\mathrm{Hot}}
\let\D\relax\newcommand{\D}{\mathbf{D}}
\let\uHom\relax\newcommand{\uHom}{\underline{\mathrm{Hom}}}
\let\hm\relax\newcommand{\hm}{\mathrm{hom}}
\let\Hom\relax\newcommand{\Hom}{\mathrm{Hom}}
\let\homext\relax\newcommand{\homext}{\mathrm{hom}\,\mathrm{ext}}
\let\HOT\relax\newcommand{\HOT}{\mathrm{HOT}}
\let\Ob\relax\newcommand{\Ob}{\mathrm{Ob}\,}
\let\Ltens\relax\newcommand{\Ltens}{\overset{\mathbf{L}}{\otimes}}
\let\RHom\relax\newcommand{\RHom}{\mathbf{R}\mathrm{Hom}}
\let\ab\relax\newcommand{\ab}{\mathrm{ab}}
\let\defeq\relax\newcommand{\defeq}{\overset{\text{déf}}{=}}
\let\cA\relax\newcommand{\cA}{\mathcal{A}}
\let\bDelta\relax\newcommand{\bDelta}{\boldsymbol{\Delta}}
\let\id\relax\newcommand{\id}{\mathrm{id}}
\let\Fl\relax\newcommand{\Fl}{\operatorname{Fl}}
\let\uCh\relax\newcommand{\uCh}{\underline{\mathrm{Ch}}}
\let\bPhi\relax\newcommand{\bPhi}{\boldsymbol{\Phi}}
\let\uFl\relax\newcommand{\uFl}{\underline{\mathrm{Fl}}}
\let\Ch\relax\newcommand{\Ch}{\underline{\mathrm{Ch}}}
\let\card\relax\newcommand{\card}{\mathrm{card}}
\let\cM\relax\newcommand{\cM}{\mathcal{M}}
\let\Chh\relax\newcommand{\Chh}{\underline{\underline{\mathrm{Ch}}}}
\let\KK\relax\newcommand{\KK}{\underline{\underline{K}}}
\let\dfeq\relax\newcommand{\dfeq}{\overset{\text{déf}}{=}}
\let\Cathot\relax\newcommand{\Cathot}{\mathrm{Cat\,hot}}
\let\Ind\relax\newcommand{\Ind}{\operatorname{Ind}}
\let\Ens\relax\newcommand{\Ens}{\mathrm{Ens}}
\let\Cat\relax\newcommand{\Cat}{\mathrm{Cat}}
\let\uuK\relax\newcommand{\uuK}{\underline{\underline{K}}}

% ---- from Der-ChIII-567_p001-005.tex
% ===================== page 1 (title page) =====================
\printedpage{--}
% no page number printed on this page
% head: (none)

\begin{center}
{\LARGE\sffamily LES DÉRIVATEURS}

\bigskip

\textbf{ALEXANDRE GROTHENDIECK}

\bigskip\bigskip

\textbf{Chapitre III}

\medskip

\textbf{Hom externes dans les dérivateurs}
\end{center}

\clearpage

% ===================== page 2 (editorial note) =====================
\printedpage{--}
% no page number printed on this page
% head: (none)

Ce texte a été déchiffré et transcrit en \LaTeXe{} par M. Künzer. Il a été édité par M. Künzer, J. Malgoire, et G. Maltsiniotis. La transcription est aussi fidèle que possible au manuscrit. Pour les quelques corrections évidentes, ou rares commentaires des éditeurs, ainsi que pour la numérotation originale des pages du manuscrit, les caractères de machine à écrire \texttt{[typewriter]} entre crochets sont utilisés. Un point d'interrogation entre crochets signifie que l'on n'est pas sûr du mot qui précède.

Cette édition est provisoire. Les remarques, commentaires et corrections sont bienvenus. Envoyer un message à :

\textbf{maltsin@math.jussieu.fr}
\begin{flushright}
G. Maltsiniotis\qquad\qquad
\end{flushright}

\clearpage

% ===================== page 3 (printed 1) =====================
\printedpage{1}
% head: A. Grothendieck : Dérivateurs, Ch. III, Édition des Universités de Montpellier II et Paris VII
% foot: Version 2 février 2002

\msp{page 1}

\section*{1\quad Relation des dérivateurs avec la catégorie $\Hot$}

Soit $\D$ un dérivateur, $\mathcal{A}$ sa catégorie fondamentale,
\[
\mathcal{A} = \D(e).
\]
Je me propose de définir, si faire se peut, des foncteurs
\[
\left\{
\begin{array}{rccccccl}
\Hot & \times & \mathcal{A} & \xrightarrow{\pi_{\D}} & \mathcal{A} & (\xi, M) & \longmapsto & \xi \times M \ \text{ ou } \ \xi \otimes M \quad \text{(suivant les cas)}\\
\Hot^o & \times & \mathcal{A} & \xrightarrow{\varpi_{\D}} & \mathcal{A} & (\xi, M) & \longmapsto & \hm(\xi, M).
\end{array}
\right.
\tag*{(1)}
\]
Ces foncteurs doivent s'échanger par dualité : si on prend les foncteurs $\pi^o$, $\varpi^o$ sur les catégories opposées, on trouve
\[
\left\{
\begin{array}{rcccc}
\Hot & \times & \mathcal{A}^o & \xrightarrow{\varpi^o_{\D}} & \mathcal{A}^o\\
\Hot^o & \times & \mathcal{A}^o & \xrightarrow{\pi^o_{\D}} & \mathcal{A}^o,
\end{array}
\right.
\tag*{(1$'$)}
\]
et on doit avoir
\[
\left\{
\begin{array}{rcl}
\varpi^o_{\D} & = & \pi_{\D^o}\\
\pi^o_{\D} & = & \varpi_{\D^o}.
\end{array}
\right.
\]
À vrai dire, considérant (1) comme des foncteurs
\[
\begin{array}{rcl}
\Hot & \xrightarrow{\pi_{\D}} & \uHom(\mathcal{A}, \mathcal{A})\\
\Hot^o & \xrightarrow{\varpi_{\D}} & \uHom(\mathcal{A}, \mathcal{A}),
\end{array}
\]
on s'attend à ce que ces foncteurs s'étendent en
\[
\left\{
\begin{array}{rcl}
\Hot & \xrightarrow{\pi_{\D}} & \uHom(\D, \D)\\
\Hot^o & \xrightarrow{\varpi_{\D}} & \uHom(\mathcal{A}, \mathcal{A})
\end{array}
\right.
\tag*{(2)}
\]
\texttt{[plutôt deux fois} $\uHom(\D, \D)$\texttt{]}. Cela signifie, essentiellement, que pour $\xi \in \mathrm{Ob}\,\Hot$ donné, les opérations $\xi \times_{\D} M$ et $\uHom_{\D}(\xi, M)$ \texttt{[plutôt} $\hm_{\D}(\xi, M)$\texttt{]} sont définies non seulement dans $\mathcal{A} = \D(e)$, mais dans tous les $\D(I)$, et ceci de façon à commuter aux foncteurs $f^*$ :

\dgrblock{dgrb147}{Spaces acting on coefficients}
\begin{dcomment}{exposition}{p1:Ces foncteurs doivent s'échanger par dualité}{Spaces acting on coefficients}
Grothendieck is looking for the part of ordinary homotopy theory that must act on every theory of coefficients. A space should provide an instruction for assembling copies of an object, while the same space, with the variance reversed, should provide an instruction for taking their limit. This is the reason for placing the two operations together before attempting to construct an external Hom.

Write \(B X=|N X|\) for the classifying space of a small category, and temporarily assume that all the categories used below belong to the permitted diagram category. If \(p_X:X\to e\), put
\[
 T_X(M)=(p_X)_!p_X^*M,\qquad
 C_X(M)=(p_X)_*p_X^*M.
\]
These expressions are already defined by the derivator. The problem is to show that they depend on \(X\) only through \(B X\), and then to understand their compatibility as \(X\) and \(M\) vary. In the derivator of spaces they are \(B X\times M\) and the derived function space \(\operatorname{Map}(B X,M)\). In a derived category they are derived tensor and Hom constructions with the integral chains of \(B X\). The common formalism therefore sees more than the set of components of \(X\).

There are three different Hom constructions in play. The ordinary set \(\Hom_{\D(e)}(M,N)\) records morphisms after localization. The cotensor \(C_X(N)\) is another coefficient object, hence belongs to \(\D(e)\). The proposed external Hom \(\homext(M,N)\) is a homotopy type. Its intended property is
\[
 \Hom_{\D(e)}(T_XM,N)
 \cong [B X,\homext(M,N)].
\]
In a model presentation this last object is the derived mapping space. It is not the internal Hom of a monoidal category of coefficients: no multiplication of two coefficient objects has been assumed. For example, the derived category of modules over a noncommutative ring still has an external mapping space, although a tensor product of two left modules over that ring has no such general meaning.

The extension to every \(\D(I)\) is essential. With \(p:X\times I\to I\), the same formulas define \(T_{X,I}=p_!p^*\) and \(C_{X,I}=p_*p^*\). Base change along products supplies their compatibility with restriction \(f^*\). Thus the request in (2) is for morphisms of derivators, including their pseudonaturality isomorphisms, rather than unrelated endofunctors of the individual categories. The passage from \(\D(e)\) to the shifted derivator \(\D_I\) is the mechanism already used in the commentary on Chapter~I, manuscript p.~4 (\emph{From derived coefficients to diagram categories}); here its new role is to make the action itself coherent in the diagram variable.

For a represented derivator, the calculation reduces to an instructive limiting case:
\[
 T_X(M)=\coprod_{\pi_0 X}M,\qquad
 C_X(M)=\prod_{\pi_0 X}M.
\]
Indeed a cone on a constant ordinary diagram is a family of maps constant on every connected component. The richer dependence on \(B X\) belongs to the homotopy theory, not to the notation \(p_!p^*\) alone.
\end{dcomment}


\msp{page 2}
Si
\[
f : I' \longrightarrow I,
\]
\printedpage{2}
% head: A. Grothendieck : Dérivateurs, Ch. III, Édition des Universités de Montpellier II et Paris VII
% foot: Version 2 février 2002
alors pour $F \in \mathrm{Ob}\,\D(I)$, on doit avoir des isomorphismes canoniques
\[
\left\{
\begin{array}{rcl}
f^*(\xi \times_{\D} M) & \simeq & \xi \times_{\D} f^*(M)\\
f^*\hm_{\D}(\xi, M) & \simeq & \hm_{\D}(\xi, f^*(M)),
\end{array}
\right.
\tag*{(3)}
\]
avec les compatibilités habituelles pour un composé $fg$, relatif à $I'' \xrightarrow{g} I' \xrightarrow{f} I$.

Mais je me place de préférence dans $\mathcal{A}$, ce qui devrait impliquer les constructions similaires dans tous les $\D(I)$, quitte à appliquer ce qui a été dit au dérivateur $\D_I$.

L'intuition algébrique fondamentale est celle-ci, en désignant par
\[
X \longmapsto \mathrm{hot}(X) : (\mathrm{Cat}) \longrightarrow (\Hot)
\]
le foncteur de localisation de (Cat) dans (Hot), obtenu en inversant les flèches dans (Cat) qui sont des `équivalences faibles'. Supposons $X$ dans Diag, soit
\[
p : X \longrightarrow e
\]
le foncteur à valeurs dans la catégorie ponctuelle. On veut que
\[
\left\{
\begin{array}{rcccl}
\mathrm{hot}(X) \times_{\D} M & \simeq & p_!(p^*(M)) & = & H^{\D}_{\bullet}(X, M)\\
\hm_{\D}(\mathrm{hot}(X), M) & \simeq & p_*(p^*(M)) & = & H^{\bullet}_{\D}(X, M)
\end{array}
\right.
\tag*{(4)}
\]
(isomorphismes canoniques).

\msp{page 3}
Ceci donnerait un principe de définition pour les opérations cherchées. Une difficulté vient du fait qu'en général, Diag n'est pas égal à (Cat), et que les types d'homotopie provenant de Diag ne donnent pas (même à isomorphisme près) \textit{tous} les types d'homotopie. Intuitivement, on ne trouve que des types d'homotopie `assez petits'. Mais on peut localiser Diag par rapport aux équivalences faibles (il vaut mieux dire `quasi-isomorphismes'), on trouve alors une catégorie $(\Hot)_0$, et un foncteur canonique
\[
\Hot_0 \longrightarrow \Hot,
\]
si on a de la chance, ce foncteur sera pleinement fidèle, et on essayera de décrire les foncteurs (1) ou (2) en y remplaçant Hot par la sous-catégorie pleine des `petits' types d'homotopie. On peut aussi oublier (provisoirement!) Hot, et travailler avec $\Hot_0$. C'est ce qu'on va faire pour le moment.

Pour pouvoir définir les foncteurs
\[
\left\{
\begin{array}{lccccccl}
\Hot_0 & \times & \mathcal{A} & \longrightarrow & \mathcal{A} & (\xi, F) & \longmapsto & \xi \times_{\D} F\\
\Hot_0^o & \times & \mathcal{A} & \longrightarrow & \mathcal{A} & (\xi, F) & \longmapsto & \hm_{\D}(\xi, F),
\end{array}
\right.
\tag*{(5)}
\]
je note que pour $F$ \texttt{[plutôt} $M$\texttt{]} fixé, $X$ variable dans Diag, les
\printedpage{3}
% head: A. Grothendieck : Dérivateurs, Ch. III, Édition des Universités de Montpellier II et Paris VII
% foot: Version 2 février 2002
\msp{page 4}
derniers membres dans (4)
\[
H^{\D}_{\bullet}(X, M), \qquad H^{\bullet}_{\D}(X, M)
\tag*{(6)}
\]
sont bien, l'un covariant, l'autre contravariant en $X$. Donc il y a , à isomorphisme unique près, une seule paire de bifoncteurs (5) satisfaisant (bifonctoriellement) les isomorphismes (4), si et seulement si les expressions (6), considérées comme fonctorielles en $X$, transforment quasi-isomorphismes en isomorphismes. Cela équivaut donc à la validité du % [sic] "il y a ," -> "il y a,", espace imprimée avant la virgule

\textbf{Théorème} (?) (\footnote{Nouvel axiome ?}). \textit{Soit $f : X' \longrightarrow X$ dans Diag. Si $f$ est un quasi-isomorphisme, c'est une $\D$-équivalence.}

Quand $X$ est la catégorie ponctuelle, ce n'est là autre que l'axiome Der 3 (p. 67). J'ignore si c'est une conséquence des axiomes Der 1 à Der 5 (d). Si ce n'est pas le cas, il faudrait l'ajouter aux axiomes, je crois. Je n'ai pas même prouvé qu'une équivalence de catégories dans Diag (au sens des catégories!) est une $\D$-équivalence!

\dgrblock{dgrb148}{The descent theorem and its domain}
\begin{dcomment}{status}{p4:J'ignore si c'est une conséquence des axiomes}{The descent theorem and its domain}
The question marked \emph{Théorème (?)} is precisely the obstruction to making the preceding formulas into operations of homotopy types. It has a positive answer for derivators on all small categories with the usual homotopy Kan extension axioms. The scope of the indexing category matters: that answer must not be silently substituted for a theorem about every restricted \(\mathrm{Diag}\) occurring in the manuscript.

First, the claimed equivalence with descent is formal. Let \(W\) be the nerve weak equivalences in \(\mathrm{Diag}\). A bifunctor \((X,M)\mapsto T_XM\) factors through \(\mathrm{Diag}[W^{-1}]\) exactly when every natural transformation \(T_f:T_{X'}\to T_X\), \(f\in W\), is invertible. The localization universal property applies also to natural transformations, giving the factorization with its specified comparison uniquely up to unique compatible isomorphism. The same argument applies contravariantly to \(C_X\). The two conditions are equivalent because \(T_X\dashv C_X\), and the comparison transformations are mates: a natural isomorphism between left adjoints has an invertible right mate, and conversely.

The easy case that Grothendieck says he has not yet checked needs no nerve theorem. If \(f:X'\to X\) is an equivalence of categories, a quasi-inverse and its invertible unit and counit give, by the prederivator structure, an equivalence \(f^*:\D(X)\to\D(X')\). It carries \(p_X^*M\) to \(p_{X'}^*M\). Hence it induces bijections
\[
 \Hom_{\D(X)}(p_X^*M,p_X^*N)
 \longrightarrow
 \Hom_{\D(X')}(p_{X'}^*M,p_{X'}^*N),
\]
which, through the two adjunctions, imply exactly the required comparison isomorphisms. This argument is already available for a prederivator possessing the displayed adjoints.

The passage from categorical equivalences to all nerve weak equivalences is substantial. For a derivator on all small categories, its constant-coefficient equivalences form a fundamental localizer \citep[Proposition 3.7]{Cisinski2008}; \citet[Theorem 2.2.11]{Cisinski2004Minimal} proves that nerve weak equivalences form the smallest fundamental localizer. Their combination is precisely \citet[Corollary 3.10]{Cisinski2008}: every nerve weak equivalence is a derivator equivalence. The more structural answer is the same paper's Theorem 3.24 and Corollary 3.26: the derivator of spaces is freely generated under homotopy colimits by a point. Its universal action yields the tensor construction; duality yields the cotensor construction. These results settle the modern full-domain version of the question.

The universal property also identifies the operation, not merely its existence. In its cocontinuous form, evaluation at the point identifies cocontinuous morphisms from the space derivator to \(\D\) with objects of \(\D(e)\). Let \(F_M\) correspond to \(M\). The classifying space \(B X\) is the homotopy colimit of the constant point diagram over \(X\). Preservation of that homotopy colimit forces
\[
 F_M(B X)\cong (p_X)_!p_X^*M.
\]
The equivalence of categories makes this construction functorial in maps \(M\to N\), and the statement for shifted derivators supplies its diagram versions. This explains why the modern theorem answers the particular question on the page: the descended action is constrained to have Grothendieck's formula (4). It also explains why full-domain hypotheses enter the literature. The universal object is the entire space derivator; replacing its diagram domain by an arbitrary subcategory is an additional operation whose effect must be checked.

For a restricted \(\mathrm{Diag}\), two further assertions must be kept separate. The internal localization \(\Hot_0\) exists under the usual size convention, but its comparison with \(\Hot\) need not be fully faithful merely by its definition. Even if it is fully faithful, it need not contain all homotopy types. One can work perfectly well with the action of \(\Hot_0\) once descent has been established there. What one cannot do is invoke Yoneda with test objects outside that category. This is exactly the limitation that returns on manuscript pages 7--10.
\end{dcomment}


Notons la \textit{formule d'adjonction}
\[
\mathrm{Hom}(\xi \times_{\D} M, N) \xrightarrow{\sim} \mathrm{Hom}(M, \hm_D(\xi, N)).
\tag*{(7)}
\] % [sic] \hom_D (D italique maigre) -> \hom_{\mathbf{D}}, le dérivateur est noté en gras partout ailleurs
Si $\xi$ provient de $X \in \mathrm{Diag}$, alors cela s'écrit
\[
\mathrm{Hom}(p_!p^*M, N) \xrightarrow{\sim} \mathrm{Hom}(M, p_*p^*N),
\tag*{(8)}
\]
\msp{page 5}
or par adjonction les deux membres sont canoniquement isomorphes à
\[
\mathrm{Hom}_{\D(X)}(p^*M, p^*N).
\]

\textbf{Exemple.} Prenant pour $\xi$ le type d'homotopie du cercle standard (représenté p.ex. par l'ensemble ordonné
\[
\begin{tikzcd}[row sep=small, column sep=small]
 & b & \\
0 \arrow[ur] \arrow[dr] & & 0' \arrow[ul] \arrow[dl] \ , \\
 & a &
\end{tikzcd}
\]
où on convient de prendre 0 comme origine, et d'orienter de 0 vers $a$), les deux foncteurs en question sont (sauf erreur) les foncteurs suspension $\Sigma$ et `cosuspension' ou foncteur `espace des lacets' $\Omega$. Sous réserve de vérifier que $\Omega F$ coïncide bien avec ce que j'ai décrit ailleurs, à propos des carrés substantiels cartésiens dans $\mathcal{A}$.

\dgrblock{dgrb149}{What the circle actually gives}
\begin{dcomment}{examples}{p5:où on convient de prendre 0 comme origine}{What the circle actually gives}
Grothendieck wants the circle action to recover the suspension and loop constructions, thereby linking the abstract action to the elementary squares of the next chapter. The four-object poset drawn here does have the homotopy type of a circle: its nerve has four vertices and four nondegenerate edges and no higher nondegenerate simplices. The difficulty is therefore not a defective model of \(S^1\). It is the distinction between an unpointed tensor and a reduced one, which he himself identifies on page 12.

In spaces,
\[
 T_{S^1}(M)=S^1\times M,\qquad
 C_{S^1}(N)=\operatorname{Map}(S^1,N).
\]
The second space is the free loop space. Choosing \(s\in S^1\) gives the evaluation map \(\operatorname{Map}(S^1,N)\to N\); its homotopy fibre over a chosen point \(n\) is \(\Omega_nN\). A basepoint in the indexing circle alone does not turn every free loop into a loop based at \(n\).

The algebraic calculation makes the excess term completely visible. With cohomological grading and a chosen orientation,
\[
 C_*(S^1;\mathbf Z)\simeq\mathbf Z[0]\oplus\mathbf Z[1].
\]
Here \(\mathbf Z[1]\) has its nonzero term in degree \(-1\). Consequently, in \(D(k\text{-Mod})\),
\[
 T_{S^1}(M)\simeq M\oplus M[1],
 \qquad C_{S^1}(M)\simeq M\oplus M[-1].
\]
The inclusion of the basepoint and the augmentation split off \(M\). Their reduced complements are \(M[1]\) and \(M[-1]\), the suspension and loop functors of this stable derivator. Taking \(M=k\ne0\), the extra degree-zero copy of \(k\) proves that the unqualified identification on the present page is false.

In a pointed homotopy theory the corresponding operation uses the smash tensor with a based circle; its cotensor is the based mapping object. The ordinary unpointed action is recovered by adjoining a disjoint basepoint to the space first. Formula (14) is therefore correct with \(\Sigma,\Omega\) understood as the unpointed circle tensor and cotensor, but it is not then the familiar reduced suspension-loop identity. With the reduced meanings, the right side is the loop space of the pointed external Hom at its zero morphism. The explicit Kan extension construction of those reduced functors appears in \citet[Definition 3.16 and Proposition 3.17]{Groth2013}; Chapter IV is already its essential diagrammatic pattern.
\end{dcomment}


Dans $\Hot_0$ on a une opération produit, en fait ça doit être le produit catégorique dans $\Hot_0$. Mais sur les modèles il est donné par
\[
\mathrm{hot}_0(X) \times \mathrm{hot}_0(Y) \ \simeq \ \mathrm{hot}_0(X \times Y)
\tag*{(8)}
\] % [sic] (8) -> (9), le numéro (8) est déjà employé plus haut sur cette page

% ---- from Der-ChIII-567_p006-010.tex
\printedpage{4}
% head: A. Grothendieck : Dérivateurs, Ch. III, Édition des Universités de Montpellier II et Paris VII
% foot: Version 2 février 2002

\texttt{[deux fois} (8)\texttt{]}. Ceci posé, on a les isomorphismes canoniques, pour $\xi,\eta\in\Ob\Hot_0$,
\[
\left\{
\begin{aligned}
\xi\times_{\D}(\eta\times_{\D}M) &\;\simeq\; (\xi\times\eta)\times_{\D}M\\
\hm_{\D}(\xi,\hm_{\D}(\eta,M)) &\;\simeq\; \hm_{\D}(\xi\times\eta,M).
\end{aligned}
\right.
\tag*{(9)}
\]

Si $\xi$, $\eta$ sont représentés par $X$, $Y$, la première formule p.ex.
\msp{page 6}
s’établit en calculant le deuxième membre comme
\[
\varphi_!(M_{X\times Y})\simeq p_!(\;\underbrace{q'_!(M_{X\times Y})}_{\substack{\simeq\ q_!(M_Y)_X\\ \text{car $q$ est $\mathbf{D}$-lisse}}}\;)\simeq\xi\times_{\D}(\eta\times_{\D}M),
\]
O.K.
\[
\begin{tikzcd}
X \ar[d,"p"'] & X\times_e Y \ar[l,"q'"'] \ar[d,"p'"] \ar[dl,dotted,"\varphi"'] \\
e & Y \ar[l,"q"]
\end{tikzcd}
\]

J’ai aussi envie de prouver :

\dgrblock{dgrb150}{Associativity and exactness from product squares}
\begin{dcomment}{exposition}{p6:J’ai aussi envie de prouver :}{Associativity and exactness from product squares}
Grothendieck next checks that the action has the algebra expected of an action. The calculation is a Fubini argument for homotopy Kan extensions, and its force is that it works before any mapping-space enrichment has been supplied. We retain his convention that coefficients on \(I\) have underlying diagrams of shape \(I^o\).

First, the product on \(\Hot_0\) mentioned just above can be justified on the stated models. Assume that \(\mathrm{Diag}\) is closed under binary products. The nerve preserves those products, and weak equivalences of spaces are preserved by taking a product with a fixed space. Hence the product functor sends a pair of nerve weak equivalences to a nerve weak equivalence and descends to the localization. The diagonal functor descends too. Before localization the diagonal is left adjoint to the binary product; their unit and counit descend as natural transformations, and still satisfy the triangle identities. The product functor on the localization is therefore right adjoint to its diagonal, which is exactly the universal property of a categorical product. This proves the sentence beginning \emph{en fait ça doit être}, under the required closure assumption, without assuming that \(\Hot_0\to\Hot\) is fully faithful. Existence of this product should not be confused with closure under the right adjoints to products, which becomes relevant to (13).

For \(X,Y\in\mathrm{Diag}\), compose the projections \(X\times Y\to X\to e\). Pseudofunctoriality of left Kan extension identifies their composite with the extension to \(e\). Base change for the product square identifies the first intermediate extension of a constant coefficient with the constant coefficient \(T_YM\) on \(X\). Thus
\[
 T_XT_YM\cong T_{X\times Y}M.
\]
The terminal category supplies \(T_e\cong\mathrm{id}\). The analogous right Kan calculation gives \(C_XC_Y\cong C_{X\times Y}\), with the harmless interchange of \(X\) and \(Y\) supplied by the product symmetry. The associativity and unit comparisons are canonical mates. Their coherence follows by pasting product squares and using the composition constraints for Kan extension; it is not a second independent choice of associators.

The same argument proves more than preservation of finite colimits or finite limits. Fix \(X\), let \(f:I\to J\), and write \(p:X\times I\to I\), \(q:X\times J\to J\), \(f'=1_X\times f\). The product square is homotopy exact in both directions. Therefore
\[
\begin{aligned}
 f_*C_{X,I}F
 &=f_*p_*p^*F\\
 &\cong q_*f'_*p^*F\\
 &\cong q_*q^*f_*F=C_{X,J}f_*F .
\end{aligned}
\]
The dual calculation gives \(f_!T_{X,I}\cong T_{X,J}f_!\). Every Kan extension permitted by \(\mathrm{Diag}\) is covered. This is the precise content of \emph{exact à gauche} and \emph{exact à droite} here.

For completeness, the decisive product comparison can be checked directly by the pointwise formula. At \((x,j)\), the comma category computing the appropriate Kan extension is a product of the comma category for \(f\) with a slice of \(X\) having an extremal object. Homotopy Kan extension over the latter factor reduces to evaluation. The resulting isomorphism is the mate of the product square, so its pointwise verification also identifies the comparison that must be used. This is the special case of base change for Grothendieck fibrations and opfibrations formalized in \citet[Proposition 1.24]{Groth2013}.

Finally,
\[
 \Hom(T_XM,N)
 \cong\Hom_{\D(X)}(p_X^*M,p_X^*N)
 \cong\Hom(M,C_XN)
\]
proves the tensor-cotensor adjunction without representability in \(\Hot\). This separation is useful: (7), (9), and Proposition 1 are consequences of the Kan formalism once descent is available; the external Hom sought next is an additional representation problem.
\end{dcomment}


\textbf{Proposition 1.} \textit{Soit $\xi\in\Ob\Hot_0$. Alors $F\longmapsto\xi\times_{\D}F$, en tant que morphisme de dérivateurs $\D\longrightarrow\D$, est} exact à droite, \textit{et $F\longmapsto\hm_{\D}(\xi,F)$ est exact à gauche.}

On peut supposer que $\xi$ est défini par $X$ dans Diag. Soit $f:I\longrightarrow J$, soit $F$ dans $\D(I)$, il faut prouver
\[
f_*(\hm_{\D}(\xi,F))\;\simeq\;\hm_{\D}(\xi,f_*(F)).
\]
Considérons donc
\[
\begin{tikzcd}
X\times I \ar[r,"f'"] \ar[d,"{p=(x,i)\longmapsto i}"'] & X\times J \ar[d,"{q=(x,j)\longmapsto j}"] \\
I \ar[r,"f"] & J.
\end{tikzcd}
\]

On a
\[
\hm_{\D}(\xi,F)=p_*(p^*F),
\]
donc
\[
f_*(\hm_{\D}(\xi,F))=\underbrace{f_*p_*}_{(fp)_*}(p^*(F))=(qf')_*p^*(F)\simeq q_*(f'_*p^*(F)).
\]
Mais
\[
f'_*p^*(F)\;\simeq\;q^*f_*(F),
\]
formule de changement de base des $f_*$, pour le changement de base $q$, qui est lisse, car déduit de $X\longrightarrow e$, qui est lisse, par changement de base $J\longrightarrow e$. On trouve
\printedpage{5}
% head: A. Grothendieck : Dérivateurs, Ch. III, Édition des Universités de Montpellier II et Paris VII
% foot: Version 2 février 2002
\msp{page 7}
donc
\[
f_*(\hm_{\D}(\xi,F))\simeq q_*q^*(f_*(F))\defeq\hm_{\D}(\xi,f_*(F)),
\]
q.e.d. L’énoncé pour $F\longmapsto\xi\times_{\D}F$ est dual de celui qu’on vient de prouver.

Reprenons, pour $\xi\in(\mathrm{hot})_0$ \texttt{[}$=\Hot_0$\texttt{]}, la formule d’adjonction
\[
\Hom(\xi\times_{\D}M,N)\simeq\Hom(M,\hm_{\D}(\xi,N))\defeq\Hom_{\xi}(M,N).
\tag*{$(*)$}
\]

Si $\xi$ est défini par $X$ dans Diag, les deux membres de $(*)$ s’identifient à
\[
\Hom_{\D(X)}(p^*M,p^*N),
\]
où $p:X\longrightarrow e$ est le foncteur structural.

Fixons $M$, $N$, et considérons l’expression $(*)$ comme contrafoncteur
\[
(\Hot_0)^o\;\longrightarrow\;\mathrm{Ens}.
\]

Il semblerait que dans tous les cas connus, ce foncteur soit ‘représentable’ par un objet, sinon de $\Hot_0$ lui-même, du moins un objet dans Hot. Si on ne peut le définir comme objet $H$ de $\Hot_0$, le fait qu’on ait un isomorphisme canonique
\[
\Hom_{\xi}(M,N)\;\simeq\Hom_{\Hot}(\xi,H),
\tag*{(10)}
\]
fonctoriel en $\xi\in\Ob\Hot_0$, ne \texttt{[le]} définit pas forcément
\msp{page 8}
à isomorphisme unique près. En fait, quand
\[
\D=\D_{\mathcal{M},\Sigma}
\]
est défini par une catégorie de modèles $(\mathcal{M},\Sigma)$, alors il y a une définition naturelle d’une opération $\homext$ (Hom ‘externe’) à valeurs dans Hot, savoir un foncteur
\[
\left\{
\begin{aligned}
&\mathcal{A}^o\times\mathcal{A} &&\longrightarrow &&\mathrm{Hot}\\
&(M,N) &&\longmapsto &&\homext(M,N),
\end{aligned}
\right.
\tag*{(11)}
\]
où $\mathcal{A}=\mathcal{M}\Sigma^{-1}$ est la catégorie fondamentale de $\D=\D_{\mathcal{M},\Sigma}$. Il est vrai que je n’ai pas vérifié, si p. ex. $(\mathcal{M},\Sigma)$ est sous-jacente à une catégorie de modèles quillénienne, de sorte que (par \textsc{Thomason}) $\D$ est bien un dérivateur, que cette notion de $\homext$ satisfait bien à une formule de la forme (10). Ici je vais le poser [?], au besoin, ces données supplémentaires : donnée de (11), et d’un isomorphisme fonctoriel
\[
\underbrace{\Hom_{\xi}(M,N)}_{\begin{subarray}{l}=\ \Hom(\xi\times_{\D}M,N)\\ =\ \Hom(M,\hm_{\D}(\xi,N))\end{subarray}}
\quad\simeq\;\Hom_{\Hot}(\xi,\homext(M,N)).
\tag*{(12)}
\]

\dgrblock{dgrb151}{External Hom as a derived mapping space}
\begin{dcomment}{later}{p8:Il est vrai que je n’ai pas vérifié}{External Hom as a derived mapping space}
Grothendieck is now trying to recover the space of maps from the already available operations on coefficients. In a model category this space had been constructed independently: the hammock localization of \citet[Corollary 4.2, Proposition 4.4, and Corollary 4.7]{DwyerKan1980Function} supplies mapping complexes, identifies their components with localized morphisms, and compares them with the enrichment of a simplicial model category. This work predates the manuscript. Chapter V will approach the same issue through categories of zigzags.

Here is the verification of (12) in the simplicial case. Choose a cofibrant representative \(QM\) and a fibrant representative \(RN\). A simplicial set \(K\) representing \(\xi\) is cofibrant in the usual model structure. The simplicial tensor-Hom adjunction gives
\[
 \operatorname{Map}_{\mathcal M}(K\otimes QM,RN)
 \cong
 \operatorname{Map}_{\mathrm{sSet}}
       (K,\operatorname{Map}_{\mathcal M}(QM,RN)).
\]
Taking components identifies the left side with
\(\Hom_{\mathrm{Ho}(\mathcal M)}(K\otimes^{\mathbf L}M,N)\).
The right side is
\([K,\operatorname{Map}_{\mathcal M}(QM,RN)]\), because the inner mapping complex is fibrant. Moreover \(K\otimes^{\mathbf L}M\) is the homotopy colimit of the constant \(M\)-diagram indexed by the space \(K\). For \(K=B X\), the simplicial replacement of the constant \(X^o\)-diagram gives exactly that tensor; its simplices are copies of \(QM\) indexed by strings of arrows in \(X^o\). Reversal identifies its realized indexing space with \(B X\). This proves the formula with the particular tensor defined in (4).

The cotensor identity follows either from the dual model construction or from the tensor-cotensor adjunction already proved. Derived function complexes and framings give the corresponding statement for a general Quillen model category; the tensor, cotensor, and enrichment compatibilities are treated systematically in \citet[Sections 10.2--10.3]{Riehl2014}. Thus the extra datum is supplied by the model presentation and is independent, up to homotopy equivalence, of the chosen resolutions.

For a presentable \(\infty\)-category \(\mathcal C\), the same calculation has the particularly transparent form
\[
 \operatorname{Map}_{\mathcal C}
 \bigl(\operatorname*{colim}_{K}M,N\bigr)
 \simeq
 \operatorname{Map}_{\mathcal S}
 \bigl(K,\operatorname{Map}_{\mathcal C}(M,N)\bigr).
\]
It is the universal property of the colimit indexed by the space \(K\). Passing to homotopy categories gives (12); no new choice of external Hom is needed once \(\mathcal C\) is given. The universal role of space-valued presheaves is made precise in \citet[Theorem 5.1.5.6]{Lurie2009HTT}.

For a completely abstract derivator, one should distinguish an action by spaces from a closed action possessing these Hom objects. \citet[Section 1, pp.~3--4, and Definition 4.2]{GrothShulman2017} explicitly distinguish the two structures. The cited results settle the model and presentable \(\infty\)-categorical cases; they do not justify erasing Grothendieck's representability hypothesis for every restricted diagram domain considered here.
\end{dcomment}


\printedpage{6}
% head: A. Grothendieck : Dérivateurs, Ch. III, Édition des Universités de Montpellier II et Paris VII
% foot: Version 2 février 2002
\msp{page 9}
Si maintenant $\eta$ est un deuxième élément de $\Hot_0$, on a envie d’écrire
\[
\left\{
\begin{aligned}
\homext(\eta\times_{\D}M,N) &\;\simeq\;\homext(M,\hm_{\D}(\eta,N))\\
&\;\simeq\;\hm_{\HOT}(\eta,\homext(M,N)),
\end{aligned}
\right.
\tag*{(13)}
\]
où dans le deuxième membre, $\hm_{\HOT}$ est l’opération (1) dans le cas où $\D=\mathrm{HOT}$, et qui est défini aussi en termes du produit cartésien $\xi\times\eta$ dans Hot, par
\[
\Hom_{\Hot}(\xi,\hm_{\HOT}(\eta,\zeta))\;\simeq\;\Hom_{\Hot}(\xi\times\eta,\zeta),
\]
en d’autres termes, pour $\eta$ fixé, le foncteur
\[
\zeta\;\longmapsto\;\hm_{\HOT}(\eta,\zeta)
\]
est défini comme le foncteur adjoint à droite du foncteur
\[
\xi\;\longmapsto\;\xi\times\eta.
\]

Pour vérifier heuristiquement la formule (13) (vérification valable ‘en forme’ si les trois objets envisagés sont en fait définis comme objets de $\Hot_0\ \ldots$), on prend $\Hom_{\Hot}(\xi,-)$ à valeurs dans les trois membres, pour $\xi$ donné dans $\Hot_0$. Pour le premier, on trouve
\[
\begin{aligned}
\Hom_{\xi}(\eta\times_{\D}M,N) &\;=\;\Hom(\xi\times_{\D}(\eta\times_{\D}M),N)\\
&\;=\;\Hom((\xi\times\eta)\times_{\D}M,N).
\end{aligned}
\]
Pour le deuxième
\[
\begin{aligned}
\Hom_{\xi}(M,\hm_{\D}(\eta,N)) &\;=\;\Hom_{\mathcal{A}}(\xi\times_{\D}M,\hm_{\D}(\eta,N))\\
&\;\simeq\;\Hom_{\mathcal{A}}(\eta\times_{\D}(\xi\times_{\D}M),N)\\
&\;\simeq\;\Hom_{\mathcal{A}}((\eta\times\xi)\times_{\D}M,N),
\end{aligned}
\]
\msp{page 10}
et enfin pour le troisième
\[
\begin{aligned}
\Hom_{\Hot}(\xi\times\eta,\homext(M,N)) &\;\simeq\;\Hom_{\xi\times\eta}(M,N)\\
&\;\simeq\;\Hom((\xi\times\eta)\times_{\D}M,N),
\end{aligned}
\]
donc toujours la même chose ! Pas de doute, ça doit marcher.

S’il faut se donner le $\homext$ (11) comme donnée supplémentaire, avec de plus un isomorphisme trifonctoriel (12), il faut aussi se donner les isomorphismes trifonctoriels (13), \textit{‘compatibles’} avec l’isomorphisme (12) au sens du calcul précédent.

\dgrblock{dgrb152}{What the Yoneda calculation proves}
\begin{dcomment}{exposition}{p10:S’il faut se donner le}{What the Yoneda calculation proves}
Grothendieck's three calculations are intended to identify the external Hom of a tensor with the cotensor of an external Hom. The calculation is correct. His qualification about the test category is also correct, and is part of the argument rather than an unnecessary hesitation.

Suppose the external Hom represents (12), and suppose that all three candidate objects in (13) belong to \(\Hot_0\). For every \(\xi\in\Hot_0\), tensor associativity and the tensor-cotensor adjunction give natural identifications of their represented functors with
\[
 \xi\longmapsto
 \Hom_{\mathcal A}((\xi\times\eta)\times_{\D}M,N).
\]
Yoneda in \(\Hot_0\) then gives a unique isomorphism respecting these identifications. Naturality in \(\eta,M,N\) follows from the same uniqueness. The last candidate additionally requires the relevant function object to exist in \(\Hot_0\); closure under products alone does not imply that closure under function objects.

If the candidates lie only in \(\Hot\), the calculation instead identifies their \emph{restricted} represented functors. Such an identification need not determine the objects. For a concrete example, take \(\Hot_0\) to be the full subcategory of discrete spaces. For any discrete set \(S\) and any nonempty connected space \(H\), the set \([S,H]\) has one element: a map is a choice of a component for each element of \(S\), and there is only one component. Consequently the point and \(S^1\) define isomorphic functors on this test category, although they are not isomorphic homotopy types. This example isolates the exact logical failure that would result from applying unrestricted Yoneda here.

There are two legitimate completions. One may assume the restricted Yoneda functor is fully faithful on the particular objects under consideration, and then the displayed calculation suffices. Alternatively, a model or \(\infty\)-categorical construction gives the actual comparison of mapping spaces before passage to homotopy types, proving (13) directly. The manuscript follows the first course conditionally and otherwise proposes to retain the comparisons as part of the additional structure. No gap should be charged to a proof that explicitly states this restriction.
\end{dcomment}


\printedpage{7}
% head: A. Grothendieck : Dérivateurs, Ch. III, Édition des Universités de Montpellier II et Paris VII
% foot: Version 2 février 2002

Par exemple, prenant $\eta=$ type d’homotopie du cercle $S^1$, les formules (13) se spécialisent en
\[
\begin{aligned}
\homext(\Sigma M,N) &\;\simeq\;\homext(M,\Omega N)\\
&\;\simeq\;\Omega\homext(M,N)\\
&\;\defeq\;\hm_{\HOT}(S^1,\homext(M,N)).
\end{aligned}
\tag*{(14)}
\]

Mais cette formule dans Hot n’est sans doute pas tout à fait la bonne, car le bon formalisme de $\Omega$, $\Sigma$ se fait dans les catégories \textit{ponctuées,} et en particulier dans la catégorie homotopique ponctué $\Hot_{\bullet}$. % [sic] ponctué -> ponctuée, agreement with « catégorie »

Je vais expliciter un exemple, en prenant le dérivateur
\[
\mathrm{Dér}(A)
\]
‘dérivé’ d’une catégorie abélienne $A$. Donc $\D(e)=$ catégorie dérivée de $A$. Pour éviter des
\msp{page 11}
difficultés accessoires ici, je vais supposer que $A$ est la catégorie des $k$-modules, $k$ un anneau quelconque.

À tout $\xi\in\Ob\Hot$ est associé un ‘abélianisé’
\[
\xi_{\ab}\defeq H_{\bullet}(\xi)\in\underbrace{D(\mathrm{Ab})}_{\text{ou }\mathbf{D}(\mathbf{Z}\text{-Mod})}\;,\qquad\quad\text{catégorie dérivée des groupes abéliens,}
\]
objet qui n’a de cohomologie qu’en degré $\le 0$. Ceci dit, si $M$ est dans $D(k\text{-Mod})$, on aura
\[
\left\{
\begin{aligned}
\xi\times_{\D}M &\;\simeq\;\xi_{\ab}\Ltens_{\mathbf{Z}}M\\
\hm_{\D}(\xi,M) &\;\simeq\;\RHom_{\mathbf{Z}}(\xi_{\ab},M),
\end{aligned}
\right.
\tag*{(15)}
\]
formule qui est en cheville avec la formule de Künneth pour l’homologie
\[
(\xi\times\eta)_{\ab}\simeq\xi_{\ab}\Ltens_{\mathbf{Z}}\eta_{\ab},
\tag*{(16)}
\]
ce qui redonne un aspect familier aux formules (9), qui se réduisent à
\[
\left\{
\begin{aligned}
\xi_{\ab}\Ltens_{\mathbf{Z}}(\eta_{\ab}\Ltens_{\mathbf{Z}}M) &\;\simeq\;(\xi_{\ab}\Ltens_{\mathbf{Z}}\eta_{\ab})\Ltens_{\mathbf{Z}}M\\
\RHom_{\mathbf{Z}}(\xi_{\ab},\RHom_{\mathbf{Z}}(\eta_{\ab},M)) &\;\simeq\;\RHom_{\mathbf{Z}}(\xi_{\ab}\Ltens_{\mathbf{Z}}\eta_{\ab},M).\qquad\text{(\footnotemark[2])}
\end{aligned}
\right.
\]
\footnotetext[2]{Expliciter aussi la formule d’adjonction (7) (p. 4).}

D’autre part, existe-t-il bien un $\homext(M,N)$, si $M$, $N$ dans $D(k\text{-Mod})$ ? On pense à $\RHom_k(M,N)$, mais pour en faire un type d’homotopie, on doit ‘tronquer en degrés $>0$’. On pourrait donc prendre
\[
\homext(M,N)\;\simeq\;\tau_{\le 0}(\RHom_k(M,N)).
\]

\dgrblock{dgrb153}{Integral chains and the truncation of derived Hom}
\begin{dcomment}{examples}{p11:D’autre part, existe-t-il bien}{Integral chains and the truncation of derived Hom}
Grothendieck turns to complexes because they make the external Hom computable. We can verify every part of the proposed formula, provided that \emph{abélianisation} means the derived free abelian-group construction and that the truncation is the canonical one.

For a simplicial set \(K\), let
\[
 A_K=N\mathbf Z[K]
\]
be its normalized free chain complex, written cohomologically in degrees at most zero. This is the object denoted \(H_\bullet(\xi)\) in the manuscript. It is a chain complex up to quasi-isomorphism, not merely its graded homology groups. The free simplicial abelian-group functor is left adjoint to the forgetful functor. After deriving and applying the Dold--Kan equivalence, it gives
\[
 [K,U\,\mathrm{DK}(B)]
 \cong
 \Hom_{D(\mathbf Z)}(A_K,B),
 \qquad B\in D(\mathbf Z)^{\leq0}.
\]
Here \(\mathrm{DK}\) means denormalization with the change from cohomological to homological grading, and \(U\) forgets the abelian-group structure. The underlying equivalence is \citet[Theorem 14.24.3]{StacksDoldKan}. It is this derived adjunction that answers the question in footnote 3; being only an abelian-group object in the ordinary homotopy category would not by itself specify the requisite strict or coherent abelian structure.

Put \(R=\mathbf R\Hom_k(M,N)\), viewed as an object of \(D(\mathbf Z)\). For a noncommutative \(k\), this target is still correct: there is generally no left \(k\)-module structure on this Hom complex. Since \(A_K\in D(\mathbf Z)^{\leq0}\), the truncation adjunction gives
\[
\begin{aligned}
 {}[K,U\,\mathrm{DK}(\tau_{\leq0}R)]
 &\cong\Hom_{D(\mathbf Z)}(A_K,\tau_{\leq0}R)\\
 &\cong\Hom_{D(\mathbf Z)}(A_K,R)\\
 &\cong\Hom_{D(k)}(A_K\otimes_{\mathbf Z}^{\mathbf L}M,N).
\end{aligned}
\]
The final equality is the derived tensor-Hom adjunction; \(k\) acts on the factor \(M\). This proves (10) and (12) for the claimed external Hom, with the otherwise suppressed Dold--Kan and forgetful functors restored.

The truncation is the smart truncation: its degree-zero term is \(\ker(d:R^0\to R^1)\), and its negative terms are those of \(R\). Simply deleting the positive degrees would replace degree-zero cocycles by all of \(R^0\), generally giving a wrong set of components. The homotopy groups of the resulting simplicial abelian group are
\[
 \pi_n\homext(M,N)
 \cong H^{-n}R
 \cong\Hom_{D(k)}(M,N[-n]),\qquad n\geq0.
\]
For instance, with \(k=\mathbf Z\), \(M=\mathbf Z\), and \(N=\mathbf Z[r]\), \(r\geq0\), one obtains \(K(\mathbf Z,r)\). The adjunction then reads
\([K,K(\mathbf Z,r)]\cong H^r(K;\mathbf Z)\), including degree zero with its usual unreduced interpretation.

The other formula in (15) follows from the same adjunction:
\[
 T_KM=A_K\otimes_{\mathbf Z}^{\mathbf L}M,\qquad
 C_KM=\mathbf R\Hom_{\mathbf Z}(A_K,M).
\]
The Eilenberg--Zilber comparison identifies \(A_{K\times L}\) with \(A_K\otimes_{\mathbf Z}^{\mathbf L}A_L\) \citep[Section 1, theorem, pp.~200--201]{EilenbergZilber1953}; normalized free chains are bounded above complexes of free abelian groups in this grading, so the ordinary tensor already computes the derived tensor. Formula (16) is this chain-level equivalence. A formula involving only the tensor product of homology groups would lose the Tor terms and would not suffice for Grothendieck's argument.
\end{dcomment}


\printedpage{8}
% head: A. Grothendieck : Dérivateurs, Ch. III, Édition des Universités de Montpellier II et Paris VII
% foot: Version 2 février 2002
\msp{page 12}
La formule fondamentale (10) est-elle alors valable ? Le deuxième membre est en effet
\[
\Hom_{\Hot}(\xi,H)\;\simeq\;\Hom_{\mathrm{Hotab}}(\xi_{\ab},H)\qquad\quad\text{(\footnotemark[3])}
\]
\footnotetext[3]{? ? est-ce vrai ? ? ?}
car $H$ est un objet abélien de Hot (cf. le tapis sur l’abélianisation \ldots), à partir de là c’est gagné :
\[
\begin{aligned}
\Hom_{\underbrace{\scriptstyle\mathrm{Hotab}}_{\begin{subarray}{l}\approx\ D(\mathbf{Z}\text{-Mod})^{\le 0},\\ \text{sous-catégorie \textit{pleine}}\\ \text{de }D(\mathbf{Z}\text{-Mod})\end{subarray}}}(\xi_{\ab},\tau_{\le 0}\RHom_k(M,N)) &\;\simeq\;\Hom_{D(\mathbf{Z}\text{-Mod})}(\xi_{\ab},\RHom_k(M,N))\\
&\;\simeq\;\Hom_{D(k\text{-Mod})}(\xi_{\ab}\times M,N),
\end{aligned}
\]
O.K.

D’autre part, le bifoncteur
\[
(M,N)\;\longmapsto\;\homext(M,N)\;=\;\tau_{\le 0}\RHom_k(M,N)
\]
se factorise canoniquement par Hotab et a fortiori par $\Hot_{\bullet}$. Le ‘bon’ $\homext(M,N)$ serait ici plutôt celui dans Hotab, voire même dans la catégorie $D^{-}(\mathbf{Z}\text{-Mod})$, déduite de celle là ‘en rendant le foncteur translation inversible’. Il est évident d’ailleurs qu’ici le $\xi_{\ab}\times_{\D}M$, où $\xi=S^1$, n’est pas ‘la bonne’ suspension, qui doit être le décalage des degrés – $\xi_{\ab}$ ne devrait avoir de la cohomologie qu’en degré $-1$, et non aussi en degré $0$.

\dgrblock{dgrb154}{Keeping all degrees of the mapping complex}
\begin{dcomment}{status}{p12:D’autre part, le bifoncteur}{Keeping all degrees of the mapping complex}
Grothendieck's last suggestion seeks an external Hom that remembers the linear structure and the information restored by suspension. That is the correct direction, but the proposed bounded-above target requires a qualification when the coefficient category is the unbounded derived category.

The homotopy category of simplicial abelian groups is equivalent to \(D(\mathbf Z)^{\leq0}\). Allowing finite shifts of its objects produces \(D^-(\mathbf Z)\): an object in the union has cohomology bounded above, and conversely any bounded-above complex can be shifted into nonpositive degrees. This observation correctly describes the stated formal inversion of translation at the level of objects and morphisms. It does not imply that every derived Hom complex belongs to this union.

A two-term algebra already shows the difficulty. Let \(F\) be a field, let \(k=F[\epsilon]/(\epsilon^2)\), and let \(L=k/(\epsilon)\). A projective resolution of \(L\) is
\[
 \cdots\xrightarrow{\epsilon}k
 \xrightarrow{\epsilon}k
 \xrightarrow{\epsilon}k\longrightarrow L\longrightarrow0.
\]
It is exact because the kernel and image of multiplication by \(\epsilon\) are both the ideal \((\epsilon)\). Applying \(\Hom_k(-,L)\) gives one copy of \(F\) in every nonnegative cohomological degree and zero differentials. Thus
\[
 H^n\mathbf R\Hom_k(L,L)=\operatorname{Ext}^n_k(L,L)
 \cong F\qquad(n\geq0).
\]
The full derived Hom is not in \(D^-(\mathbf Z)\). The present external mapping space is nevertheless well defined: its nonpositive truncation is the discrete group \(F\). The omitted extension groups reappear as
\(\pi_0\homext(L,L[n])\), for every \(n\geq0\). This example explains why a single mapping space and the whole stable mapping object are different levels of information.

For complexes, the full enhancement is by the unbounded complex \(\mathbf R\Hom_k(M,N)\); in a general stable homotopy theory it is by a mapping spectrum. The classical differential graded formulation and its relation to triangulated categories are discussed by \citet[Sections 2--3]{Keller2006DG}. The modern theory of monoidal and enriched derivators distinguishes the unpointed, pointed, and stable actions and the existence of their right adjoints \citep[Sections 3 and 8--9]{GrothPontoShulman2014Additivity}.

The correction here concerns the generality of a tentative proposal, not a theorem whose proof has collapsed. The formulas involving \(\tau_{\leq0}\mathbf R\Hom\) are valid. The earlier circle identification is repaired by reduced chains. An additional boundedness hypothesis, such as suitable finite projective dimension and boundedness on the inputs, is needed only if one insists that the entire linear external Hom take values in \(D^-(\mathbf Z)\).
\end{dcomment}
\dgrchapter{IV}{Chapter IV}{Diagrammes substantiels}
\dgrpart{0}{III-V}{Manuscript pages 73--83}{III-V\enspace\textperiodcentered\enspace Chapter IV\enspace\textperiodcentered\enspace manuscript pages 73--83}{Chapter IV \textperiodcentered\ \textup{III-V}}
% ---- from Der-ChIV-567_p001-005.tex
\printedpage{}% no page number printed on this page
% head: (none)

\begin{center}
{\LARGE\textsf{LES DÉRIVATEURS}}

\bigskip

\textbf{ALEXANDRE GROTHENDIECK}

\bigskip\bigskip

\textbf{Chapitre IV}

\medskip

\textbf{Diagrammes substantiels}
\end{center}

\clearpage

\printedpage{}% no page number printed on this page
% head: (none)

\noindent Ce texte a été déchiffré et transcrit en \LaTeXe\ par M. Künzer. Il a été édité par M. Künzer, J. Malgoire, et G. Maltsiniotis. La transcription est aussi fidèle que possible au manuscrit. Pour les quelques corrections évidentes, ou rares commentaires des éditeurs, ainsi que pour la numérotation originale des pages du manuscrit, les caractères de machine à écrire \texttt{[typewriter]} entre crochets sont utilisés. Un point d'interrogation entre crochets signifie que l'on n'est pas sûr du mot qui précède.

\noindent Cette édition est provisoire. Les remarques, commentaires et corrections sont bienvenus. Envoyer un message à :

\noindent\textbf{maltsin@math.jussieu.fr}

\hfill G. Maltsiniotis\hspace*{2cm}

\clearpage

\printedpage{1}
% head: A. Grothendieck : Dérivateurs, Ch. IV, Édition des Universités de Montpellier II et Paris VII
% foot: Version 2 février 2002

\noindent\texttt{[Il n'y a pas de pages 1-72]}

\msp{page 73}

\subsection*{Construction de diagrammes substantiels pour un dérivateur.}

\noindent Soit $\D$ un prédérivateur, $I$ dans Diag, considérons

\noindent $\varphi_I : \mathcal{D}(I) \longrightarrow \uHom(I^o,\cA)$,\ \ où $\cA=\D(e)$ est la catégorie fondamentale des coefficients. % [sic] \mathcal{D}(I) -> \mathbf{D}(I), elsewhere the prederivator is the bold D

Les objets de $\uHom(I^o,\cA)$ s'interprètent comme des \textit{diagrammes de type $I^o$} dans $\cA$. Ce langage est surtout commode si $I$ est ordonnée. Notons, pour $F\in\mathrm{Ob}\,\D(I)$, par $(F_i)_{i\in\mathrm{Ob}\,I}$ le $I^o$-diagramme associé dans $\cA$. Par l'axiome Dér 2, le foncteur
\[
\varphi : F \longmapsto (F_i)_{i\in\mathrm{Ob}\,I}
\]
est conservatif. Mais il n'a aucune envie d'être fidèle, encore moins pleinement fidèle. Quand on connaît le diagramme $(F_i)_i$, on est loin de connaître $F$, pas même à isomorphisme non unique près, je présume. Il n'est nullement clair que si $F$, $F'$ dans $\D(I)$ sont tels que les diagrammes associés soient isomorphes, ils soient eux-mêmes isomorphes.

\dgrblock{dgrb155}{Why the skeleton is insufficient}
\begin{dcomment}{context}{p73:Il n'est nullement clair que}{Why the skeleton is insufficient}
Grothendieck's term \emph{diagramme substantiel} names the information on which the whole project depends: a diagram before its coherent relations have been discarded. The point is not that an ordinary commutative diagram has been drawn inaccurately. It is that several different homotopies and compatibilities can produce the same drawing in the homotopy category.

In the model example,
\[
 \D(I)=\mathcal M^{I^o}[(W^{I^o})^{-1}],
 \qquad
 \varphi_I:\D(I)\longrightarrow(\mathrm{Ho}\mathcal M)^{I^o}.
\]
Localization is performed on the category of whole diagrams. The target instead localizes the values first and subsequently records arrows and their equations. These operations need not commute. Axiom Der~2 says that a morphism of coherent diagrams which is an isomorphism at every vertex is invertible. It says neither that every skeletal morphism lifts nor that two lifts agree. The distinction was illustrated in the commentary on Chapter~I, manuscript p.~11 (\emph{Conservative is not faithful}); there is no need to repeat that counterexample here.

For the walking arrow one can make a useful positive statement. A \emph{strong} derivator satisfies fullness and essential surjectivity of
\[
 \D(I\times[1])\longrightarrow\D(I)^{[1]^o}
\]
for every \(I\). This is the usual modern Der~5, not an identification of the numbered axioms appearing earlier in the manuscript \citep[Definition 1.8]{Groth2013}. Fullness together with conservativity implies that two coherent arrows whose skeletal arrows are isomorphic are themselves isomorphic: lift the skeletal isomorphism by fullness, then use its invertibility at both endpoints. It does not imply a unique lift of that isomorphism. Thus Grothendieck's question about objects with isomorphic skeletons has a different answer for a single arrow in a strong derivator than for an unrestricted diagram shape.

The historical motivation is particularly close to the present construction. In \emph{Pursuing Stacks}, \citet[Section 69, pp.~1--2 of the edited extract]{Grothendieck1983} recalls the need, arising around SGA~5, to supplement a derived category by derived categories of triangles and filtered complexes. Here he seeks a procedure that constructs the supplements and their maps from Kan extension itself. Heller's parallel treatment packages homotopy theories as systems of diagram categories \citep[Introduction and Chapter I]{Heller1988}. Neither perspective asks one to manufacture a functorial cone on the bare homotopy category.

The two examples pursued below are deliberately chosen to retain a universal property throughout the construction. The loop square is determined by a zero extension followed by a right Kan extension. The unit and counit arrows are determined by a category that encodes the functor \(f\). This is why canonical coherent versions of these particular arrows can be obtained even though a general skeletal arrow has no specified canonical lift.
\end{dcomment}


Pour insister sur ce point, les objets de $\D(I)$ s'appelleront les `\textit{diagrammes substantiels} de type $I^o$ dans $\cA$' (ou $\D$-substantiel, s'il y a risque de confusion). Pour étayer l'intuition et simplifier les explications, un objet $F$ de $\D(I)$ ($\D$-coefficients sur $I$) sera représenté graphiquement par le diagramme `faible' ou `squelettique' correspondant -- on le regarde
\msp{page 74}
comme une sorte de (pâle) `squelette' du diagramme substantiel dont il provient. P. ex. si $I=\bDelta_0 : 0\longrightarrow 1$ \texttt{[plutôt $\bDelta_1$]}, on note un diagramme substantiel de type $(\bDelta_0)^o$ \texttt{[plutôt $(\bDelta_1)^o$]} par le même graphisme, que le diagramme ordinaire :
\[
F_0 \longleftarrow F_1.
\]
Si
\[
I = \begin{tikzcd}
a \arrow[r] & c \\
0 \arrow[u] \arrow[r] & b, \arrow[u]
\end{tikzcd}
\]
les diagrammes substantiels de type $I$ se notent
\[
\begin{tikzcd}
F_a \arrow[d] & F_c \arrow[l] \arrow[d] \\
F_0 & F_b. \arrow[l]
\end{tikzcd}
\]

\clearpage

\printedpage{2}
% head: A. Grothendieck : Dérivateurs, Ch. IV, Édition des Universités de Montpellier II et Paris VII
% foot: Version 2 février 2002

Il faudra bien sûr préciser chaque fois si ce graphisme désigne un diagramme squelettique, ou substantiel. Souvent on part d'un squelette, et on essaye de `l'étoffer' avec un peu de chair, pour en faire un diagramme substantiel -- si possible, de façon canonique. C'est là un point de philosophie dérivatique essentiel pour la pratique. On procédera souvent de proche en proche, en `substantifiant' des sous-diagrammes correspondants à une suite strictement croissante $I_0\subset I_1\subset\cdots\subset I_n=I$ de $I$, au mieux de la situation particulière.

\msp{page 75}

En voici un exemple typique, en supposant $\D$ un dérivateur \textit{ponctué} (les $\D(I)$ sont ponctuées, et les $f^*$ transforment objets nuls en objets nuls). Il s'agit de la définition du foncteur `cosuspension' ou `espace de lacets' dans $\cA$,
\[
\Omega : \cA \longrightarrow \cA.
\]
Il est entendu que pour $M$ dans $\cA$, $\Omega(M)$ doit s'insérer dans un carré commutatif
\[
\begin{tikzcd}
e \arrow[d] & \Omega M \arrow[l] \arrow[d] \\
M & e, \arrow[l]
\end{tikzcd}
\tag*{(C)}
\]
où $e$ désigne les éléments unités. Ce carré cependant devrait être substantiel, plus précisément, être un carré substantiel `$\D$-\textit{cartésien}'. Étant entendu qu'un objet de $\D(I)$ (carré substantiel dans $\cA$) est dit $\D$-cartésien, s'il est dans l'image essentielle de $\D(I_1)$ par $i_{1*}$, où
\[
i_1 : I_1 \hookrightarrow I,
\]
$I_1$ étant la sous-catégorie ouverte de $I$ définie par
\[
I_1 = \begin{tikzcd}
a \\
0 \arrow[u] \arrow[r] & b.
\end{tikzcd}
\]

\dgrblock{dgrb156}{Constructing the loop square without choosing homotopies}
\begin{dcomment}{exposition}{p75:En voici un exemple typique}{Constructing the loop square without choosing homotopies}
Grothendieck wants to obtain a functor \(\Omega\), not merely an object which happens to have the expected homotopy type. The two successive Kan extensions supply that functoriality. The direction of the arrows is governed by the presheaf convention: the indexing square has initial vertex \(0\) and final vertex \(c\), while its coefficient diagram has arrows pointing toward \(0\).

Let \(I_0=\{0\}\), let \(I_1=\{0,a,b\}\) with \(0\to a,0\to b\), and let \(I\) be the completed square. Denote the inclusions by \(i_0:I_0\to I_1\) and \(i_1:I_1\to I\), as in the text. Since \(I_0\) is a sieve in \(I_1\), the \emph{left} Kan extension in the presheaf convention is extension by the initial object outside \(I_0\). Pointedness identifies this initial object with zero. Hence
\[
 (i_0)_!M|_0=M,\qquad
 (i_0)_!M|_a=(i_0)_!M|_b=0.
\]
The pointwise assertion follows from empty comma categories at \(a,b\), and the assertion at \(0\) follows from full faithfulness of the inclusion. Thus the cospan \(0\to M\leftarrow0\) has been made coherent without choosing representatives or nullhomotopies.

Apply the right Kan extension \((i_1)_*\). Its value at \(c\) is the homotopy limit of that cospan. Set
\[
 \Omega M=c^*(i_1)_*(i_0)_!M.
\]
The object \((i_1)_*(i_0)_!M\) carries, at once, the two maps from \(\Omega M\) to zero and the specified coherent commutativity with the two zero maps to \(M\). It is cartesian by construction. Every morphism \(M\to N\) induces a morphism of these squares by the functoriality of the composite, so the construction really produces an endofunctor.

The stronger equivalence stated on page 77 follows with little extra work. Let \(\mathcal Q\subset\D(I)\) be the full subcategory of cartesian squares with values zero at \(a,b\). The functor
\[
 L=(i_1)_*(i_0)_!:\D(e)\longrightarrow\mathcal Q
\]
is fully faithful, since both Kan extensions are. It is essentially surjective as well. If \(Q\in\mathcal Q\), cartesianness gives \(Q\cong(i_1)_*i_1^*Q\). The counit
\((i_0)_!i_0^*i_1^*Q\to i_1^*Q\) is an isomorphism at \(0\), and at \(a,b\) both sides are zero. Der~2 therefore makes it an isomorphism. Substitution gives \(Q\cong L(Q_0)\), and restriction to \(0\) is a quasi-inverse.

The phrase \emph{unique à isomorphisme unique près} has its usual universal-property meaning: after specifying the identification of the input with \(M\), a comparison is unique when required to induce that identification. A square may have automorphisms induced by automorphisms of \(M\). Also, one should not add a requirement that the induced comparison fix an independently chosen identification of the fourth vertex. The theorem determines that vertex as part of the construction.

For complexes, the homotopy fibre of \(0\to M\) is \(M[-1]\), so this recovers the reduced loop functor, with no unwanted summand \(M\). This is the construction needed to repair the circle calculation of Chapter III.
\end{dcomment}


\msp{page 76}

La restriction du carré substantiel cherché (C) à la sous-catégorie $I_1$ est donc un diagramme substantiel de squelette
\[
\begin{tikzcd}
e \arrow[d] \\
M & e, \arrow[l]
\end{tikzcd}
\tag*{(Sq)}
\]
\printedpage{3}
% head: A. Grothendieck : Dérivateurs, Ch. IV, Édition des Universités de Montpellier II et Paris VII
% foot: Version 2 février 2002
dont il se déduit (à isomorphisme canonique près) en appliquant $i_*$ à ce dernier, de façon à le compléter en $\begin{tikzcd}
e \arrow[d] & \Omega M \arrow[l, dotted] \arrow[d, dotted] \\
M & e \arrow[l]
\end{tikzcd}$. Il reste donc à décrire une substantiation canonique du squelette (Sq), en termes de $M$. Cela nous conduit à regarder la sous-catégorie ouverte
\[
I_0 = \{0\} :\quad
\begin{tikzcd}
{} \arrow[r, dotted] & {} \\
0 \arrow[u, dotted] \arrow[r, dotted] & {} \arrow[u, dotted]
\end{tikzcd}
\]
de $I_1$ (donc aussi de $I$) réduite à l'objet initial, et l'inclusion
\[
I_0 \overset{i_0}{\hookrightarrow} I_1,
\]
et de substantifier le squelette (Sq) par $(i_0)_!(M)$.

On peut dire a posteriori que, si on connaissait le carré substantiel, il n'y a, à isomorphisme unique près, qu'une seule façon de l'obtenir comme un $i_{1*}i_{0!}(?)$, avec ? dans $\D(I_0)\simeq\cA$, car $i_{0!}$ et $i_{1*}$ étant pleinement fidèles, il en est de même de leur composé.

\msp{page 77}

En fait, on voit que le foncteur
\[
i_{1*}i_{0!} :\ \cA\simeq\D(I_0) \longrightarrow \D(I)
\]
induit une équivalence de $\cA$ avec la sous-catégorie pleine de $\D(I)$ formée des carrés $\D$-cartésiens dont les fibres en $a$ et $b$ sont nulles. (Le foncteur quasi-inverse évident est $i^*$, où $i : U_0 \hookrightarrow I$ \texttt{[plutôt $I_0 \hookrightarrow I$]}.) Cela signifie donc que dans ce cas, la donnée du squelette de nature très particulière ($F_a\simeq F_b\simeq 0_{\cA}$), plus la contrainte supplémentaire $\D$-cartésienne, sur le carré cherché, détermine celui-ci à isomorphisme unique près en termes de son squelette : il y a une substantiation non seulement `canonique' de celui-ci, mais même unique à isomorphisme unique près, à cause de la contrainte $\D$-cartésienne.

Je voudrais commencer par regarder le cas plus simple encore où $I=\bDelta_1=(0\longrightarrow 1)$, donc les squelettes envisagés sont simplement les flèches
\[
F_1 \longrightarrow F_0 \qquad \text{dans } \cA.
\]
La philosophie nous oblige, pour les flèches les plus fondamentales dans la théorie des dérivateurs $\D$ (flèches dans $\cA$, ou par extension dans une $\D(I)\simeq\cA_{\D_I}$, catégorie fondamentale du dérivateur induit $\D_I$) de se demander s'il n'y a pas

% ---- from Der-ChIV-567_p006-009.tex
\printedpage{4}
% head: A. Grothendieck : Dérivateurs, Ch. IV, Édition des Universités de Montpellier II et Paris VII
% foot: Version 2 février 2002
\msp{page 78}
une façon ‘canonique’, ‘naturelle’, ‘voulue par Dieu’, de le ‘remplumer’ en un diagramme substantiel. Les flèches les plus évidentes qui se présentent sont relatives à un foncteur
\[ f : I \longrightarrow J \]
dans Diag, et aux deux morphismes d'adjonction pour le couple $(f^*, f_*)$, savoir
\[
\begin{aligned}
G &\longrightarrow f_*f^*G && \text{dans } D(J)\\
f^*f_*F &\longrightarrow F && \text{dans } D(I). \quad (\text{\footnotemark})
\end{aligned}
\]
\footnotetext{Bien sûr, problème dual pour $(f_!, f^*)$.}

Je vais essayer de remplumer ces deux flèches, et vais généraliser en envisageant la situation ci-contre \texttt{[ci-dessous]}
\[
\begin{tikzcd}
I \arrow[d, "f"'] & F \\
J & G, \arrow[u, "u"']
\end{tikzcd}
\]
où
\[ \Hom^f(G,F) \overset{\text{déf}}{=} \Hom(f^*G,F) \overset{\text{adj.}}{\simeq} \Hom(G,f_*F). \]
Donc le même $u$ s'interprète par \textit{deux} flèches squelettiques
\[
\left\{
\begin{aligned}
f^*(G) &\longrightarrow F && \text{dans } \mathbf{D}(I)\\
G &\longrightarrow f_*(F) && \text{dans } \mathbf{D}(J).
\end{aligned}
\right.
\tag*{(1)}
\]
Les remplumer, c'est les interpréter (si faire se peut) en termes de deux objets (si possible canoniques) dans $\mathbf{D}(I \times \bDelta_1)$ et $\mathbf{D}(J \times \bDelta_1)$ respectivement.

Je vais me guider sur le cas particulier d'un dérivateur trivial
\[ \mathbf{D} = \mathbf{D}_{\mathcal{M}}, \]
où $\mathcal{M}$ est une catégorie de modèles ‘nue’ $(\Sigma = \emptyset)$.

\msp{page 79}
Pour $f$ fixé, on va construire une catégorie auxiliaire $K(f)$, de telle façon que pour toute catégorie $\mathcal{M}$, les systèmes $(G,F,u)$ d'un $\mathcal{M}$-préfaisceau $G$ sur $J$, d'un autre $F$ sur $I$, et d'un $f$-homomorphisme $G \longrightarrow F$, s'identifient (fonctoriellement) aux $\mathcal{M}$-préfaisceaux sur
\printedpage{5}
% head: A. Grothendieck : Dérivateurs, Ch. IV, Édition des Universités de Montpellier II et Paris VII
% foot: Version 2 février 2002
$K(f)$. Ceci dit, on va essayer d'interpréter les flèches (1) squelettiques, mais peu importe ici -- tous les $\mathbf{D}$-squelettes sont maintenant remplumés) en termes de la géométrie de $K(f)$. % [sic] remplumés) -> remplumés, closing parenthesis printed without any opening parenthesis

\[
\left\{
\begin{aligned}
&\Ob K(f) = \Ob K(I) \amalg \Ob K(J) \qquad \text{\texttt{[plutôt} $\Ob I \amalg \Ob J$\texttt{ ]}}\\
&\text{$I$, $J$ s'identifiant à des sous-catégories pleines de $K(f)$,}\\
&\text{il n'y a pas des flèches d'un $j \in J$ dans un $i \in I$.}\\
&\text{(Donc $I$ est ouverte dans $K(f)$, $J$ son complémentaire fermé.)}\\
&\text{Pour $i \in \Ob I$, $j \in \Ob J$, on pose $\Hom_{K(f)}(i,j) = \Hom_J(f(i),j)$.}
\end{aligned}
\right.
\]

La composition des flèches dans $K(f)$ est évidente. On va définir une rétraction canonique
\[ \bar{f} : K(f) \longrightarrow J \]
de $K(f)$ sur sa sous-catégorie fermée $J$. Donc $\bar{f}$ est connu sur $\Ob J$, et sur $\Ob I$ ce sera $f$ :
\[
\left\{
\begin{aligned}
\bar{f}(j) &= j && \text{si } j \in \Ob J\\
\bar{f}(i) &= f(i) && \text{si } i \in \Ob I.
\end{aligned}
\right.
\]
Il reste à définir, pour $x, y \in \Ob K(f)$,
\msp{page 80}
\[ \Hom_{K(f)}(x,y) \longrightarrow \Hom_J(\bar{f}(x), \bar{f}(y)). \]
Si $x, y \in \Ob J$, c'est l'identité, si $x, y \in I$, c'est $f$. Si $x \in I$, $y \in J$, on doit définir une flèche
\[ \Hom(f(x),y) \longrightarrow \Hom(f(x),y), \]
et on prend encore l'identité. On a
\[
\begin{tikzcd}
J \arrow[r, hook, "j"] \arrow[dr, equal] & K(f) \arrow[d, "\bar{f}"] & I \arrow[l, hook', "i"'] \\
& J &
\end{tikzcd}
\]
\[
\left\{
\begin{aligned}
\bar{f} \circ j &= \id_J\\
j \circ \bar{f} &= \left\{
\begin{aligned}
&j && \text{sur } J\\
&jf && \text{sur } I,
\end{aligned}
\right.
\end{aligned}
\right.
\qquad \text{i.e.} \quad
\boxed{
\begin{aligned}
j\bar{f}j &= j && \text{sur } J\\
j\bar{f}i &= jf && \text{sur } I
\end{aligned}
}\ ,
\]
et on a \texttt{[}$K = K(f)$\texttt{]}
\[
\begin{aligned}
\alpha :\ \id_K &\longrightarrow j\bar{f} && \text{i.e. homomorphisme fonctoriel}\\
x &\xrightarrow{\alpha(x)} j(\bar{f}(x)), && x \in \Ob K,
\end{aligned}
\]
\printedpage{6}
% head: A. Grothendieck : Dérivateurs, Ch. IV, Édition des Universités de Montpellier II et Paris VII
% foot: Version 2 février 2002
défini ainsi :
\[
\left\{
\begin{aligned}
&\text{Si } x \in J,\ \alpha(x) = \id_x\ :\ x \xrightarrow{\sim} j(\bar{f}(x)) = j(x) = x\\
&\text{Si } x \in I,\ \alpha(x) : \underbrace{x \longrightarrow f(x)}_{\text{dans } K(f)},\ \text{ou } \underbrace{f(x) \longrightarrow f(x)}_{\text{dans } J},\ \alpha(x) = \id_{f(x)}.
\end{aligned}
\right.
\]
On doit vérifier la fonctorialité pour $x$ variable dans $K$ \dots

\dgrblock{dgrb157}{The collage records a relative morphism}
\begin{dcomment}{exposition}{p79:La composition des flèches}{The collage records a relative morphism}
Grothendieck introduces \(K(f)\) to turn a morphism between coefficients on two different bases into an ordinary diagram on one base. This is a categorical form of the mapping cylinder, or more precisely the collage of the profunctor \(\Hom_J(f(-),-)\). The construction encodes the direction of the relative morphism in its mixed arrows.

For clarity, distinguish the copy of \(I\) from the copy of \(J\), even when they have the same objects. Internal arrows are those of \(I\) and \(J\); a mixed arrow \(x\to y\), with \(x\in I,y\in J\), is an arrow \(f(x)\to y\) in \(J\); there are no mixed arrows in the other direction. If \(a:x'\to x\) lies in \(I\), \(b:y\to y'\) lies in \(J\), and \(\beta:f(x)\to y\) represents a mixed arrow, composition is
\[
 \beta\circ a=\beta\circ f(a),\qquad
 b\circ\beta=b\beta
\]
in \(J\). Associativity follows directly from associativity in \(I,J\) and functoriality of \(f\). These formulas also prove that \(I\) is a sieve and \(J\) is its complementary cosieve.

Here is the universal property in full. Given presheaves \(F:I^o\to\mathcal M\), \(G:J^o\to\mathcal M\), and a natural transformation \(u:f^*G\to F\), define \(\Phi:K(f)^o\to\mathcal M\) to equal \(F\) and \(G\) on the two pieces. On the mixed arrow represented by \(\beta:f(x)\to y\), set
\[
 \Phi(\beta)=
 \bigl(G(y)\xrightarrow{G(\beta)}G(f(x))
       \xrightarrow{u_x}F(x)\bigr).
\]
Naturality of \(u\) proves compatibility with composition on the \(I\)-side; functoriality of \(G\) proves compatibility on the \(J\)-side. Conversely, restriction gives \(F,G\), and evaluation on the mixed arrow \(x\to f(x)\) represented by \(\mathrm{id}_{f(x)}\) gives \(u_x\). The two constructions are inverse.

A morphism between triples is a pair \(r:F\to F'\), \(s:G\to G'\) with \(r\,u=u'\,f^*s\). This equation is precisely naturality on the distinguished mixed arrows, and hence on every mixed arrow by the formula above. We have therefore obtained an isomorphism of categories, natural in \(\mathcal M\):
\[
 \mathrm{Fun}(K(f)^o,\mathcal M)
 \cong\{(F,G,u:f^*G\to F)\}.
\]
The isolated superscript \(o\) on the category of triples on manuscript page 82 is inconsistent with the preceding displayed universal property if the same variance is retained. One may either omit it in the modern formulation or dualize both sides. The printed reading is preserved in the transcription and recorded in the corrections table.

The retraction \(\bar f:K(f)\to J\) and transformation \(\alpha:\mathrm{id}\to j\bar f\) need no additional coherence data. In fact \(\bar f\dashv j\): for \(x\in K(f),y\in J\), the equality
\(\Hom_K(x,jy)=\Hom_J(\bar f x,y)\) follows from the defining cases. Its unit is precisely \(\alpha\), and its counit is the identity. This proves the naturality left for verification on page 80 and explains why the identities involving \(\bar f\) fit together.
\end{dcomment}


Soit $\Phi$ un $\mathcal{M}$-préfaisceau sur $K$. Il définit
\[
\begin{aligned}
F &= i^*(\Phi) && \in \Ob \mathbf{D}_{\mathcal{M}}(I)\\
G &= j^*(\Phi) && \in \Ob \mathbf{D}_{\mathcal{M}}(J).
\end{aligned}
\]
D'autre part $f^*(G)$ s'interprète comme $(j\bar{f})^*(\Phi)|I$, car
\[
\begin{aligned}
f^*(G)(i) &= G(\ \underbrace{f(i)}_{= \bar{f}(i) \,\in\, J}\ )\\
&= j^*(\Phi)(\bar{f}(i))\\
&= \Phi(j\bar{f}(i))\\
&= (j\bar{f})^*(\Phi)(i).
\end{aligned}
\]

\msp{page 81}
Donc
\[ f^*(G) \simeq (j\bar{f})^*(\Phi)|I. \]
Or la flèche
\[ \alpha : \id_K \longrightarrow j\bar{f} \]
donne
\[ \underbrace{\id_K(\Phi)}_{= \Phi} \xleftarrow{\Phi(\alpha)} (j\bar{f})^*(\Phi), \qquad \text{fonctoriel en } \Phi, \]
d'où, en prenant les restrictions à $I$,
\[ i^*(\Phi) \overset{\text{déf}}{=} F \xleftarrow{\ u\ } f^*G. \]

En fait, sans tellement se tortiller, on devrait pouvoir montrer que pour toute catégorie $\mathcal{M}$, la catégorie $\uHom(K^o, \mathcal{M})$ des $\mathcal{M}$-préfaisceaux sur $K$ est isomorphe à celle des triples $(F,G,u)$, avec $F \in \Ob \uHom(I^o, \mathcal{M})$, $G \in \uHom(J^o, \mathcal{M})$, et $u : f^*(G) \longrightarrow F$. Ainsi, $K$ apparaît comme solution d'un problème universel
\[ \uHom(K^o, \mathcal{M}) \longrightarrow \{F,\ G,\ u \text{ comme dessus}\}(\mathcal{M}). \]

Mais pour y préciser la flèche $\longrightarrow$ de façon conceptuelle simple, il est bon de se donner $K$ avec les structures supplémentaires \begin{tikzcd}[cramped, column sep=normal] I \arrow[r, hook, shift left, "j"] & K \arrow[l, shift left, "\bar{f}"] & I \arrow[l, hook', "i"'] \end{tikzcd}, $\bar{f}j = \id_J$, $\id_K \xrightarrow{\alpha} j\bar{f}$. Et ce système % [sic] I (first object, source of j and target of f-bar) -> J, since j : J -> K(f) and f-bar : K(f) -> J
\printedpage{7}
% head: A. Grothendieck : Dérivateurs, Ch. IV, Édition des Universités de Montpellier II et Paris VII
% foot: Version 2 février 2002
\msp{page 82}
$\{K,\ i,\ j,\ \bar{f},\ \alpha\}$ est (sauf erreur) défini à isomorphisme canonique près par la condition de représenter le foncteur $\mathcal{M} \longmapsto \{F,\ G,\ u\}^o(\mathcal{M})$ \dots

En particulier, le système défini par $G \in \Ob \mathbf{D}_{\mathcal{M}}(J)$
\[ (f^*(G),\ G,\ \id : f^*(G) \longrightarrow f^*(G)) \]
est défini sauf erreur par $\bar{f}^*(G)$ sur $K(f)$. On se propose [?] de trouver une interprétation similaire pour
\[ (F,\ f_*(F),\ f^*f_*F \longrightarrow F) \]
$(F \in \Ob \mathbf{D}(I))$. Le candidat naturel serait ici $\Phi = i_*(F)$. Admettons ce point (i.e. surtout la vérification du fait que $j^*(i_*(F)) \simeq f_*(F)$). Tout ceci montre que si $\mathbf{D}$ est à nouveau un dérivateur quelconque, on peut, par les formules précédentes, trouver un objet canonique $\Phi$ de $\mathbf{D}(K)$, en termes soit d'un $\mathbf{D}$-coefficient $F$ sur $I$ (en prenant $\Phi = i_*(F)$), soit d'un $\mathbf{D}$-coefficient $G$ sur $J$ (en prenant $\Phi = \bar{f}^*(G)$). D'autre part, on peut espérer définir canoniquement
\[
\begin{aligned}
\mathbf{D}(K) &\longrightarrow \mathbf{D}(I \times \bDelta_1)\\
\mathbf{D}(K) &\longrightarrow \mathbf{D}(J \times \bDelta_1)
\end{aligned}
\]
dont les composés avec les foncteurs ‘squelettes’
\[
\begin{aligned}
\mathbf{D}(I \times \bDelta_1) &\longrightarrow \uHom(\bDelta_1^o, \mathbf{D}(I))\\
\mathbf{D}(J \times \bDelta_1) &\longrightarrow \uHom(\bDelta_1^o, \mathbf{D}(J))
\end{aligned}
\]
\msp{page 83}
redonnent (pour $\Phi \in \Ob \mathbf{D}(K)$) les flèches squelettiques
\[
\left\{
\begin{aligned}
f^*(G) &\xrightarrow{\ u\ } F\\
G &\xrightarrow{\ v\ } f_*F
\end{aligned}
\right.
\]
associées à $(F,G,u)$. De façon précise, $f^*(G) \longrightarrow F$ se définit en termes de $\mathbf{D}$, et $\Phi \in \Ob \mathbf{D}(K)$, comme ci-dessus, et $v$ s'en déduit par adjonction.

\dgrblock{dgrb158}{The missing Kan extension comparison}
\begin{dcomment}{exposition}{p82:Le candidat naturel serait ici}{The missing Kan extension comparison}
Grothendieck's candidate \(i_*F\) is correct, and the comparison \(j^*i_*F\cong f_*F\) follows directly from the geometry of the mixed arrows. This is the point at which the collage begins to do mathematical work: a right Kan extension along an arbitrary functor is recovered from one along an inclusion.

For \(y\in J\), the comma category \((i/jy)\) has objects \(x\in I\) together with arrows \(i(x)\to j(y)\) in \(K(f)\). By definition, such an arrow is an arrow \(f(x)\to y\) in \(J\). Morphisms are the same commuting triangles in both descriptions. Hence there is a canonical isomorphism
\[
 (i/jy)\cong(f/y)
\]
commuting with the projections to \(I\), naturally in \(y\). The pointwise right Kan extension axiom in the presheaf convention computes both
\((i_*F)_{j(y)}\) and \((f_*F)_y\) as the homotopy limit of \(F\) over \((f/y)^o\). The comparison is therefore an isomorphism at every \(y\), and Der~2 makes it an isomorphism in \(\D(J)\).

There is a precise canonical map behind this pointwise calculation. The transformation
\(i\to jf\) induced by \(\alpha\) yields
\[
 f^*j^*i_*F\longrightarrow i^*i_*F\cong F.
\]
Its mate is \(j^*i_*F\to f_*F\). The comma-category identification proves that this mate is the isomorphism just computed. Consequently the displayed arrow on \(I\), after that identification, is exactly the counit \(f^*f_*F\to F\); agreement on objects is not being used as a substitute for agreement of the relevant natural transformations.

The other candidate is immediate but equally canonical. If \(\Phi=\bar f^*G\), then \(j^*\Phi=G\), \(i^*\Phi=f^*G\), and the transformation on \(I\) is the identity, since \(\bar f(\alpha_x)=\mathrm{id}_{\bar f(x)}\). These assertions hold for a derivator whenever the required indexing shapes are permitted. They use only restriction, right Kan extension, and the prederivator action on natural transformations.

In modern language the comma square expressing this argument is homotopy exact. The general calculus of such squares, with the comparison specified as a mate, is developed in \citet[Proposition 1.26]{Groth2013}; the relevant earlier convention was fixed in the commentary on Chapter~I, manuscript p.~31 (\emph{The base-change map as a mate}). The argument above completes the local proof rather than appealing to a stronger base-change statement with unchecked hypotheses.
\end{dcomment}


\dgrblock{dgrb159}{Completing both coherent adjunction arrows}
\begin{dcomment}{development}{p83:et $v$ s'en déduit par adjonction.}{Completing both coherent adjunction arrows}
Grothendieck leaves the chapter just after finding the characteristic map \(K(f)\to[1]\). His construction can be completed explicitly. The point is to use that map together with the retraction to \(J\), rather than to select a coherent representative of an arbitrary arrow in a homotopy category. We give the completion because it explains why both members of an adjoint pair can be treated by the same category.

\textbf{Collage completion.}\ours\
Let \(\D\) be a derivator in the presheaf convention, and let \(f:I\to J\). Assume its diagram category contains \(K=K(f)\), the products \(I\times[1]\), \(J\times[1]\), and the comma categories needed for the indicated right Kan extensions. There are functors
\[
 H:\D(K)\longrightarrow\D(I\times[1]),\qquad
 V:\D(K)\longrightarrow\D(J\times[1])
\]
whose skeletal arrows are, respectively, \(u:f^*G\to F\) and its adjoint \(v:G\to f_*F\), for \(F=i^*\Phi\), \(G=j^*\Phi\), and \(u\) induced by \(\alpha\).

\emph{Proof.}
Define \(h:I\times[1]\to K\) by
\[
 h(x,0)=i(x),\qquad h(x,1)=j(f(x)).
\]
On an arrow \(a:x\to x'\) at level \(0\) use \(i(a)\), and at level \(1\) use \(j(f(a))\). A mixed arrow \((x,0)\to(x',1)\), determined by \(a\), is sent to the mixed arrow represented by \(f(a)\). The definition of composition in the collage shows that this is a functor. Its vertical arrow at \(x\) is \(\alpha_{i(x)}\). Therefore \(H=h^*\) has restrictions \(F,f^*G\), and its arrow in the presheaf direction is exactly \(u:f^*G\to F\). This proves the first assertion without any Kan extension.

To construct the adjoint arrow, let \(\chi:K\to[1]\) be \(0\) on \(I\) and \(1\) on \(J\), and set
\[
 q=(\bar f,\chi):K\longrightarrow J\times[1],
 \qquad V=q_*.
\]
For \(y\in J\), the category \((q/(y,0))\) contains only objects from \(I\); its definition identifies it with \((f/y)\). Thus
\[
 (q_*\Phi)|_{J\times\{0\}}\cong f_*F.
\]
The category \((q/(y,1))\) has the terminal object given by \(j(y)\) and the identity of \((y,1)\). Indeed every object of that comma category specifies precisely one arrow to \(j(y)\), either an internal arrow of \(J\) or the mixed arrow represented by its \(J\)-component. In the presheaf convention the right Kan formula takes a homotopy limit over the opposite of this comma category, so the terminal object becomes an initial object and the limit is evaluation there. Consequently
\[
 (q_*\Phi)|_{J\times\{1\}}\cong G.
\]
These comparisons are natural, and their pointwise invertibility implies invertibility as coefficient diagrams.

It remains to identify the arrow, since its two endpoints alone would not determine it. Put \(E=q_*\Phi\), let \(t:E_1\to E_0\) be its skeletal arrow, and use the counit
\(\varepsilon:q^*E\to\Phi\).
Its restriction to \(J\) gives an isomorphism \(\lambda:E_1\to G\). Its restriction to \(I\) gives \(f^*E_0\to F\); let \(\theta:E_0\to f_*F\) be its mate. The preceding comma calculations say exactly that \(\theta\) is an isomorphism. Naturality of \(\varepsilon\) with respect to the mixed arrows \(x\to f(x)\) gives
\[
 (f^*E_0\longrightarrow F)\circ f^*t
 =
 u\circ f^*\lambda.
\]
After substituting the definitions of \(\theta\) and the counit of \(f^*\dashv f_*\), this says that
\(\theta\,t\,\lambda^{-1}:G\to f_*F\)
has adjoint \(u\). The adjunction bijection identifies it with \(v\). This proves the second assertion and the proposition. \(\square\)

Taking \(\Phi=\bar f^*G\), the first arrow is the identity on \(f^*G\), so the second is the unit \(G\to f_*f^*G\). Taking \(\Phi=i_*F\), the first is the counit \(f^*f_*F\to F\), while the second is the identity on \(f_*F\). Hence both adjunction transformations have canonical coherent realizations, natural in their coefficients. Applying the construction after a shift by \(A\) gives the corresponding statement in every \(\D(A)\); the compatibility with restriction in \(A\) follows from the product-square base change already used in Chapter III.

No strongness axiom was used. The input was the coherent \(\Phi\), and for the unit and counit that input was itself constructed by a specified Kan extension or restriction. The result does not supply a canonical section of the skeleton functor on all arrows.

The collage and the mate calculus are standard constructions. We checked the treatments of homotopy exact squares in \citet[Section 1]{Groth2013}, enriched derivator adjunctions in \citet[Sections 8--9]{GrothPontoShulman2014Additivity}, and pointwise Kan extension through modules in \citet[Section 5]{RiehlVerity2017}. We did not find in these sources the precise pair \(h^*,q_*\) written as a completion of this unfinished chapter. The marker above identifies our explicit completion and proof; it makes no claim of priority.
\end{dcomment}


L'idée naturelle ici, serait de ramener le problème de l'étoffement des flèches d'adjonction sqelettiques, au cas des immersions ouvertes ou fermées. Notons que les sous-catégories ouvertes de $K$ correspondent justement aux foncteurs $K \longrightarrow \bDelta_1$, en prenant les images inverses de l'ouvert générique. Il y a donc un unique % [sic] sqelettiques -> squelettiques, missing u (spelled squelettiques elsewhere in the text)
\[ \varphi :\ K(f) \longrightarrow \bDelta_1 \]
avec $I = \varphi^{-1}(\{0\})$. Donc \texttt{[fin du chapitre]}

\dgrblock{dgrb160}{A diagonal and its two projections}
\begin{dcomment}{examples}{p83:L'idée naturelle ici, serait de ramener}{A diagonal and its two projections}
Grothendieck's geometric construction can be seen completely in the case where \(I=\{a,b\}\) is discrete, \(J=e\), and \(f\) is the unique functor. Here \(K(f)\) is the cospan \(a\to *\leftarrow b\). Its presheaf coefficients have the opposite shape: an object \(\Phi\) is a coherent pair of arrows from \(G\) to \(F_a,F_b\).

The right Kan extension \(f_*F\) is the product \(F_a\times F_b\) in \(\D(e)\). Indeed Der~1 identifies \(\D(I)\) with \(\D(e)\times\D(e)\); the right adjoint to the diagonal restriction is the categorical product. The map \(q:K(f)\to[1]\) sends \(a,b\) to \(0\) and \(*\) to \(1\). Its right Kan extension sends the coherent pair to the single coherent arrow
\[
 G\longrightarrow F_a\times F_b
\]
whose two components are the original arrows. This proves directly, in this example, that the second construction is the adjoint of the first.

For \(\Phi=\bar f^*G\), both original arrows are identities. The output is the diagonal \(G\to G\times G\), as a specified object of \(\D([1])\). For \(\Phi=i_*F\), the value at \(*\) is \(F_a\times F_b\), and restriction along \(h\) gives the two projection arrows. The same collage has therefore produced the unit and counit of the diagonal-product adjunction, including their coherent representatives.

The empty case checks the direction of every construction. If \(I=\varnothing\), then \(K(f)=J\), \(q\) is the level-\(1\) inclusion, and \(q_*\Phi\) has value \(\Phi\) at level \(1\) and the terminal coefficient at level \(0\). Its presheaf arrow is \(\Phi\to1\), the unit of the adjunction between restriction to the empty category and its right Kan extension. Interchanging the two endpoints would give the wrong arrow even in this simplest case.
\end{dcomment}
\dgrchapter{V}{Chapter V}{Catégories de chemins et localisation}
\dgrpart{0}{III-V}{Manuscript pages 1--52}{III-V\enspace\textperiodcentered\enspace Chapter V\enspace\textperiodcentered\enspace manuscript pages 1--52}{Chapter V \textperiodcentered\ \textup{III-V}}
% ---- from Der-ChV-234_p001-005.tex
\printedpage{}
% head: (none); no page number printed on this page

\begin{center}
\vspace*{6cm}

{\LARGE\textsf{LES DÉRIVATEURS}}

\vspace{1.5cm}

{\large\textbf{ALEXANDRE GROTHENDIECK}}

\vspace{2.5cm}

{\large\textbf{Chapitre V}}

\vspace{0.4cm}

{\large\textbf{Catégories de chemins et localisation}}
\end{center}

\clearpage

\printedpage{}
% head: (none); no page number printed on this page

\vspace*{6cm}

Ce texte a été déchiffré et transcrit en \LaTeXe{} par M. Künzer. Il a été édité par M. Künzer, J. Malgoire, et G. Maltsiniotis. La transcription est aussi fidèle que possible au manuscrit. Pour les quelques corrections évidentes, ou rares commentaires des éditeurs, ainsi que pour la numérotation originale des pages du manuscrit, les caractères de machine à écrire \texttt{[typewriter]} entre crochets sont utilisés. Un point d'interrogation entre crochets signifie que l'on n'est pas sûr du mot qui précède.

Cette édition est provisoire. Les remarques, commentaires et corrections sont bienvenus. Envoyer un message à :

\textbf{maltsin@math.jussieu.fr}\\
\hspace*{\fill}G. Maltsiniotis\hspace*{2cm}

\clearpage

\printedpage{1}
% head: \textbf{A. Grothendieck :} \textit{Dérivateurs,} Ch. V, Édition des Universités de Montpellier II et Paris VII
% foot: Version 9 septembre 2001

\msp{page 1}

\section*{1\quad La catégorie des chemins}

Soit $\Phi$ la catégorie suivante :
\[
\left\{
\begin{array}{rcll}
\Ob\Phi & = & \mathbf{Z} & \\
\Fl\Phi & = & \Ob\Phi\sqcup\mathbf{Z} & \qquad (\Ob\Phi\leftrightarrow\text{identités})
\end{array}
\right.
\]
On note $\{\alpha\}$, ou simplement $\alpha$, l'objet qui correspond à $\alpha\in\mathbf{Z}$. On note $u_\alpha$ la flèche qui correspond à l'entier $\alpha\in\mathbf{Z}$.
\[
\begin{array}{rcl@{\qquad}l@{\quad}l@{\quad}l}
s(u_\alpha) & = & \alpha & \text{si } \alpha \text{ pair}, & \alpha+1 & \text{si } \alpha \text{ impair}\\[1ex]
b(u_\alpha) & = & \alpha+1 & \text{si } \alpha \text{ pair}, & \alpha-1 & \text{si } \alpha \text{ impair}
\end{array}
\]
\[
\texttt{[plutôt}\qquad b(u_\alpha) \;=\; \alpha+1 \quad \texttt{si}\ \alpha\ \texttt{pair,}\quad \alpha \quad \texttt{si}\ \alpha\ \texttt{impair]}
\]

Il n'y pas à définir de compositions, car si deux flèches sont composables, elles sont l'identité. En fait, $\Phi$ est une catégorie \textit{ordonnée}, $\{\alpha\}<\{\beta\}$ si et seulement si $\alpha$ pair, $\beta=\alpha+1$ \texttt{[plutôt} $\beta=\alpha+1$ \texttt{ou} $\beta=\alpha-1$\texttt{]}. % [sic] Il n'y pas -> Il n'y a pas, verb "a" missing

On note, pour $n\in\mathbf{N}$,
\[
\Phi_n\subseteq\Phi,\quad \Ob\Phi_n=\{\alpha\in\Ob\Phi\mid -n\le\alpha\le+n\},
\]
la sous-catégorie pleine de $\Phi$ dont les objets sont les $\{\alpha\}$ avec $-n\le\alpha\le+n$. Ainsi
\[
\Phi_1=\{-1\longleftarrow 0\longrightarrow 1\},
\]
la duale est
\[
\Psi_1=\Phi_1^{o}=\{-1\longrightarrow 0\longleftarrow 1\}.
\]
On peut considérer $\Phi$ comme somme amalgamée d'une suite ordonnée infinie (de deux côtés) de copies de $\Phi_1$, ou aussi d'une suite similaire de copies de $\Psi_1$.
\[
\begin{tikzcd}[column sep=small, execute at end picture={
\draw[decorate,decoration={brace,amplitude=4pt}] ([yshift=10pt]\tikzcdmatrixname-1-2.center) -- ([yshift=10pt]\tikzcdmatrixname-1-4.center) node[midway,above=4pt]{$\scriptstyle\Phi_1$};
\draw[decorate,decoration={brace,amplitude=4pt}] ([yshift=10pt]\tikzcdmatrixname-1-4.center) -- ([yshift=10pt]\tikzcdmatrixname-1-6.center) node[midway,above=4pt]{$\scriptstyle\Phi_1$};
\draw[decorate,decoration={brace,amplitude=4pt}] ([yshift=10pt]\tikzcdmatrixname-1-6.center) -- ([yshift=10pt]\tikzcdmatrixname-1-8.center) node[midway,above=4pt]{$\scriptstyle\Phi_1$};
\draw[decorate,decoration={brace,amplitude=4pt}] ([yshift=10pt]\tikzcdmatrixname-1-8.center) -- ([yshift=10pt]\tikzcdmatrixname-1-10.center) node[midway,above=4pt]{$\scriptstyle\Phi_1$};
\draw[decorate,decoration={brace,amplitude=4pt}] ([yshift=10pt]\tikzcdmatrixname-1-10.center) -- ([yshift=10pt]\tikzcdmatrixname-1-12.center) node[midway,above=4pt]{$\scriptstyle\Phi_1$};
\draw[decorate,decoration={brace,mirror,amplitude=4pt}] ([yshift=-10pt]\tikzcdmatrixname-1-3.center) -- ([yshift=-10pt]\tikzcdmatrixname-1-5.center) node[midway,below=4pt]{$\scriptstyle\Psi_1$};
\draw[decorate,decoration={brace,mirror,amplitude=4pt}] ([yshift=-10pt]\tikzcdmatrixname-1-5.center) -- ([yshift=-10pt]\tikzcdmatrixname-1-7.center) node[midway,below=4pt]{$\scriptstyle\Psi_1$};
\draw[decorate,decoration={brace,mirror,amplitude=4pt}] ([yshift=-10pt]\tikzcdmatrixname-1-7.center) -- ([yshift=-10pt]\tikzcdmatrixname-1-9.center) node[midway,below=4pt]{$\scriptstyle\Psi_1$};
\draw[decorate,decoration={brace,mirror,amplitude=4pt}] ([yshift=-10pt]\tikzcdmatrixname-1-9.center) -- ([yshift=-10pt]\tikzcdmatrixname-1-11.center) node[midway,below=4pt]{$\scriptstyle\Psi_1$};
}]
\cdots & {} & -4 \arrow[l] \arrow[r] & -3 & -2 \arrow[l] \arrow[r] & -1 & 0 \arrow[l] \arrow[r] & 1 & 2 \arrow[l] \arrow[r] & 3 & 4 \arrow[l] \arrow[r] & {} & \cdots
\end{tikzcd}
\]

\dgrblock{dgrb161}{A path remembers the arrows that will be inverted}
\begin{dcomment}{exposition}{p1:comme somme amalgamée d'une suite ordonnée infinie}{A path remembers the arrows that will be inverted}
Grothendieck is preparing a geometric description of localization. A morphism in a category of fractions is represented by a finite word of arrows, some traversed in their original direction and some traversed backwards. The alternation imposed here is a way of giving all such words a common shape. Inserting identity arrows makes any finite zigzag alternating; adding identities at its ends gives an interval symmetric about zero. The identities are bookkeeping, whereas the arrows traversed backwards carry the mathematical restriction imposed by the chosen class of denominators.

After the edition's correction of the target of $u_\alpha$, the category $\Phi$ is the poset on $\mathbf Z$ whose only strict relations are
\[
  2i<2i-1,\qquad 2i<2i+1.
\]
Every even vertex is a source and every odd vertex a sink. There is no composite of two nonidentity arrows. The sentence saying that two composable arrows ``sont l'identit\'e'' must therefore mean that at least one of them is an identity. Composing an identity with an arbitrary arrow supplies the elementary exception to its literal wording.

The subcategory $\Phi_n$ has $2n$ edges, not $n$. Thus the manuscript's later references to the parity of a path concern the index $n$, or half its number of edges. This distinction explains the repair on manuscript page~14. A path of index~1 has shape $\cdot\leftarrow\cdot\rightarrow\cdot$; two such paths joined end to end have the opposite orientation from $\Phi_2$. The correct assertion at the end of the present page is that $\Phi_n$, for even $n>0$, is assembled from copies of $\Psi_1=\Phi_1^o$. The printed $\Psi_n$ has the wrong orientation. We keep that reading in the source and use the corrected statement in the mathematics.

The nerve of every finite $\Phi_n$ is an ordinary subdivided interval. Nevertheless the category itself, including the directions of its arrows, matters: a functor from it to $X$ is an actual zigzag in $X$, and its natural transformations are actual commutative ladders. Passing immediately to the geometric interval would discard precisely the information needed to distinguish denominators from the other arrows.
\end{dcomment}


Les $\Phi_n$, pour $n$ impair, sont de même somme de copies de $\Phi_1$ mises bout à bout, et les $\Psi_n$, pour $n$ pair $>0$, somme de copies de $\Psi_1$.

\clearpage

\printedpage{2}
% head: \textbf{A. Grothendieck :} \textit{Dérivateurs,} Ch. V, Édition des Universités de Montpellier II et Paris VII
% foot: Version 9 septembre 2001

Les $\Phi_n$ forment un système inductif avec foncteurs de transition pleinement fidèles, dont $\Phi$ est la limite inductive
\[
\Phi=\varinjlim_n\Phi_n.
\]
Mais pour la suite on va faire des $\Phi_n$ un système \textit{projectif}, en définissant pour tout $n\le m$ une \textit{rétraction} $\varphi_{n,m}$ de $\Phi_m$ sur $\Phi_n$, par

\msp{page 2}
\[
\left\{
\begin{array}{l}
\varphi_{n,m}:\Phi_m\longrightarrow\Phi_n\subset\Phi_m\\[1ex]
\varphi_{n,m}(\{\alpha_i\})=
\left\{
\begin{array}{rcll@{\quad}l@{\ \ }r}
\alpha_i & = & i & \in\Ob\Phi_n, & \text{i.e.} & -n\le i\le n\\[1ex]
\alpha_i & = & n & & \text{si} & n\le i\le m\\[1ex]
\alpha_i & = & -n & & \text{si} & -m\le i\le -n
\end{array}
\right.
\end{array}
\right.
\]

\texttt{[plutôt}
\[
\begin{array}{rcl}
\varphi_{n,m}:\Phi_m & \longrightarrow & \Phi_n\subset\Phi_m\\[1ex]
\varphi_{n,m}(\{i\}) & = &
\left\{
\begin{array}{l@{\qquad}l@{\quad}l@{\ \ }r}
\{i\} & \texttt{si}\ i\in\Ob\Phi_n, & \texttt{i.e} & -n\le i\le n\\[1ex]
\{n\} & & \texttt{si} & n\le i\le m\\[1ex]
\{-n\} & & \texttt{si} & -m\le i\le -n\ \texttt{]}
\end{array}
\right.
\end{array}
\]
On note le système projectif d'ensembles ordonnés ainsi obtenu par
\[
\bPhi=(\Phi_n)_{n\in\mathbf{N}}\in\mathrm{Pro}(\mathrm{Cat}).
\]
Soit $X$ une catégorie, alors pour $n\in\mathbf{N}$
\[
\uCh_n(X)\overset{\text{déf}}{=}\uHom(\Phi_n,X)
\]
est la catégorie des diagrammes de $X$ de la forme
\[
A_{-n}\cdots\longleftarrow A_{-2}\xrightarrow{u_{-2}}A_{-1}\xleftarrow{u_{-1}}A_0\xrightarrow{u_0}A_1\xleftarrow{u_1}A_2\xrightarrow{u_2}A_3\longleftarrow\cdots A_n
\tag*{$(*)$}
\]
avec $2n$ flèches bout à bout $u_i$ $(-n\le i\le n)$ \texttt{[plutôt (}$-n\le i\le n-1$\texttt{)]}, deux flèches consécutives allant toujours en sens opposé,
\[
\text{soit}\ \ \cdot\longrightarrow\cdot\longleftarrow\cdot,\ \ \text{soit}\ \ \cdot\longleftarrow\cdot\longrightarrow\cdot\,.
\]
Le foncteur canonique déduit de $\varphi_{n,m}:\Phi_m\longrightarrow\Phi_n$
\[
\varphi_X^{m,n}:\uCh_n(X)\longrightarrow\uCh_m(Y),\qquad n\le m,
\]
\printedpage{3}
% head: \textbf{A. Grothendieck :} \textit{Dérivateurs,} Ch. V, Édition des Universités de Montpellier II et Paris VII
% foot: Version 9 septembre 2001
\texttt{[plutôt} $\uCh_m(X)$\texttt{]} est le `foncteur de dégénérescence', consistant à compléter un diagramme $(*)$, par $m-n$ termes $A_{-n}$ avec des identités à gauche, et par $m-n$ termes \texttt{[}$A_n$\texttt{]} avec des identités à droite, donc pour $n$ pair on trouve
\[
\underbrace{A_{-n}\cdots\xrightarrow{\text{id}}A_{-n}}_{m-n\text{ termes}}\xleftarrow{\text{id}}A_{-n}\longrightarrow A_{-(n-1)}\cdots\cdots A_{n-1}\longleftarrow A_n\xrightarrow{\text{id}}\underbrace{A_n\xleftarrow{\text{id}}A_n\longrightarrow\cdots A_n}_{m-n\text{ termes}}.
\]
C'est un foncteur pleinement fidèle (ce qui correspond au fait que $\varphi_{n,m}$ est surjectif sur les objets et sur les flèches. % [sic] flèches. -> flèches)., opening parenthesis never closed

\msp{page 3}
On pose
\[
\underbrace{\uCh_\infty(X)}_{\text{ou simplement }\uCh(X)}\qquad=\uHom(\bPhi,X)\overset{\text{déf}}{=}\varinjlim_n\uCh_n(X)
\]
(\footnote{L'ensemble des objets de cette catégorie se notera $\mathrm{Ch}_\infty(X)$ ou $\mathrm{Ch}(X)$, de même on définit les $\mathrm{Ch}_n(X)$.}) (\footnote{dans cette formule: $\bPhi$ procatégorie}). On a une inclusion pleinement fidèle
\[
\uCh_\infty(X)\lhook\joinrel\longrightarrow\uHom(\Phi,X),
\]
où la deuxième catégorie est celle des diagrammes infinis (dans les deux directions)
\[
\cdots\longleftarrow A_{-2}\xrightarrow{u_{-2}}A_{-1}\xleftarrow{u_{-1}}A_0\xrightarrow{u_0}A_1\xleftarrow{u_1}A_2\xrightarrow{u_2}\cdots
\]
et où $\uCh_\infty(X)$ s'identifie à la sous-catégorie pleine formées des objets pour lesquels les flèches $u_i$ sont des identités, pour $|i|$ assez grand. % [sic] formées -> formée, agreement with "la sous-catégorie"

On appellera $\uCh_\infty(X)$ la \textit{catégorie des chemins} dans $X$, ses objets seront les chemins dans $X$. Ceux qui proviennent de $\uCh_n(X)$ seront appelés \textit{chemins standards de longueur} $2n$. (Le terme `standard' s'applique ici à nos conventions de notations. Ainsi nous ne considérons que des chemins de longueur paire, et pour une longueur déterminée, nous choisissons l'une parmi les deux espèces raisonnables, par exemple pour les chemins de longueur 2 nous avons pris seulement
\[
A_{-1}\longleftarrow A_0\longrightarrow A_1
\]
et non pas
\[
A_{-1}\longrightarrow A_0\longleftarrow A_1.)
\]

\dgrblock{dgrb162}{The pro-interval and the eventually constant diagram}
\begin{dcomment}{exposition}{p3:ses objets seront les chemins dans}{The pro-interval and the eventually constant diagram}
Grothendieck wants paths of arbitrary finite length without making their length part of the final object. He obtains them by extending a finite diagram with identity arrows at both ends. The projective system of intervals and the inductive system of path categories describe the same operation from opposite sides.

Write $r_{n,m}:\Phi_m\to\Phi_n$ for the numerical clamping map
\[
 r_{n,m}(i)=\max(-n,\min(i,n)),\qquad n\leq m.
\]
It is a functor: an edge inside $[-n,n]$ is unchanged, and every other edge is sent to an identity. The equality $r_{n,m}r_{m,l}=r_{n,l}$ is literal, so precomposition gives a strict inductive system. Its transition functors are fully faithful, for a reason slightly more precise than mere surjectivity on objects and arrows. If $F,G:\Phi_n\to X$ and $\eta:Fr_{n,m}\to Gr_{n,m}$, naturality on each collapsed edge says that the two adjacent components of $\eta$ agree. They therefore agree throughout either constant tail. Restriction to $\Phi_n$ determines $\eta$ uniquely and yields an arbitrary natural transformation $F\to G$. This proves the required bijection on Hom sets.

Surjectivity alone would not prove this assertion for an arbitrary functor. The map from a discrete two-point category to the point is surjective on objects and arrows, while precomposition gives the diagonal $X\to X\times X$, which need not be full. What makes clamping work is the chain of edges sent to identities inside each fibre.

The filtered colimit is consequently realized as the full subcategory of $\operatorname{Fun}(\Phi,X)$ consisting of diagrams whose edges are identities outside a finite interval. A natural transformation between two such diagrams also becomes constant in both tails: choose one interval containing the nonidentity edges of both diagrams and apply naturality to the identity edges beyond it. Thus both objects and morphisms descend to a common finite stage. No assertion about commuting an arbitrary functor category with filtered colimits is being used here.

The notation $\operatorname{Fun}(\boldsymbol\Phi,X)$ is a convenient enriched version of the usual pro-object formula: it means $\mathop{\rm colim}_n\operatorname{Fun}(\Phi_n,X)$ as a category. Its set of objects is not a space of arbitrary infinite paths. The eventual identity condition is essential, and differs even from eventual invertibility.

Reflection $i\mapsto-i$ is the unique nonidentity automorphism of $\Phi_n$ when $n>0$, since an automorphism of the underlying finite path is either the identity or reflection and both preserve the alternating orientation. For $n=0$ the automorphism group is trivial. This harmless boundary case will also matter when the two endpoints coincide.
\end{dcomment}


\msp{page 4}
Je me rends compte que mes notations $u_i$ ne sont pas commodes pour la dualité, je vais les changer, en $u_i$ pour les flèches entre objets $\{i\}$ avec $i\ge0$, et $v_i$ la flèche symétrique de $u_i$ par rapport à l'origine. (Le groupe des automorphismes des $\Phi_n$ est réduit au groupe à deux éléments, le seul automorphisme non trivial étant la symétrie $\{i\}\longrightarrow\{-i\}$. Donc nos chemins seront notés désormais % [sic] {-i}. -> {-i}). , opening parenthesis "(Le groupe" not closed on this page
\[
A_{-2}\xrightarrow{v_1}A_{-1}\xleftarrow{v_0}A_0\xrightarrow{u_0}A_1\xleftarrow{u_1}A_2\longrightarrow\cdots.
\tag*{$(*)$}
\]

% ---- from Der-ChV-234_p006-010.tex
\printedpage{4}
% head: A. Grothendieck : Dérivateurs, Ch. V, Édition des Universités de Montpellier II et Paris VII
% foot: Version 9 septembre 2001

(Ce n'est pas non plus enthousiasmant \dots)

Soit
\[
\varepsilon=\{-1,+1\}
\]
l'ensemble ordonné \textit{discret,} réduit aux deux éléments $-1$, $+1$, et définissons des foncteurs d'inclusion (pleinement fidèles)
\[
\begin{aligned}
\alpha_n:\varepsilon\longrightarrow\Phi_n\quad & \alpha_n(-1)=-n\\
& \alpha_n(1)=n,
\end{aligned}
\]
ces foncteurs sont compatible avec les $\varphi_{n,m}$, d'où % [sic] compatible -> compatibles, accord avec « ces foncteurs »
\[
\alpha_\infty:\varepsilon\longrightarrow\boldsymbol{\Phi}\in\mathrm{ProCat},
\]
d'où pour tout $X$ un foncteur
\[
\kappa:\uHom(\boldsymbol{\Phi},X)\longrightarrow\uHom(\varepsilon,X)\simeq X\times X
\]
i.e.

\msp{page 5}
\[
\begin{aligned}
\uCh(X) &\longrightarrow X\times X\\
c &\longrightarrow (\sigma(x),\beta(x)).
\end{aligned}
\]

\texttt{[plutôt} $(\sigma(c),\beta(c))$\texttt{]} Pour tout chemin $c$ dans $X$ i.e. un diagramme $(*)$ (satisfaisant à la condition de dégénérescence à l'infini), on a une source ou origine du chemin, et un but ou une \textit{extrémité}
\[
\left\{
\begin{aligned}
\sigma(x) &= A_i\quad\text{pour } i\leq -\mathrm{long}(x)\\
\beta(x) &= A_i\quad\text{pour } i\geq \mathrm{long}(x)
\end{aligned}
\right.
\]
\texttt{[remplacer} $x$ \texttt{par} $c$\texttt{; pour long}$(c)$ \texttt{défini convenablement]}. La fibre du foncteur précédent en $x$, $y$ sera notée
\[
\uCh(X;x,y)\qquad\text{ou simplement }\uCh(x,y),
\]
\textit{catégorie des chemins de $x$} à \textit{$y$.}

\dgrblock{dgrb163}{Fixing the endpoints also fixes the vertical arrows}
\begin{dcomment}{examples}{p5:La fibre du foncteur précédent en}{Fixing the endpoints also fixes the vertical arrows}
Grothendieck's source--target functor allows a path to become a candidate morphism from a specified object to another. Its fibre over $(x,y)$ is a strict categorical fibre. An object is an eventually constant zigzag with tails $x$ and $y$; a morphism is a natural transformation whose two tail components are $1_x$ and $1_y$. This last condition is indispensable for the interpretation by a Hom set.

There are two different full-subcategory statements in these pages. The category of $\Sigma$-paths is full in the category of all paths, because it imposes a condition only on the objects. Likewise its fixed-endpoint fibre is full in the fixed-endpoint fibre of all paths. But a fixed-endpoint fibre is generally not full in the entire functor category. On pages~15--16 and~24 the expressions ``sous-cat\'egorie pleine'' must be read with this correction: the endpoint components of a natural transformation are required to be identities.

For a concrete test take the one-object category $BG$ of a nontrivial group $G$, and let $\Sigma$ contain only the identity. The paths representing $1$ and $g\ne1$ may be written, at index~2,
\[
 (1,1,1,1),\qquad (g,1,1,1),
\]
where the entries label the four alternating arrows. There is a natural transformation from the first diagram to the second with vertex components $(1,g,g,g,g)$. It satisfies all four naturality squares, but its final component is $g$, not the identity. Thus it is not a morphism in the strict fibre. If one used the full subcategory on objects with the prescribed endpoints, this transformation would identify $1$ with $g$ on connected components, although localization at identities leaves $BG$ unchanged.

For an ordinary arrow category, the same convention makes the fibre over $(x,y)$ discrete: a morphism from $f:x\to y$ to $g:x\to y$ with both vertical sides identities exists precisely when $f=g$, and is then unique. This is the precise assertion used on page~8.

The map $\varepsilon\to\Phi_n$ is fully faithful for $n>0$. At $n=0$ it sends both endpoint labels to the same object and is not fully faithful. It still defines the source--target functor, whose index-zero stage is the diagonal $X\to X\times X$. Consequently a path of index zero occurs in the fibre only when $x=y$, exactly as the identity construction requires.
\end{dcomment}


\printedpage{5}
% head: A. Grothendieck : Dérivateurs, Ch. V, Édition des Universités de Montpellier II et Paris VII
% foot: Version 9 septembre 2001

\section*{2\quad Calcul d'une catégorie de fractions $X\Sigma^{-1}$}

Avant de continuer le formalisme des chemins (composition etc.), je vais tester l'utilité de la notion de catégorie des chemins, pour le `calcul' des $X\Sigma^{-1}$, où
\[
\Sigma\subset\mathrm{Fl}(X),\qquad\Sigma\text{ contient les identités}
\]
est une partie autrement \textit{quelconque} de $\mathrm{Fl}(X)$. On sait que $\mathrm{Ob}\,X\Sigma^{-1}=\mathrm{Ob}\,X$, et tout revient à définir les
\[
\Hom_{X\Sigma^{-1}}(x,y),\qquad x,y\in\mathrm{Ob}\,X,
\]
\msp{page 6}
et pour $x,y,z\in\mathrm{Ob}\,X$ la composition
\[
\Hom_\Sigma(x,y)\times\Hom_\Sigma(y,z)\xrightarrow{(u,v)\mapsto vu}\Hom_\Sigma(x,z),
\]
(où on note $\Hom_\Sigma$ au lieu de $\Hom_{X\Sigma^{-1}}$, pour alléger les notations).

Grâce à $\Sigma$, on définit une sous-catégorie pleine
\[
\uCh_\Sigma(X)\hookrightarrow\uCh(X),
\]
(\footnote[3]{L'ensemble de ses objets est noté $\mathrm{Ch}(X)$. Ses éléments sont les $\Sigma$-chemins de $x$ à $y$ \texttt{[de} $x$ \texttt{à} $y$\texttt{?]}.}) par la condition sur $c\in\mathrm{Ch}(X)$ : Pour $i$ impair, \texttt{[plutôt pair]} i.e. $i-1\xleftarrow{u_i}i$ ($u_i$ est dite `rétrograde'), $c(u_i)\in\Sigma$, i.e. les `tronçons' rétrogrades dans le chemin $c$ doivent être dans $\Sigma$. On définit de même les sous-catégories pleines
\[
\uCh_\Sigma(X;x,y)=\{c\in\uCh(X;x,y)\mid c\in\uCh_\Sigma(X)\},
\]
ou encore la fibre en $(x,y)$ du foncteur source-but pour $\uCh_\Sigma(X)$
\[
\uCh_\Sigma(X)\longrightarrow X\times X.
\]

\dgrblock{dgrb164}{Why a commutative ladder gives equality after localization}
\begin{dcomment}{exposition}{p6:les `tronçons' rétrogrades dans le chemin}{Why a commutative ladder gives equality after localization}
Grothendieck places no calculus-of-fractions hypothesis on $\Sigma$: it is any set of arrows containing the identities. Its closure under composition is unnecessary for the path construction. The hypothesis concerns only the backwards horizontal arrows; the vertical arrows of a morphism of paths may be arbitrary arrows of $X$.

To see why this is sufficient, let $Q:X\to X[\Sigma^{-1}]$ be the localization and consider a square around a backwards edge,
\[
\begin{tikzcd}[column sep=large]
 A_{i-1}\arrow[d,"a_{i-1}"'] & A_i\arrow[l,"s"']\arrow[d,"a_i"]\\
 B_{i-1} & B_i\arrow[l,"t"']
\end{tikzcd}
\qquad s,t\in\Sigma.
\]
Its commutativity gives $Q(a_{i-1})Q(s)=Q(t)Q(a_i)$, hence
\[
 Q(a_i)Q(s)^{-1}=Q(t)^{-1}Q(a_{i-1}).
\]
The square therefore remains commutative after reversing its horizontal arrows. Apply this to every backwards edge and then compose the resulting forward squares. The endpoint components give a commutative square between the two composite arrows in the localization. In a strict endpoint fibre those components are identities, so the two composites are equal.

This proves the functor to the arrow category and the induced function on $\pi_0$ without assuming that any vertical arrow becomes invertible. It also explains the deliberate use of connected components: an equality may be witnessed by a finite zigzag of natural transformations rather than a single transformation. The later reparametrization argument is designed to produce just such zigzags.

The parity notation changes during the construction. It is safest to number the edges consecutively from the left whenever discussing a fixed even-index interval; its even-numbered edges are then the backwards ones. This numerical convention must not be confused with the original label $u_\alpha$ on the doubly infinite interval, or with the symmetric $u_i,v_i$ notation introduced on page~4. The invariant definition is simply that an edge traversed right-to-left belongs to $\Sigma$.
\end{dcomment}


On définit une application canonique
\[
\mathrm{Ch}_\Sigma(X,x,y)\longrightarrow\Hom_\Sigma(x,y),
\]
\msp{page 7}
en associant `a tout chemin de $x$ dans $y$ le $\Sigma$-homomorphisme de $x$ dans $y$ obtenu en \textit{inversant} (dans $X\Sigma^{-1}$) les tronçons rétrogrades (ce qui est possible, puisqu'ils sont dans $\Sigma$), de sorte qu'on trouve dans $X\Sigma^{-1}$ un `chemin direct' % [sic] `a -> à, accent imprimé comme guillemet ouvrant
\[
x=A_{-n}\longrightarrow\cdots\longrightarrow A_{-1}\longrightarrow A_0\longrightarrow A_1\longrightarrow A_2\cdots A_n=A_{n+1}=\cdots=y.
\]
et on compose tous ces morphismes.

En fait, on trouve aussitôt que l'on définit ainsi un \textit{foncteur}
\[
\uCh_\Sigma(X)\longrightarrow\uFl(X\Sigma^{-1}),
\]
\texttt{[}$\uFl$ \texttt{catégorie des flèches]} du fait que si on a un diagramme

\printedpage{6}
% head: A. Grothendieck : Dérivateurs, Ch. V, Édition des Universités de Montpellier II et Paris VII
% foot: Version 9 septembre 2001
\[
\begin{tikzcd}
A_{i-1} \arrow[d, "f_{i-1}"] & A_i \arrow[l, "u_i"'] \arrow[d, "f_i"] \\
A'_{i-1} & A'_i \arrow[l, "u'_i"']
\end{tikzcd}
\qquad\qquad (i\text{ pair})
\]
avec $u_i$, $u'_i$ dans $\Sigma$, ce diagramme, en tant que diagramme dans $X\Sigma^{-1}$, est commutatif avec flèches horizontales des \textit{isomorphismes,} et reste donc commutatif en inversant lesdites flèches. D'où un foncteur induit
\[
\uCh_\Sigma(X;x,y)\longrightarrow\uFl(\underbrace{X\Sigma^{-1}}_{X'};x,y).
\]

Mais pour toute catégorie $X'$, et des objets $x$, $y$ dans $X'$, la catégorie $\uFl(X',x,y)$ est
\msp{page 8}
une catégorie \textit{discrète,} ayant comme ensemble d'objets $\Hom_{X'}(x,y)$. Ainsi le foncteur précédant définit en fait une application canonique % [sic] précédant -> précédent, adjectif
\[
\boxed{\pi_0\uCh_\Sigma(X;x,y)\longrightarrow\Hom_\Sigma(x,y).}
\]
Je dis que cette application est toujours bijective. Pour le voir, nous allons \textit{constuire} `from scratch' la catégorie $X\Sigma^{-1}$, en prenant comme objets ceux de $X$, et en \textit{définissant} % [sic] constuire -> construire, lettre omise
\[
\boxed{\Hom_\Sigma(x,y)=\pi_0\uCh_\Sigma(X;x,y).}
\]
Il faut donc définir la composition
\[
\Hom_\Sigma(x,y)\times\Hom_\Sigma(y,z)\longrightarrow\Hom_\Sigma(x,z).
\]
Pour ceci, on va définir des foncteurs `composition de chemins'
\[
\begin{aligned}
\uCh_m(X;x,y)\times\uCh_n(X;y,z) &\longrightarrow \uCh_{m+n}(X;x,z)\\
(c,c') &\longmapsto c'\circ c,
\end{aligned}
\]
mais en nous bornant pour simplifier au cas où $m$, $n$ ont \textit{la même parité,} de telle façon qu'en `mettant bout à bout' $c$ et $c'$ (dans cet ordre), on trouve bien un chemin de type $n+m$.

\dgrblock{dgrb165}{Localization without a short-fraction hypothesis}
\begin{dcomment}{context}{p8:Je dis que cette application est toujours bijective.}{Localization without a short-fraction hypothesis}
Grothendieck now proposes to recover the ordinary localization directly from categories of paths. The basic localization by composable words and inverse symbols is classical; compare \citet[Chapter~I, pp.~6--21]{GabrielZisman1967} and the explicit word formulation in \citet[\S3.1--3.3]{DwyerKan1980Simplicial}. His additional object of study is the category of presentations of a morphism, whose connected components recover those words modulo relations.

A short fraction such as $x\leftarrow x'\to y$ need not represent every localized morphism unless one imposes extra hypotheses. Arbitrary finite zigzags do. The use of all indices thus avoids an Ore condition, a model structure, and any saturation assumption on the original set $\Sigma$. Enlarging $\Sigma$ to the subcategory it generates changes neither the ordinary localization nor the final set of components, but the unlocalized path categories can change. They should not be silently identified.

There are three tasks in the proposed proof. First, concatenation must descend from particular finite representatives to their classes. Second, the original arrows of $X$ must determine a functor to those classes, and arrows in $\Sigma$ must become invertible. Finally, every functor inverting $\Sigma$ must factor uniquely by evaluating zigzags. The first task occupies most of the combinatorics. The second is where a generalized interval lemma is temporarily assumed. We supply its relative version below and also give direct witnesses for the particular identities needed here. The universal property then follows by an explicit argument at the end of the section.
\end{dcomment}


\msp{page 9}
Cette opération est `transposée' du couple de foncteurs
\[
\begin{aligned}
\Phi_n &\xrightarrow{\alpha} \Phi_{n+m}\quad i\longmapsto (i-m)\\
\Phi_m &\xrightarrow{\beta} \Phi_{n+m}\quad i\longmapsto (i+n),
\end{aligned}
\]
\printedpage{7}
% head: A. Grothendieck : Dérivateurs, Ch. V, Édition des Universités de Montpellier II et Paris VII
% foot: Version 9 septembre 2001
où on aura
\[
\alpha(\{n\})=\beta(\{-m\})=n-m,
\]
et on aura
\[
\Phi_{n+m}\simeq\Phi_n\sqcup_e\Phi_m,
\]
i.e. $\Phi_{n+m}$ est la `somme amalgamée' dans (Cat), obtenue en `recollant' $\Phi_n$ et $\Phi_m$, par identification de l'`élément but' de $\Phi_n$ avec l'`élément source' de $\Phi_m$. Le foncteur composition de chemins induit \texttt{[dans ce qui précède, à partir du début de la page 9, il faut inverser} $m$ \texttt{et} $n$\texttt{]}
\[
\uCh_{m,\Sigma}(x,y)\times\uCh_{n,\Sigma}(y,z)\longrightarrow\uCh_{m+n,\Sigma}(x,z),
\]
d'où en passant aux $\pi_0$
\[
\Hom_{\Sigma,m}(x,y)\times\Hom_{\Sigma,n}(y,z)\longrightarrow\Hom_{\Sigma,m+n}(x,z),
\]
où on dénote, pour $a,b\in\mathrm{Ob}\,X$, par
\[
\Hom_{\Sigma,n}(a,b)\overset{\text{déf}}{=}\pi_0(\uCh_{m,\Sigma}(a,b)).
\]

\texttt{[plutôt} $\uCh_{n,\Sigma}$\texttt{]}

\msp{page 10}
Pour en déduire une application sur les $\varinjlim$
\[
\underbrace{\Hom_\Sigma(x,y)}_{=\ \varinjlim\ \Hom_{\Sigma,m}(x,y)}\quad\times\quad\underbrace{\Hom_\Sigma(y,z)}_{=\ \varinjlim\ \Hom_{\Sigma,n}(y,z)}\quad\longrightarrow\quad\underbrace{\Hom_\Sigma(x,z)}_{=\ \varinjlim\ \Hom_{\Sigma,n+m}(x,z)}
\]
il faut prouver, pour $n\leq n'$, $m\leq m'$, la commutativité dans
\[
\begin{tikzcd}
\Hom_{\Sigma,m}(x,y)\times\Hom_{\Sigma,n}(y,z) \arrow[r] \arrow[d] & \Hom_{\Sigma,m+n}(x,y) \arrow[d] \\
\Hom_{\Sigma,m'}(x,y)\times\Hom_{\Sigma,n'}(y,z) \arrow[r] & \Hom_{\Sigma,m'+n'}(x,y)
\end{tikzcd}
\]
\texttt{[plutôt} $\Hom_{\Sigma,m+n}(x,z)$\texttt{,} $\Hom_{\Sigma,m'+n'}(x,z)$\texttt{]} En d'autres termes, que pour $c:x\longrightarrow y$ et $c':y\longrightarrow z$ deux $\Sigma$-chemins composables de longueur $m$ et $n$ respectivement, le composé $c'c:x\longrightarrow z$ ne change pas, `à homotopie près' dans $\uCh_\Sigma(x,z)$, quand on remplace $c$, $c'$ par des chemins $\bar{c}$, $\bar{c}'$ de longueur $m'$, $n'$ déduits de $c$ et $c'$ par `dégénérescence aux bouts'.

Plus précisement, % [sic] précisement -> précisément, accent omis
\[
\bar{c}'\bar{c}\sim\overline{c'c},
\]
le chemin composé est homotope à celui déduit de $c'c$ par `dégénérescence aux bouts'.

Je viens de regarder pour $n=m=1$, $n'=m'=2$, et de constater qu'il n'y a pas de flèche naturelle entre $\bar{c}'\bar{c}$ et $\overline{c'c}$ - l'homotopie cherchée doit
\printedpage{8}
% head: A. Grothendieck : Dérivateurs, Ch. V, Édition des Universités de Montpellier II et Paris VII
% foot: Version 9 septembre 2001
\msp{page 11}
se définir à l'aide d'un `chemin' dans $\uCh_{\Sigma,m',n'}(X;x,z)$. Dans le cas général, on va en fait trouver une homotopie
\[
h=h_{(n,m),(n',m')}
\]
`rendant commutatif' le diagramme de foncteurs
\[
\begin{tikzcd}[column sep=large]
\mathcal{A}=\uCh_{\Sigma,m}(X;x,y)\times\uCh_{\Sigma,n}(X;y,z) \arrow[r, "*_{m,n}"] \arrow[d, "i_{m,m'}\times j_{n,n'}"] & \uCh_{\Sigma,m+n}(X;x,z) \arrow[d, "k_{m+n,m'+n'}", ""'{name=K}] \\
\uCh_{\Sigma,m'}(X;x,y)\times\uCh_{\Sigma,n'}(X;y,z) \arrow[r, "*_{m',n'}", ""{name=S}] & \uCh_{\Sigma,m'+n'}(X;x,z)=\mathcal{B},
\arrow[from=K, to=S, bend right, "h"']
\end{tikzcd}
\]
où les $i_{\alpha,\beta}$, $j_{\alpha,\beta}$, $k_{\alpha,\beta}$ sont les foncteurs de transition
\[
\uCh_{\Sigma,\alpha}(X;a,b)\longrightarrow\uCh_{\Sigma,\beta}(X;a,b)
\]
$(\alpha\leq\beta)$ pour les $(x,y)$, $(y,z)$ et $(x,z)$ (en guise de $(a,b)$) respectivement. L'homotopie $h$ sera un chemin dans la catégorie $\uHom(\mathcal{A},\mathcal{B})$, allant de
\[
k_{m+n,m'+n'}\circ{*}_{m+n}\longrightarrow{*}_{m'+n'}\circ(i_{m,m'}\times j_{n,n'}).
\]
L'existence d'une telle homotopie signifie ni plus ni moins que le diagramme carré ci-dessus est commutatif, en tant que diagramme dans la catégorie (Cat\,hot) des `catégories à homotopie près,' dont les objets sont
\msp{page 12}
les (petites) catégories, et où on définit
\[
\Hom_{\mathrm{Cat\,hot}}(\mathcal{A},\mathcal{B})\overset{\text{déf}}{=}\underbrace{\pi_0(\uHom(\mathcal{A},\mathcal{B}))}_{\begin{array}{@{}l@{}}\text{classes d'homotopie}\\ \text{de foncteurs de }\mathcal{A}\\ \text{dans }\mathcal{B}\end{array}}\,.
\]
C'est cette conclusion qui sera pour nous utile, et non la description explicite du chemin $h$. Cette existence admise (cf. ci-dessous la preuve), on voit donc que si on regarde, pour $a,b\in\mathrm{Ob}\,X$, le système des catégories $\uCh_{\Sigma,r}(X;a,b)$ $(r\in\mathbf{N})$ comme un \textit{ind-objet de} Cat\,hot\textit{,} qu'on désigne par
\[
\uCh^{\mathrm{ind}}_{\Sigma}\,(X;a,b)\in\mathrm{Ob}\,\mathrm{Ind}(\mathrm{Cat\,hot}),
\]
on aura un morphisme canonique dans Ind(Cat\,hot)
\[
\uCh^{\mathrm{ind}}_{\Sigma}\,(X;x,y)\times\uCh^{\mathrm{ind}}_{\Sigma}\,(X;y,z)\xrightarrow{\ *\ }\uCh^{\mathrm{ind}}_{\Sigma}\,(X;x,z),
\]
appelé \textit{`ind-morphisme de composition de chemins.'} On aurait aimé le déduire d'un honnête \textit{foncteur} `composition de chemins'
\[
\begin{aligned}
\uCh_\Sigma(X;x,y)\times\uCh_\Sigma(X;y,z) &\xrightarrow{\ ?\ } \uCh_\Sigma(X;x,z),\\
\text{qu'on noterait }(c,c') &\longmapsto c'\circ c\text{ ou }c'\cdot c,
\end{aligned}
\]
mais je doute qu'un tel foncteur existe,

\dgrblock{dgrb166}{What composition in an ind-category actually asserts}
\begin{dcomment}{status}{p12:On aurait aimé le déduire d'un honnête}{What composition in an ind-category actually asserts}
Grothendieck has reached a useful level of coherence before obtaining a single composition functor on the eventual path categories. For even indices $m,n$, concatenation gives an actual functor at a finite stage. Increasing either index changes its comparison with the larger-stage concatenation by a finite zigzag of natural transformations. This is exactly the relation imposed in $\mathrm{Cat\,hot}$, where
\[
 \operatorname{Hom}_{\mathrm{Cat\,hot}}(A,B)
   =\pi_0\operatorname{Fun}(A,B).
\]
Products descend to this category because one can multiply each natural transformation in a zigzag by an identity transformation.

For ind-objects the assertion is governed by the formula
\[
 \operatorname{Hom}_{\operatorname{Ind}(C)}
       ((A_i),(B_j))
   =\lim_i\mathop{\rm colim}_j\operatorname{Hom}_C(A_i,B_j).
\]
Take $A_{m,n}=\operatorname{Ch}_{\Sigma,m}(x,y)\times
\operatorname{Ch}_{\Sigma,n}(y,z)$ and $B_r=\operatorname{Ch}_{\Sigma,r}(x,z)$, with all indices even. The stagewise concatenation maps land in $B_{m+n}$. Their compatibility as elements of the displayed inverse limit is precisely what the reparametrization lemma proves. The passage to components then gives an ordinary associative composition law.

A functor between the ordinary filtered unions would be additional data. Even if such a functor exists, its restriction to one source stage need not land in a bounded target stage: the source stage can have infinitely many objects. The parenthesis on page~13 explicitly adds the necessary boundedness condition. The equivalence asserted there is therefore conditional on this extra factorization, not a formal property of every functor between filtered colimits.

We do not infer a negative theorem from Grothendieck's doubt. The strict composition demanded on his particular category with symmetric end-padding is not constructed in these pages, and we did not locate a resolution of that exact requirement in the localization sources consulted. Other choices of path presentations do support strict composition. Their existence does not by itself provide the requested functor on this fixed category. The distinction is comparable to the one already emphasized in the commentary on Chapter~I, manuscript p.~7 (``Coherent squares and what the axioms remember''): an equality after passage to a homotopy category does not specify its coherent witness.
\end{dcomment}


% ---- from Der-ChV-234_p011-015.tex
\printedpage{9}
% head: A. Grothendieck : Dérivateurs, Ch. V, Édition des Universités de Montpellier II et Paris VII

\msp{page 13}
qui par passage aux ind-objets redonne le morphisme précédent, ou, ce qui revient au même, tel que pour tout $m,n\in\mathbf{N}$ (de même parité), il existe un $N\geq m+n$ assez grand tel que l'on ait un diagramme
\[
\begin{tikzcd}
\Ch_{\Sigma,m}(X;x,y)\times\Ch_{\Sigma,n}(X;y,z) \arrow[r,"*_{m,n}"] \arrow[dd,hook] \arrow[dr] & \Ch_{\Sigma,m+n}(X;x,z) \arrow[d,hook] \\
 & \Ch_{\Sigma,N}(X;x,z) \arrow[d,hook] \\
\Ch_{\Sigma}(X;x,y)\times\Ch_{\Sigma}(X;y,z) \arrow[r,"?"] & \Ch_{\Sigma}(X;x,z)
\end{tikzcd}
\]
et où le triangle supérieur est \textit{commutatif dans} Cat\,hot. (Il faut supposer que pour $N$ assez grand, le composé
\begin{tikzcd}[cramped, sep=small]
{} \arrow[d,hook] & \\
{} \arrow[r,"?"] & {}
\end{tikzcd}
se factorise par $\Ch_{\Sigma,N}(X;x,z)$ \dots)
Pour voir que les foncteurs composition $*_{m,n}$ définissent un homomorphisme des ind-objets dans (Cat\,hot), il suffit de le prouver dans les cas particuliers $(m',n')=(m+1,n+1)$. On va faire la figure en supposant $m$, $n$ \textit{impairs.}

Je m'aperçois ici que je me suis canulé stupidement, la composition de deux chemins de longueur \textit{impaire $m$, $n$ n'est}

\msp{page 14}
\textit{pas un chemin} de longueur $m+n$, i.e. un diagramme de type $\Phi_n$, mais un diagramme de type \textit{dual} $\Phi_n^o=\Psi_n$. Par exemple, pour $m=n=1$
\[
\underbrace{A_0\longleftarrow A_1\longrightarrow A_2}_{c}=\underbrace{B_0\longleftarrow B_1\longrightarrow B_2}_{c'}
\]
est de type $\Phi_2^o$! Donc pour ne pas trop compliquer la vie, on va dans les développements précédents se borner à $m$, $n$ \textit{pairs,} donc des chemins de longueur au moins 4. C'est un peu ridicule! On est ramené à prouver la commutativité à homotopie près dans le cas où on laisse $c'$ fixé et où on envoie $\Ch_{\Sigma,m}(X;x,y)$ dans $\Ch_{\Sigma,m+2}(X;x,y)$, ou inversement $c$ fixe et $\Ch_{\Sigma,n}(X;y,z)$ dans $\Ch_{\Sigma,n+2}$ \texttt{[lire:}$\Ch_{\Sigma,n+2}(X;y,z)$\texttt{]}. Prenons le premier cas, et notons un chemin $c$ sans indices négatifs (ce n'est pas une notation bien heureuse!)
\[
A_0\longrightarrow A_1\longleftarrow A_2\cdots\longleftarrow A_{2m}
\]
(NB à cause de la condition $n$ pair, tous nos chemins commencent désormais par une flèche de $A_0$ dans $A_1$, jamais en sens inverse).

Je vois que l'explicitation de l'homotopie, sans offrir de difficulté, est très fastidieuse à expliquer. Je vais essayer de donner une démonstration conceptuelle, sans `calcul' fastidieux.

\dgrblock{dgrb167}{Even indices and reparametrization relative to the ends}
\begin{dcomment}{exposition}{p14:Je vais essayer de donner une démonstration conceptuelle}{Even indices and reparametrization relative to the ends}
Grothendieck's correction on this page is part of the proof. The earlier statement that equal parity suffices for concatenation is false for two odd indices. Translation by an odd integer reverses the source--sink parity in $\Phi$, so the proposed embeddings of the two odd intervals into $\Phi_{m+n}$ are not functors with the advertised orientation. Restricting to even indices repairs the construction and loses no path, since every finite index is bounded by an even one.

For even $m,n$ use the embeddings
\[
 \Phi_m\longrightarrow\Phi_{m+n},\quad i\longmapsto i-n,
 \qquad
 \Phi_n\longrightarrow\Phi_{m+n},\quad j\longmapsto j+m.
\]
They identify the final vertex of the first interval with the initial vertex of the second. At the join both incident arrows have the same source--sink type, so no new composite of two nonidentity arrows appears. Consequently the resulting pushout in categories is $\Phi_{m+n}$. Restricting a functor on that pushout is the same as giving the two composable diagrams with their common object literally identified. Natural transformations with identity components at the exterior endpoints concatenate as well.

To compare padding before concatenation with padding afterwards, regard the two resulting diagrams as precompositions with maps from a larger interval to a smaller one. Both maps are nondecreasing in the ordinary numerical order of the vertices and fix the two endpoint labels. Their categorical orders are still the alternating orders. This second monotonicity condition, explicitly restored in the edition's footnote on page~17, prevents a backwards source edge from travelling forwards in the smaller interval.

The boundary equation is $h\alpha'=\alpha$, with $\alpha'$ the endpoint map into the larger interval and $\alpha$ the endpoint map into the smaller one. The corrected Hom expression on page~16 fixes the order of those intervals, but its description of the restriction as $\alpha'$ still has the wrong codomain; the displayed equation is the precise interpretation.

The intermediate maps in a homotopy must have that same property. If $c$ is a $\Sigma$-path and $h$ is a numerically nondecreasing functor between alternating intervals, each backwards edge is sent either to a backwards edge or to an identity. Hence $ch$ remains a $\Sigma$-path. Every natural transformation between such $h$'s induces a ladder with identity endpoint components. The forthcoming connectedness argument therefore proves compatibility inside the required path category, with the denominator condition intact.

The reduction to simultaneous increments $(m,n)\mapsto(m+1,n+1)$ before the correction is insufficient for arbitrary independent changes of the two indices. After restriction to even indices, separate increments by two in the first or second variable generate all transitions. Grothendieck's revised formulation on this page makes exactly this repair.
\end{dcomment}


% foot: Version 9 septembre 2001

\printedpage{10}
% head: A. Grothendieck : Dérivateurs, Ch. V, Édition des Universités de Montpellier II et Paris VII

\msp{page 15}
Les deux foncteurs composés
\[
f=k\circ *_{m,n}\ ,\quad g=*_{m',n'}\circ(i\times j)
\]
dans le carré page 11, dont il faut prouver qu'ils sont homotopes dans $\uHom(\mathcal{A},\mathcal{B})$, peuvent être vus comme obtenus de la façon suivante, en termes de deux foncteurs
\[
\varphi,\psi:\Phi_{m'+n'}\rightrightarrows\Phi_{m+n}, \tag*{(1)}
\]
d'où deux foncteurs `transposés'
\[
\begin{tikzcd}
\uHom(\Phi_{m,n},X) \arrow[r,shift left,"F"] \arrow[r,shift right,"G"'] & \uHom(\Phi_{m',n'},X)
\end{tikzcd} \tag*{(2)}
\]
\texttt{[plutôt} $\Phi_{m+n}$\texttt{,} $\Phi_{m'+n'}$\texttt{]} lesquels foncteurs `induisent' $f$ et $g$. De façon précise, désignant par $a$, $b$ les objets source et but de $\Phi_{m+n}$, par $c$ l'objet qui correspond au but du chemin $c$ de longueur $m$ \texttt{[trop de} $c$\texttt{'s]}, à la source du chemin $c'$ de longueur $n$ qu'il s'agit de composer, enfin par $a_0$, $b_0$ les éléments source et but de $\Phi_{m'+n'}$, les foncteurs envisagés (1) appliquent $a_0$ dans $a$, $b_0$ dans $b$, donc définissent un diagramme commutatif
\[
\begin{tikzcd}
\Phi_{m'+n'} \arrow[r,shift left,"\varphi"] \arrow[r,shift right,"\psi"'] & \Phi_{m+n} \\
\underset{\text{\footnotesize catégorie discrète}}{\varepsilon=\{a_0,b_0\}.} \arrow[u,hook,"\alpha'"] \arrow[ur,hook,"\alpha"] &
\end{tikzcd} \tag*{(3)}
\]
Ceci dit, la catégorie but $\mathcal{B}$ pour $f$, $g$, à savoir $\Ch_{\Sigma,m'+n'}(X;x,z)$, s'interprète comme la sous-catégorie pleine de $\uHom(\Phi_{m+n},X)$ formée

\msp{page 16}
des foncteurs $c':\Phi_{m'+n'}\longrightarrow X$ dont le composé avec $\alpha'$ est le foncteur
\[
\beta:\varepsilon\longrightarrow X\qquad\qquad
\begin{aligned}
a_0&\longrightarrow x\\
b_0&\longrightarrow z,
\end{aligned} \tag*{$(*)$}
\]
et de même la catégorie source $\mathcal{A}$ s'interprète comme la sous-catégorie pleine de $\uHom(\Phi_{m+n},X)$ formée des $c:\Phi_{m+n}\longrightarrow X$ dont la restriction à $\{a,b,c\}$ est le foncteur
\[
a\longrightarrow x,\ b\longrightarrow y,\ c\longrightarrow z.
\]
\texttt{[plutôt} $b\longrightarrow z$\texttt{,} $c\longrightarrow y$\texttt{]} Elle est contenue dans la sous-catégorie pleine $\mathcal{A}'$ des foncteurs dont la restriction à $\{a,b\}$ soit $a\longrightarrow x,\ b\longrightarrow z$, i.e. dont le composé avec $\alpha$ soit le foncteur $\beta:\varepsilon\longrightarrow X$ déjà envisagé ci-dessus. Ainsi les foncteurs $f,g:\mathcal{A}\rightrightarrows\mathcal{B}$ se factorisent par $\mathcal{A}'$
% foot: Version 9 septembre 2001
\printedpage{11}
% head: A. Grothendieck : Dérivateurs, Ch. V, Édition des Universités de Montpellier II et Paris VII
\[
\begin{tikzcd}
\mathcal{A} \arrow[r,hook] \arrow[dr,shift left,"f"] \arrow[dr,shift right,"g"'] & \mathcal{A}' \arrow[d,shift right,"g'"'] \arrow[d,shift left,"f'"] \\
 & \mathcal{B},
\end{tikzcd}
\]
foncteurs déduits de $\varphi$, $\psi$ ci-dessus, donnant lieu aux triangles commutatifs (3), et du foncteur $\beta:\varepsilon\longrightarrow X$, de la façon dite. Pour prouver que $f$ et $g$ sont homotopes, il suffit de le prouver pour $f'$, $g'$. Et pour ceci, il suffit de montrer que $\varphi$ et $\psi$ sont homotopes en tant qu'objets de la catégorie $\uHom(\Phi_{m+n},\Phi_{m'+n'};\alpha)$ formée des foncteurs dont la restriction à $\varepsilon$ est $\alpha$ \texttt{[plutôt de la catégorie} $\uHom(\Phi_{m'+n'},\Phi_{m+n};\alpha')$ \texttt{formée des foncteurs dont la restriction à} $\varepsilon$ \texttt{est} $\alpha'$ \texttt{]}.

\msp{page 17}
D'ailleurs, les deux assertions (sur $f'$, $g'$, quels que soient $X$, $x$, $y$, $z$, \texttt{[plutôt} $x$\texttt{,} $z$\texttt{]} et sur $\varphi$, $\psi$, $\alpha$ \texttt{[plutôt} $\alpha'$\texttt{]}) sont en fait équivalentes comme on voit en faisant $X=\Phi_{m+n}$, $x=a$, $y=c$, $z=b$ \texttt{[plutôt} $x=a$\texttt{,} $z=b$\texttt{]}. Ceci sera contenu dans le lemme plus général. (\footnote[4]{Dans les considérations qui suivent, on considère en fait tacitement que les foncteurs envisagés sont \textit{croissants} sur les sommets pour l'ordre de leur numérotation. Donc il faudrait écrire $\uHom_{cr}$ au lieu de $\uHom$, dans les énoncés qui suivent, y compris dans le lemme de connexité ci-contre.})

\textbf{Lemme.} \textit{Soient $k$ et $l$ deux entiers. Alors la sous-catégorie pleine de $\uHom(\Phi_k,\Phi_l)$ formée des foncteurs $\Phi_k\longrightarrow\Phi_l$ appliquant source en source, but en but, est connexe.} (NB Elle est non vide si et seulement si $k\geq l$.)

On va décrire entièrement ces foncteurs (= applications croissantes) $\Phi_k\xrightarrow{\varphi}\Phi_l$, en identifiant $\Fl(\Phi_k)$ \texttt{[plutôt l'ensemble des flèches non identiques de} $\Phi_k$\texttt{]} à l'intervalle d'entiers $[1,2k]$, et en associant à $\varphi$ la partie
\[
E_\varphi\subset\Fl(\Phi_k)=[1,2k]\qquad\qquad E_\varphi=\{u\in\Fl(\Phi_k)\mid\varphi(u)\text{ une identité}\}
\]
formée des $u\in\Fl(\Phi_k)$ qui sont `contractées' par $\varphi$ en une identité. Si $E$ est un ensemble d'entiers, on y introduit la notion de composantes connexes (par la relation d'équivalence dans $E$ engendrée par la relation de succession $y=x+1$). On voit aussitôt que les composantes connnexes de $E_\varphi$ qui ne contiennent ni la première ni la dernière flèche de $\Phi_k$ sont paires, et que % [sic] connnexes -> connexes, triple n
\[
\underbrace{\card\Fl(\Phi_k)}_{2k}=\card\,(E_\varphi)+\underbrace{\card\Fl(\Phi_l)}_{2l},
\]

\dgrblock{dgrb168}{A contraction is determined by the edges it erases}
\begin{dcomment}{exposition}{p17:On va décrire entièrement ces foncteurs}{A contraction is determined by the edges it erases}
Grothendieck reduces a diagrammatic homotopy problem to finite combinatorics. Number the vertices of $\Phi_k$ from $0$ to $2k$, preserving their numerical order; the vertex numbered $t$ is the original vertex $t-k$. Number the edges from $1$ to $2k$. For a numerically nondecreasing endpoint-preserving functor $f:\Phi_k\to\Phi_l$, adjacent values differ by zero or one: the target has no nonidentity arrow joining nonadjacent vertices. Every noncollapsed edge advances the numerical position once. Therefore
\[
 f_E(t)=-l+t-\bigl|E\cap[1,t]\bigr|,
 \qquad 0\leq t\leq2k,
\]
where the values are written in the original target numbering and $E$ is the set of collapsed edges. This formula both proves uniqueness and makes the prefix and suffix observations on page~20 transparent.

When $l>0$, the admissibility conditions have a direct explanation. Exactly $2l$ edges survive, so $|E|=2(k-l)$. Between successive surviving edges, an even number of erased edges is necessary and sufficient for their alternating orientations to agree. The same condition holds at an endpoint when the initial orientations agree. If $k$ and $l$ have opposite parity, the initial and terminal erased runs must both be odd. The parity of the terminal run follows from the total cardinality together with the parity of all interior runs; the edition rightly inserts both conditions (a) and (b) at that point.

Conversely, these conditions make each consecutive pair $f_E(t-1),f_E(t)$ either equal or the endpoints of the correctly oriented target edge. They thus define a functor. It is surjective on objects and nonidentity arrows, because it starts at $-l$, ends at $l$, and advances by steps of at most one. Its only identifications are the consecutive strings of vertices joined by edges of $E$, so it is the stated quotient. This establishes necessity and sufficiency, not merely the numerical count.

For $l=0$ there is one constant functor and $E=[1,2k]$. No initial target direction exists, so the parity condition must be omitted in that case. For $k=l=0$ the same convention gives $E=\varnothing$. Endpoint-preserving functors exist exactly when $k\geq l$: necessity follows from counting surviving edges, while sufficiency follows from the clamping functor $\Phi_k\to\Phi_l$.

The footnote's restriction to numerical monotonicity is essential for the description by $E$. A general functor can retrace a target edge, and its noncollapsed edges need not number $2l$. Later we prove contractibility also for the full endpoint fibre, but that stronger statement does not make the intervening classification of contractions valid without its hypothesis.

For instance, number the source vertices of $\Phi_3$ from $0$ to $6$, and retain the target numbering $-1,0,1$ for $\Phi_1$. The row
\[
 (-1,0,-1,0,-1,0,1)
\]
defines an endpoint-preserving functor. At the odd-numbered source vertices, which are sources in this numbering, its value is $0$; every generating arrow is therefore sent to one of the two nonidentity arrows leaving $0$ in $\Phi_1$. No edge is erased. Nevertheless the difference in the indices is $3-1=2$, so the proposed cardinality would be four. This example isolates the precise role of the footnote: endpoint preservation is sufficient for the later contractibility theorem, but it is insufficient for this particular description by a set of contractions.
\end{dcomment}


\msp{page 18}
i.e.
\[
\card\,(E_\varphi)=2(k-l).
\]

% foot: Version 9 septembre 2001

\printedpage{12}
% head: A. Grothendieck : Dérivateurs, Ch. V, Édition des Universités de Montpellier II et Paris VII

Il est immédiat qu'un foncteur $\varphi$ est surjectif sur les objets et sur les flèches, donc identifie $\Phi_l$ à une catégorie quotient de $\Phi_k$ - et la relation d'équivalence dans $\Phi_k$ définie par $\varphi$ est celle obtenue en identifiant tous les objets qui sont source ou but de flèches d'une même composante connexe $I$ de $E_\varphi$
\[
\cdots\longleftarrow\ \underbrace{\longrightarrow\ \longleftarrow\ \longrightarrow\ \longleftarrow}_{I}\ \longrightarrow\cdots
\]
et les flèches dans cette composante à l'identité correspondante. Les conditions énoncées sur $E$ (indépendamment d'un $\varphi$ qui lui donnerait naissance assurent que la catégorie quotient est isomorphe (avec \textit{préservation} de source et de but) soit à $\Phi_l$, soit à la duale $\Psi_l=\Phi_l^o$, suivant la direction de la première flèche de $\bar{\Phi}_k$ (quotient de $\Phi_k$). Cette `direction initiale' est la même que celle pour $\Phi_k$ si la première flèche de $\Phi_k$ n'appartient pas à $E$, ou bien si elle y appartient et sa composante connexe dans $E$ est de cardinal pair. Dans le cas inverse ($u\in E$ et sa composante connnexe dans $E$ est de cardinal impair), la direction initiale de $\bar{\Phi}_k$ est opposée à celle de $\Phi_k$. % [sic] "(indépendamment d'un $\varphi$ qui lui donnerait naissance" -> parenthèse fermante manquante, the opening parenthesis is never closed; [sic] connnexe -> connexe, triple n

\msp{page 19}
On voit donc que $E$ doit satisfaire encore à une condition supplémentaire, pour que le $\bar{\Phi}_k$ correspondant soit isomorphe à $\Phi_l$, i.e. pour que $E$ provienne d'un foncteur $\varphi$. Je résume toutes les conditions

\begin{itemize}
\item[a)] $\card E=2(k-l)$.
\item[b)] Les composantes connexes de $E$ qui ne contiennent pas de flèche initiale ou finale sont de cardinal pair.
\item[c)] Si $\Phi_k$ et $\Phi_l$ ont même direction initiale ($k$ et $l$ de même parité), alors l'éventuelle composante de $E$ contenant la flèche initiale doit être elle aussi de cardinal pair (ce qui, à cause de a) \texttt{[plutôt a) et b)]}, équivaut à la même condition pour l'éventuelle composante connexe contenant la flèche finale de $\Phi_k$). Si $\Phi_k$ et $\Phi_l$ ont des directions initiales différentes, il faut que la flèche initiale de $\Phi_k$ appartienne à une composante connexe de cardinal \textit{impair} de $E$ (ce qui équivaut à itou pour la flèche finale).
\end{itemize}

D'autre part, le foncteur $\varphi$ est uniquement déterminé par la connaissance de $E$, car il y a un \textit{unique} isomorphisme $\bar{\Phi}_k\xrightarrow{\sim}\Phi_l$ compatible avec sources et buts, à cause du

\textbf{Lemme} (immédiat) : \textit{Un automorphisme de $\Phi_l$ qui est l'identité sur l'objet source (ou sur l'objet but) est l'identité partout.}

Les foncteurs $\varphi$, $\psi$ de tantôt correspondaient au choix, l'un de $E_\varphi=\{1,2,k-1,k\}$ \texttt{[plutôt} $E_\varphi=\{1,2,2k-1,2k\}$\texttt{]}, l'autre de $E_\psi=\{1,2,i,i+1\}$, 2 composantes connexes

\msp{page 20}
dans un cas, trois \texttt{[plutôt deux]} dans l'autre.) % [sic] "l'autre.)" -> parenthèse fermante sans parenthèse ouvrante, no matching opening parenthesis

% foot: Version 9 septembre 2001

\printedpage{13}
% head: A. Grothendieck : Dérivateurs, Ch. V, Édition des Universités de Montpellier II et Paris VII

Dans le cas général, il y a lieu d'expliciter la relation d'ordre de l'ensemble $\uHom^{!}(\Phi_k,\Phi_l)$ (! signifie : compatible avec source et but), en termes des ensembles représentatifs $E_\varphi$, $E_\psi$. Je ne vais pas expliciter cette relation d'ordre, mais seulement une propriété de celle-ci qui nous sera suffisante. Soit $\varphi_E$ le foncteur défini par une partie $E\subset[1,2k]$. Tout d'abord :

\textbf{Lemme 1} (immédiat) : \textit{Soient $E,E'\subset[1,2k]$ admissibles, et soit $j\in[1,2k]$ tel que $E\cap[1,j]=E'\cap[1,j]$. Alors $\varphi_E$ et $\varphi_{E'}$ coïncident sur le but de $u_j$ ($j$ une flèche de $\Phi_k$) et sur tous les objets de $\Phi_k$ antérieurs à celle-ci (au sens de la numérotation des objets, et non de la relation d'ordre de la catégorie $\Phi_k$ bien sûr.) Dualement, si $E\cap[j,2k]=E'\cap[j,2k]$, alors $\varphi_E$ et $\varphi_{E'}$ coïncident sur la source de $y$, et sur tous les objets de $\Phi_k$ postérieurs à celui-ci.} % [sic] la source de $y$ -> la source de $u_j$, apparently intended by duality with "le but de $u_j$"

Considérons maintenant une partie connexe $I$ de $E$ (pas forcément une composante connexe), $I=[r,s]\subset[1,2k]$ ($r$ le plus petit élément de $I$, $s$ son plus grand élément) et supposons que

\begin{itemize}
\item[a)] $\card I$ est pair (i.e. $s-r$ est pair) \texttt{[plutôt} $s-r$ \texttt{impair]}
\item[b)] $r\neq1$, et $r-1\notin E$ (i.e. $r-1\in[1,2k]\backslash E$)
\end{itemize}

et soit
\[
E'=(E\backslash\{s\})\cup\{r-1\}=E\backslash I\cup\overbrace{[r-1,s-1]}^{I'}
\]
déduit de $I$ en `translatant' sur la gauche d'un cran.

\msp{page 21}
Il est immédiat que si $E$ est admissible, $E'$ aussi (et inversement). Ceci dit, on a le

\textbf{Lemme 2.} \textit{$E$ et $E'$ étant comme ci-dessus, il existe une homomorphisme de foncteurs entre $\varphi_E$ et $\varphi_{E'}$} (je ne dis pas dans quel sens va la flèche!). % [sic] une homomorphisme -> un homomorphisme, wrong gender

\textbf{Corollaire.} Énoncé dual, quand $E'$ est déduit de $E$ en translatant à droite d'un cran un intervalle connexe de longueur paire $I=[r,s]$ de $E$, $s\neq2k$ et $s+1\notin E$.

En effet, dans ce cas $E$ est déduit de $E'$ par une opération du type explicité dans le lemme 2.

La connexité de la catégorie $\uHom^{!}(\Phi_k,\Phi_l)$ est maintenant claire. Si $k$ et $l$ ont même parité, en poussant les composantes connexes de $E$ vers la gauche, on écrit $E$ sous la forme `normale'
\[
E_0=[1,2(k-l)].
\]
Si $k$ et $l$ sont de parité inverse, on ne peut pousser toute la composante connexe de $E$ qui contient l'élément dernier $2k$, car elle est de cardinal impair. Il faut laisser $2k$ dans $E$ et pousser le reste, et on ramène $E$ à la forme normale
\[
E_0=[1,2(k-l)-1]\cup\{2k\}.
\]

\dgrblock{dgrb169}{Sliding an even block gives a natural transformation}
\begin{dcomment}{exposition}{p21:La connexité de la catégorie}{Sliding an even block gives a natural transformation}
Grothendieck's elementary move explains why different padding conventions become equivalent. The formula just obtained gives a short proof of the move, including the natural transformation he leaves implicit. Suppose $I=[r,s]\subset E$ has even cardinality, $r>1$, and $r-1\notin E$, and put $E'=(E\setminus\{s\})\cup\{r-1\}$. Then
\[
 |E'\cap[1,t]|-|E\cap[1,t]|=
 \begin{cases}1,&r-1\leq t\leq s-1,\\0,&\text{otherwise}.
 \end{cases}
\]
The map $f_E$ is constant, say equal to $a$, on the vertices $r-1,\ldots,s$, because all intervening edges belong to $E$. Thus $f_{E'}$ differs from $f_E$ only by replacing $a$ with $a-1$ at the vertices $r-1,\ldots,s-1$. These two target vertices are joined by one arrow. Its orientation determines a natural transformation in one direction between the two functors; outside the displayed interval its components are identities. Naturality is automatic in a poset once the component inequalities hold.

The endpoints are unchanged, since $r-1>0$ and $s-1<2k$. Admissibility is preserved: removing an even initial block from a collapsed run preserves the parity of its remainder, and translating that even block preserves the parity of every run it meets. The initial and terminal parity conditions are therefore unchanged. This also proves the right-moving version by reversing the move.

Repeated left-moving operations terminate: the sum of the indices in $E$ strictly decreases at each move. If the endpoint orientations agree, every collapsed run is even and all runs can be joined to the leftmost one, yielding $[1,2(k-l)]$. If they disagree, the initial and terminal runs are odd. Keep the final edge $2k$ in place and move the remaining even part of the final run, together with the interior runs, to the left. The normal form is
\[
 [1,2(k-l)-1]\cup\{2k\}.
\]
Each move is a morphism in one of the two directions, so every object is connected to the same normal form. This is a proof of connectedness, which is all the composition argument needs.

For example, with $k=4,l=2$, the set $E=\{1,2,7,8\}$ successively becomes $\{1,2,6,7\}$, $\{1,2,5,6\}$, $\{1,2,4,5\}$, and $\{1,2,3,4\}$. The associated functors form a zigzag of four natural transformations relative to the endpoints. This gives a visible instance of the homotopy comparing padding at an exterior end with padding near a join. The proof does not assert that there is always one natural transformation between the original two functors.
\end{dcomment}


\hfill cqfd!

% foot: Version 9 septembre 2001

% ---- from Der-ChV-234_p016-020.tex
\printedpage{14}
% head: A. Grothendieck : Dérivateurs, Ch. V, Édition des Universités de Montpellier II et Paris VII
Ce qui précède justifie notre raisonnement, qui donnait par composition des chemins des applications de composition bien définies
\[
\underbrace{\Hom_\Sigma(x,y)}_{=\pi_0\uCh_\Sigma(x,y)} \times \underbrace{\Hom_\Sigma(y,z)}_{=\pi_0\uCh_\Sigma(y,z)} \longrightarrow \underbrace{\Hom_\Sigma(x,z)}_{=\pi_0\underline{\textstyle\mathrm{Ch}}_\Sigma(x,z)}.
\]
% (the third "Ch" under the brace is printed in a larger size than the two others)

\msp{page 22}Je dis que cette composition est associative. Ceci résulte du fait que si on a trois chemins (d'indices \textit{pairs}!)
\[
x \xrightarrow{c} y \xrightarrow{c'} z \xrightarrow{c''} t,
\]
on a (sans avoir à effectuer des homotopies \dots) associativité stricte
\[
c''(c'c) = (c''c')c.
\]

Je dis qu'il y a des identités. On définit
\[
1_x \in \Hom_\Sigma(x,x)
\]
par le chemin d'indice 0 (qui est pair!), i.e. par le diagramme réduit à $x$, sans flèche (ou si on veut interpréter le diagramme comme un foncteur $e \longrightarrow X$, il y a la flèche identique sous-entendue). Nos raisonnements n'excluaient pas le cas $n = 0$ ou $m = 0$, et il est tautologique que pour un chemin
\[
x \xrightarrow{c} y
\]
on a
\[
c \cdot 1_x = 1_y \cdot c = c.
\]
Donc on a bien défini une catégorie, que je vais appeler $X_\Sigma$. Ja vais définir un foncteur canonique % [sic] Ja -> Je, misprint of the pronoun
\[
\theta : X \longrightarrow X_\Sigma
\]
qui sera l'identité sur les objets, il faut le définir sur les flèches
\[
\Hom_X(x,y) \longrightarrow \Hom_{X_\Sigma}(x,y) = \pi_0\mathrm{Ch}_\Sigma(x,y)
\]
\texttt{[plutôt} $\pi_0\uCh_\Sigma(x,y)$\texttt{]}. Malheureusement, une flèche
\[
x \xrightarrow{u} y
\]
\msp{page 23}ne peut s'interpréter directement comme un chemin à notre sens, il faut payer nos fuites! Mais on va à une telle flèche associer le chemin
\[
x \longrightarrow y \xleftarrow{\mathrm{id}} y \xrightarrow{\mathrm{id}} y \xleftarrow{\mathrm{id}} y,
\]
% foot: Version 9 septembre 2001
\printedpage{15}
% head: A. Grothendieck : Dérivateurs, Ch. V, Édition des Universités de Montpellier II et Paris VII
élément de $\uCh_{\Sigma,2}(x,y)$, ce qui définit l'application cherchée. Il faut vérifier la compatibilité avec les identités et sous la composition. À la flèche $1_x : x \longrightarrow x$ est associé l'élément de $\Phi_2(x,x)$
\[
x \xrightarrow{\mathrm{id}} x \xleftarrow{\mathrm{id}} x \xrightarrow{\mathrm{id}} x \xleftarrow{\mathrm{id}} x,
\]
il faut prouver que dans $\pi_0(\uCh_{\Sigma,2}(x,x))$ il devient égal à l'élément déduit de $1_x \in \Phi_{\Sigma,0}(x,x)$, par le foncteur de transition. Mais cet élément n'est autre justement que
\[
x \xrightarrow{\mathrm{id}} x \xleftarrow{\mathrm{id}} x \xrightarrow{\mathrm{id}} x \xleftarrow{\mathrm{id}} x,
\]
c'est le même (sans avoir même à passer au $\pi_0$)!

\dgrblock{dgrb170}{A denominator set need not be closed under composition}
\begin{dcomment}{examples}{p22:Je dis qu'il y a des identités.}{A denominator set need not be closed under composition}
Grothendieck's use of arbitrarily long paths is already visible in the simplest one-object example. Let $X$ have one object and endomorphism monoid $(\mathbf N,+)$, writing its generator multiplicatively as $t$. Take
\[
 \Sigma=\{1,t\}.
\]
This contains the identity but is not closed under composition: $t^2\notin\Sigma$. A backwards edge may carry $t$, but not $t^2$. For a path $c$, let $a(c)$ be the sum of the exponents on its forward edges and let $b(c)$ count the backwards edges labelled $t$. Its evaluation is determined by the integer
\[
 \nu(c)=a(c)-b(c).
\]
A commutative ladder with identity endpoints preserves this integer. Indeed, write each vertical component as $t^{h_i}$ with $h_0=h_{2n}=0$. For a forward edge of exponent $a_i$, naturality gives $h_i+a_i=a_i'+h_{i-1}$; for a backwards edge the analogous equality has the horizontal exponents reversed. Summing the signed equalities cancels all the $h_i$, leaving $\nu(c)=\nu(c')$.

Conversely, every integer is represented. A nonnegative integer $r$ is represented by the ordinary arrow $t^r$ with identity padding. For $-r<0$, concatenate $r$ backwards copies of $t$, separating them by identities, then pad to an even index. The cancellation ladders below show that a forward $t$ followed by a backwards $t$ represents the identity, in either order. Splitting a forward $t^a$ into $a$ successive copies of $t$ is allowed by the functoriality calculation for ordinary arrows. These two operations reduce any path to the representative of its signed exponent. Thus
\[
 \pi_0\operatorname{Ch}_\Sigma(X;*,*)\cong\mathbf Z,
\]
with concatenation corresponding to addition.

The resulting localization is the one-object category of the group $\mathbf Z$. Its universal property says exactly that a chosen endomorphism $T$ becomes invertible: the image of the integer $r$ must be $T^r$. Negative integers require arbitrarily many backwards edges when the denominator set is the original $\{1,t\}$. Replacing it by all powers of $t$ permits a single backwards edge to represent them, but this changes the categories of presentations. The agreement after passage to connected components is the invariant conclusion.
\end{dcomment}


Pour la compatibilité avec la composition
\[
x \xrightarrow{u} y \xrightarrow{v} z,
\]
on a d'une part
\[
x \xrightarrow{vu} z \xleftarrow{\mathrm{id}} z \xrightarrow{\mathrm{id}} z \xleftarrow{\mathrm{id}} z,
\]
d'autre part
\[
x \xrightarrow{u} y \xleftarrow{\mathrm{id}} y \xrightarrow{\mathrm{id}} y \xleftarrow{\mathrm{id}} y \xrightarrow{v} z \xleftarrow{\mathrm{id}} z \xrightarrow{\mathrm{id}} z \xleftarrow{\mathrm{id}} z.
\]
On complète le premier chemin par dégénérescence aux deux bouts, pour le comparer au deuxième en tant qu'élément de $\uCh_4(x,z)$,
\[
x \xrightarrow{\mathrm{id}} x \xleftarrow{\mathrm{id}} x \xrightarrow{vu} z \xleftarrow{\mathrm{id}} z \xrightarrow{\mathrm{id}} z \xleftarrow{\mathrm{id}} z \xrightarrow{\mathrm{id}} z \xleftarrow{\mathrm{id}} z.
\]

\dgrblock{dgrb171}{Two ordinary arrows and one comparison diagram}
\begin{dcomment}{examples}{p23:Pour la compatibilité avec la composition}{Two ordinary arrows and one comparison diagram}
Grothendieck needs to show that sending an ordinary arrow to its padded path respects composition. This particular comparison is smaller than the general interval problem and can be completed explicitly. Let $x\xrightarrow{u}y\xrightarrow{v}z$, and regard it as a functor $T:[2]\to X$ with $T(0)=x$, $T(1)=y$, and $T(2)=z$.

Number the nine vertices of $\Phi_4$ from $0$ to $8$. The concatenated path for $u$ and $v$, and the padded path for $vu$, arise from the following functors to the ordered category $[2]$:
\[
\begin{array}{c|rrrrrrrrr}
 i&0&1&2&3&4&5&6&7&8\\ \hline
 f(i)&0&1&1&1&1&2&2&2&2\\
 g(i)&0&0&0&2&2&2&2&2&2\\
 h(i)&0&0&0&1&1&2&2&2&2.
\end{array}
\]
Here $h(i)=\min(f(i),g(i))$. Every row is a functor: along each forward edge its values are nondecreasing, while the values on the two ends of every backwards edge are equal. Pointwise inequalities give natural transformations
\[
 f\longleftarrow h\longrightarrow g.
\]
Their endpoint components are identities, because all three rows take the values $0$ and $2$ at the two ends. Applying $T$ produces the desired two-step ladder in $\operatorname{Ch}_{\Sigma,4}(X;x,z)$. No nonidentity arrow is placed on a backwards edge, so the proof works for any $\Sigma$ containing the identities.

This calculation establishes $\theta(v)\theta(u)=\theta(vu)$ after passage to components. For identities the two diagrams already agree at the finite stage, as Grothendieck observes. Hence $\theta$ is a functor without assuming the general lemma. The display $\Delta_1\sqcup\Delta_1\simeq\Delta_2$ on the following manuscript page must mean amalgamation over the common object, $[1]\sqcup_{[0]}[1]\cong[2]$; the ordinary disjoint union has four objects.

The table also records why an unqualified homotopy in a functor category would be insufficient. An arbitrary map $\Phi_4\to[2]$ may send a backwards edge to a nonidentity arrow. Its image under $T$ would then require $u$, $v$, or $vu$ to lie in $\Sigma$, which has not been assumed. Our three rows avoid that problem by making every backwards edge constant.
\end{dcomment}


\msp{page 24}Il faut prouver que les deux sont homotopes. On a donc deux foncteurs
\[
f,g : \underbrace{\overbrace{\uFl(X;x,y)}^{\text{catégorie discrète}} \times \overbrace{\uFl(X;y,z)}^{\text{catégorie discrète}}}_{\mathcal{A}} \underset{g}{\overset{f}{\rightrightarrows}} \underbrace{\uCh_{4,\Sigma}(X;x,z)}_{\mathcal{B}}, \tag*{$(\ast)$}
\]
on veut prouver qu'ils sont homotopes. Pour ceci, on les regarde (comme tantôt) comme `transposés' des deux foncteurs
\[
\begin{tikzcd}[column sep=large]
\bDelta_1 \sqcup \bDelta_1 \simeq \bDelta_2 & \Phi_4 \arrow[l, shift right, "\varphi"'] \arrow[l, shift left, "\psi"] \\
& \varepsilon, \arrow[ul, "\alpha"] \arrow[u, "\alpha'"']
\end{tikzcd}
\]
vu que les catégories $\mathcal{A}$, $\mathcal{B}$ des deux membres de (*) s'interprètent comme des sous-catégories pleines de $\uHom(\bDelta_2, X)$ et $\uHom(\Phi_4, X)$ respectivement; à savoir les sous-catégories pleines formées des foncteurs dont le composé avec $\alpha$ resp. $\alpha'$ est donné par $x$ et par $z$, et de plus (pour $\mathcal{A}$) dont la valeur en l'objet médian de $\bDelta_2$ soit $y$. On peut encore remplacer $\mathcal{A}$ par la sous-catégorie pleine plus grande $\mathcal{A}'$ de $\uHom(\bDelta_2, X)$, obtenue en ignorant cette dernière condition, et il suffit de prouver que les deux foncteurs `plus gros'
% foot: Version 9 septembre 2001
\printedpage{16}
% head: A. Grothendieck : Dérivateurs, Ch. V, Édition des Universités de Montpellier II et Paris VII
\[
\mathcal{A}' \underset{g'}{\overset{f'}{\rightrightarrows}} \mathcal{B}'
\]
sont homotopes. Et pour ceci, il suffit de
\msp{page 25}voir encore que $\varphi$ et $\psi$ sont homotopes dans la catégories des foncteurs $\Phi_4 \longrightarrow \bDelta_2$ qui sont compatibles avec la source, et le but. On voit donc que notre lemme d'homotopie de tantôt n'était pas assez général. Il faudrait définir les \textit{catégories intervalles généralisés,} de longueur $n$ quelconque ($n =$ nombre de flèches bout à bout), une catégorie intervalle généralisé \textit{standard} de longueur $n$ a pour ensemble d'objets l'intervalle d'entiers % [sic] la catégories -> la catégorie, singular article with plural noun
\[
[0,n],
\]
et pour tout couple d'entiers consécutifs $i-1$, $i$ parmi les objets ($1 \leq i \leq n$), il y a exactement une flèche qui relie $i$, $i+1$ \texttt{[plutôt} $i-1$\texttt{,} $i$\texttt{]}. On affecte l'intervalle $[i,i+1]$ du signe $+$ si $i \longrightarrow i+1$ (flèche en sens direct), du signe $i$ \texttt{[plutôt} $-$\texttt{]} si $i \longleftarrow i+1$ (flèche en sens rétrograde) - donc les catégories standards sont définies, pour $n$ fixé, par un système de signes $+$ ou $-$, dont on affecte les $n$ `tronçons' consécutifs $A_i := [i-1,i]$ ($1 \leq i \leq n$) de l'intervalle d'entiers $[0,n]$. Il y en a donc exactement $2^n$.

On s'attend à ce que l'énoncé suivant soit vrai :

\dgrblock{dgrb172}{Contracting relative functor categories into an interval}
\begin{dcomment}{development}{p25:On s'attend à ce que l'énoncé suivant soit vrai}{Contracting relative functor categories into an interval}
Grothendieck's question asks whether the choices used to change a path's parametrization form a contractible category. We can prove a slightly more flexible statement that also accommodates the cancellation diagrams on page~27, where both endpoints have the same image. The proof uses the numerical order of an interval only to define minimum and maximum; the categorical order retains the prescribed directions of all its edges.

\textbf{Relative interval contraction.}\ours\ Let $J$ be the free category on an oriented path with vertices $0,\ldots,r$. Let $j:B\to A$ be a functor and $b:B\to J$ a fixed diagram. The strict fibre
\[
 \mathcal F_b=
 \{f:A\to J\mid fj=b\}
\]
of restriction, with transformations restricting to $1_b$, is either empty or categorically contractible: if $g\in\mathcal F_b$, its identity endofunctor is joined to the constant endofunctor at $g$ by a finite zigzag of natural transformations. If $A$ is itself an oriented interval and the fibre is restricted to functors nondecreasing in the numerical vertex orders, the same assertion holds whenever that restricted fibre is nonempty.

\emph{Proof.} The category $J$ is a poset. Its order relation $a\leq_J a'$ means that the unique numerical segment between $a$ and $a'$ is oriented all the way from $a$ to $a'$. For fixed $c$, each of the numerical clamping maps
\[
 a\longmapsto\min(a,c),\qquad
 a\longmapsto\max(a,c)
\]
is a functor $J\to J$: the image of a directed segment is either a point or a directed subsegment with the same orientation. Consequently numerical minimum and maximum are functors $J\times J\to J$. Indeed they are order preserving in either variable separately, and a product-order comparison factors into one comparison in each variable.

Assume that $g\in\mathcal F_b$. For each vertex $t$ of $J$ define endofunctors of $\mathcal F_b$ by pointwise operations
\[
 F_t(f)=\min\bigl(f,\max(g,t)\bigr),\qquad
 G_t(f)=\min\bigl(g,\max(f,t)\bigr),
\]
where $t$ denotes the constant functor $A\to J$. These are functors on diagrams and on natural transformations because minimum and maximum are functors. They preserve the prescribed restriction: at an object of $B$, both $f$ and $g$ have the same value $a$, and
\[
 \min(a,\max(a,t))=a.
\]
Since $J$ is a poset, the component of any resulting transformation at such an object is the identity. Thus the construction really stays inside the strict fibre.

Each generating arrow $t\to t'$ of $J$ induces natural transformations $F_t\to F_{t'}$ and $G_t\to G_{t'}$. If an edge points in the opposite direction, both induced transformations point in that opposite direction. We therefore obtain finite zigzags as $t$ runs numerically from $0$ to $r$. At the ends,
\[
 F_r=1_{\mathcal F_b},\qquad
 F_0=\min(-,g)=G_0,\qquad
 G_r=\operatorname{const}_g.
\]
Traversing the $F$ zigzag backwards and the $G$ zigzag forwards contracts the identity to $g$. Every functor in the zigzag fixes $g$ itself. On taking nerves and geometric realizations, natural transformations give homotopies, so $|N\mathcal F_b|$ is contractible as well. If $r=0$, the nonempty fibre has one object and the conclusion is immediate.

For the numerical-monotonicity variant, the pointwise minimum or maximum of two nondecreasing numerical functions is nondecreasing, and constant functions are nondecreasing. Hence all $F_t$ and $G_t$ preserve the indicated full subcategory. Their restrictions give the same contraction. This completes the proof.\hfill$\square$

The construction gives a uniform bound on the length of the contracting zigzag: there are $r$ comparisons in the $F$ family and $r$ in the $G$ family. This bound depends only on the target interval, not on the number of objects or arrows of $A$. In particular, no finite-source hypothesis was concealed in the proof. When $r=1$, the contraction has the particularly transparent form
\[
 1_{\mathcal F_b}\longleftarrow\min(-,g)
       \longrightarrow\operatorname{const}_g
\]
if the target edge points from $0$ to $1$; both arrows reverse if it points from $1$ to $0$. The fixed-boundary identity used above makes this a contraction inside the fibre even when $j$ identifies several boundary labels. That includes the two coincident endpoint labels of an excursion.

Taking $B$ to be the two endpoint labels proves both the proposed connectedness and the stronger contractibility of the generalized-interval functor category. Taking a fixed pair of equal endpoint values gives the version needed for an excursion followed by its return. For alternating intervals, the numerically monotone version strengthens the contraction-set argument on pages~17--21 and preserves the condition on backwards arrows when used for reparametrization.

The minimum--maximum preservation property is familiar in the theory of ordered digraphs; compare the definition of a Min--Max ordering in \citet[p.~1577]{HellRafiey2012}. We checked that account and the localization sources cited in this section, together with searches for relative functor categories of oriented intervals. We did not locate this fixed-restriction contraction in those sources. The result is presented as a direct continuation of the manuscript, with its full proof, and carries no claim of priority.
\end{dcomment}


\textbf{Lemme.} \textit{Soient $I$, $I'$ deux intervalles généralisés, munis d'objets source et objet but (de sorte que $I$, $I'$ soient} canoniquement \textit{isomorphes aux intervalles genéralisés standards correspondants). Soit} $\uHom^!(I,I')$ % [sic] genéralisés -> généralisés, missing accent on the first e
\msp{page 26}\textit{la catégorie des foncteurs de $I$ dans $I'$, compatibles avec source et but. Si cette catégorie est $\neq \emptyset$, elle est} 0\textit{-connexe} (et, je présume, même contractile?).

Je vais admettre provisoirement qu'il en est ainsi, du moins dans le cas particulier de $\Phi_4$ (intervalle généralisé de longueur 8) et $\bDelta_2$. Cela donne donc un foncteur
\[
\theta : X \longrightarrow X_\Sigma.
\]
On a gagné, car la propriété universelle de ce foncteur est pratiquement évidente. Mais j'ai oublié de vérifier que $\theta$ transforme éléments de $\Sigma$ en flèches inversibles. Soit donc
\[
x \xrightarrow{u} y
\]
avec $u \in \Sigma$, on va définir un chemin en sens inverse
\[
y \xrightarrow{c} x
\]
qui définira dans $X_\Sigma$ une flèche inverse de $u$. L'idée naturelle est de prendre le `chemin'
\[
y \xleftarrow{u} x
\]
% foot: Version 9 septembre 2001
\printedpage{17}
% head: A. Grothendieck : Dérivateurs, Ch. V, Édition des Universités de Montpellier II et Paris VII
de $y$ dans $x$, mais ce n'est pas un chemin au sens de nos définitions, on va le compléter en
\[
c : y \xrightarrow{\mathrm{id}} y \xleftarrow{u} x \xrightarrow{\mathrm{id}} x \xleftarrow{\mathrm{id}} x.
\]
Je dis que $c$ fait l'affaire. On doit prouver deux relations d'homotopie, de $c \cdot \bar{u}$ et de $\bar{u} \cdot c$ avec les chemins identiques en $x$ \texttt{[resp. en} $y$\texttt{]} de longueur 8 (indice 4, i.e. dans $\uCh_{\Sigma,4}(X;x,x)$ resp. $\uCh_{\Sigma,4}(X;y,y)$),
\msp{page 27}où $\bar{u}$ est le chemin associé à $u$
\[
x \xrightarrow{u} y \xleftarrow{\mathrm{id}} y \xrightarrow{\mathrm{id}} y \xleftarrow{\mathrm{id}} y.
\]

On a donc à regarder
\[
x \xrightarrow{u} y \xleftarrow{\mathrm{id}} y \xrightarrow{\mathrm{id}} y \xleftarrow{\mathrm{id}} y \xrightarrow{\mathrm{id}} y \xleftarrow{u} x \xrightarrow{\mathrm{id}} x \xleftarrow{\mathrm{id}} x
\]
et voir qu'il est homotope au chemin trivial
\[
x \xrightarrow{\mathrm{id}} x \xleftarrow{\mathrm{id}} \cdots x \xrightarrow{\mathrm{id}} x \xleftarrow{\mathrm{id}} x
\]
(des identités partout), et itou pour
\[
y \xrightarrow{\mathrm{id}} y \xleftarrow{u} x \xrightarrow{\mathrm{id}} x \xleftarrow{\mathrm{id}} x \xrightarrow{u} y \xleftarrow{\mathrm{id}} y \xrightarrow{\mathrm{id}} y \xleftarrow{\mathrm{id}} y
\]
et le chemin trivial en $y$. Cette fois, on a encore deux foncteurs
\[
\begin{tikzcd}[column sep=large]
\bDelta_1 & \Phi_4 \arrow[l, shift right, "\varphi"'] \arrow[l, shift left, "\psi"] \\
& \varepsilon, \arrow[ul, "\alpha'"] \arrow[u, "\alpha"']
\end{tikzcd}
\]
compatibles avec objets sources et objets buts, et on doit prouver qu'ils sont homotopes dans $\uHom^!(\Phi_4, \bDelta_1)$. Cela doit donc résulter du même lemme général. On a gagné - par une interprétation simple de $X\Sigma^{-1}$, et en particulier des $\Hom_{X\Sigma^{-1}}(x,y)$ et de leur composition, en termes de la géométrie des espaces de chemins :
\[
\left\{
\begin{aligned}
&\Hom_{X\Sigma^{-1}} \simeq \pi_0(\uCh_\Sigma(X;x,y)) \\ % [sic] Hom_{XΣ^{-1}} -> Hom_{XΣ^{-1}}(x,y), arguments (x,y) missing
&\text{Composition des homomorphismes dans } X\Sigma^{-1} \text{ correspond à} \\
&\text{la composition des } \Sigma\text{-chemins dans } X
\end{aligned}
\right.
\]
% foot: Version 9 septembre 2001

\dgrblock{dgrb173}{Cancellation uses the same endpoint twice}
\begin{dcomment}{examples}{p27:Cette fois, on a encore deux foncteurs}{Cancellation uses the same endpoint twice}
Grothendieck has written down the correct inverse path. The only remaining work is to make its two cancellation witnesses explicit. Let $u:x\to y$ belong to $\Sigma$, and write $T:[1]\to X$ for the corresponding functor. In the numerical ordering $0<1$, the first eight-edge excursion on the page is the image of the row
\[
 a=(0,1,1,1,1,1,0,0,0),
\]
whereas the second is the image of
\[
 b=(1,1,0,0,0,1,1,1,1).
\]
The source is $\Phi_4$, numbered from $0$ to $8$. Both rows are functors: on a forward source edge the value can increase, and on a backwards edge it can decrease. Their endpoint values are respectively $(0,0)$ and $(1,1)$.

There are pointwise natural transformations
\[
 (0,0,0,0,0,0,0,0,0)\longrightarrow a,
 \qquad
 b\longrightarrow(1,1,1,1,1,1,1,1,1).
\]
Applying $T$ gives a transformation from the constant path at $x$ to the first excursion, and one from the second excursion to the constant path at $y$. Their endpoint components are identities. Each horizontal backwards edge in these diagrams is either an identity or $u$, so it belongs to $\Sigma$. Thus both transformations are morphisms in the required $\Sigma$-path fibres, and the two cancellation identities follow immediately on $\pi_0$.

There is a small mismatch in the manuscript's invocation of the general lemma. That lemma is worded for maps taking the initial endpoint to the initial endpoint and the final endpoint to the final endpoint. These cancellation diagrams instead take both endpoints to $0$, or both to $1$. The relative-fibre theorem above includes them, but the two displayed transformations already prove exactly what is needed. This is a missing specification of the boundary condition, with an elementary repair, rather than a failure of localization.
\end{dcomment}


\dgrblock{dgrb174}{Completing the universal property}
\begin{dcomment}{exposition}{p27:par une interprétation simple}{Completing the universal property}
Grothendieck's path category now gives the promised localization. We record the complete conclusion because the construction uses less structure than many familiar fraction calculi.

\textbf{Path presentation of ordinary localization.} Let $X$ be a small category and let $\Sigma$ contain all its identity arrows. There is a canonical isomorphism of categories over $X$
\[
 X[\Sigma^{-1}]\cong X_\Sigma,
 \qquad
 \operatorname{Hom}_{X_\Sigma}(x,y)
   =\pi_0\operatorname{Ch}_\Sigma(X;x,y),
\]
where the path categories use strict endpoint fibres and composition is induced by concatenation at even indices.

\emph{Proof.} The reparametrization argument proves independence of padding. More explicitly, the maps that compare padding before and after concatenation are numerically monotone endpoint-preserving functors between alternating intervals. The connectedness proved by even-block slides supplies a zigzag between them. After precomposition with a $\Sigma$-path, every backwards edge is still a backwards edge of that path or an identity. Thus the comparison is a zigzag inside the specified path category. Morphisms of the factors concatenate to morphisms of the result, so the product descends to components. Associativity follows from strict associativity of finite concatenation, after choosing common even indices; the zero-length paths give identities.

The preceding table proves that $\theta:X\to X_\Sigma$ preserves composition, and the cancellation ladders prove that it inverts $\Sigma$. Let $F:X\to Y$ invert $\Sigma$. Evaluate each finite path in $Y$ by using $F$ on forward arrows and inverses of $F(s)$ on backwards arrows. Identity padding leaves the result unchanged. Naturality squares, with the backwards ones inverted as above, show that endpoint-fixed ladders give equal arrows of $Y$. Evaluation therefore defines maps on the proposed Hom sets. It takes concatenation to composition and constant paths to identities, so it defines a functor $\overline F:X_\Sigma\to Y$ with $\overline F\theta=F$.

To prove uniqueness, one must know that a path class is the composite of its elementary arrows and inverse arrows. Pad the original path to an even index. Replace each forward arrow by its four-edge representative under $\theta$, and each backwards arrow by the four-edge inverse representative. Concatenating these representatives gives the original list of nonidentity edges with additional identities inserted. It is the pullback of the original path by a numerically monotone contraction from a larger alternating interval. The usual end-padding is another such pullback. The two contractions belong to the same connected endpoint-preserving category, so the two paths have the same component. Hence every path class is the required composite, and any extension of $F$ must evaluate it as $\overline F$ does.

Thus $\theta$ has the defining universal property of localization. Applying the two universal properties gives inverse functors between $X_\Sigma$ and $X[\Sigma^{-1}]$ that agree with the identity on $X$. Their composites are identities by uniqueness.\hfill$\square$

The proof explains the exact status of the manuscript: its principal theorem is correct, and its proof becomes complete after the endpoint-fibre convention, the even-index repair already made by the author, and the supplied interval comparisons. No Ore condition or closure of $\Sigma$ under composition has entered.
\end{dcomment}


\dgrblock{dgrb175}{From components of path categories to mapping spaces}
\begin{dcomment}{later}{p27:composition, en termes de la géométrie des espaces de chemins}{From components of path categories to mapping spaces}
Grothendieck has recovered a set of morphisms by taking connected components of a category of presentations. The next question is whether one can retain higher homotopy information and compose it coherently. The relevant historical comparison already existed before these notes: the two papers of Dwyer and Kan appeared in 1980.

Their standard simplicial localization has the ordinary localization as its category of components \citep[\S4.2]{DwyerKan1980Simplicial}. The hammock construction gives a more diagrammatic model: its simplices are commutative ladders of variable length, with all vertical arrows and all leftward horizontal arrows in the chosen subcategory of weak equivalences. Composition is concatenation followed by reduction \citep[Definition~2.1 and Proposition~3.1]{DwyerKan1980Calculating}. These are close relatives of the present paths, but the vertical-arrow requirement is different: Grothendieck permits arbitrary vertical arrows. Therefore the formula for $\pi_0$ proved here does not itself identify the nerve of his path category with a derived mapping space.

There is also a preliminary distinction in the input data. The classical Dwyer--Kan construction starts with a subcategory $W$ containing all objects, whereas the present $\Sigma$ need contain only identities. To compare the two ordinary localizations one may take $W$ to be the subcategory generated by $\Sigma$: every functor inverting $\Sigma$ inverts its finite composites, and the converse is immediate. This proves equality of their universal properties. It does not identify the intervening categories of paths. The example with $\Sigma=\{1,t\}$ above shows concretely how adjoining composite denominators changes the allowed horizontal presentations before any homotopy comparison is proved.

The comparison becomes sharper under additional hypotheses. A homotopy calculus of fractions allows certain short zigzag categories to compute hammock mapping spaces; the exact reduction results are \citet[Proposition~6.2]{DwyerKan1980Calculating}. Their hypotheses concern the homotopy types of diagram categories, not just the existence of short ordinary fractions. Consequently a reduction that is sufficient for Hom sets need not preserve all higher homotopy groups.

One elementary check shows the difference in what is being sought. If every arrow of a category $X$ is inverted, the ordinary localization records only its groupoid completion. A higher localization retains the homotopy type of the classifying space: in the connected case the mapping spaces have the homotopy type of its loop space \citep[\S5.5]{DwyerKan1980Simplicial}. Take the poset of nonempty proper faces of a tetrahedron: its nerve is the barycentric subdivision of the boundary sphere $S^2$. Its groupoid completion is equivalent to the terminal groupoid, since the sphere is connected and simply connected, while the higher endomorphism space has the nontrivial homotopy type $\Omega S^2$. This example does not decide the homotopy type of Grothendieck's particular path category; it explains why agreement on its components is only the first comparison one would require.

The later development of simplicial and higher localizations gives a precise home for the ambition of retaining those spaces of presentations. It does not retroactively turn every category of zigzags into the same model. The admissible horizontal arrows, the vertical arrows, the allowed reparametrizations, and the endpoint conditions are part of the theorem to be proved.
\end{dcomment}


\printedpage{18}
% head: A. Grothendieck : Dérivateurs, Ch. V, Édition des Universités de Montpellier II et Paris VII
\noindent\msp{page 28}

\section*{3\quad Passage aux catégories de diagrammes}

Nous aurions dû noter la catégorie $X$ $\mathcal{M}$, car par la suite elle jouera le rôle d'une `catégorie de modèles'. On désignera par contre par $X$ les `catégories d'indices' sur lesquelles on veut construire les `catégories de diagrammes' de type $X$
\[
\mathcal{M}(X) = \uHom(X,\mathcal{M}).
\]
On se donne donc
\[
\underbrace{\Sigma}_{\text{quasi-isomorphismes}} \subseteq \mathrm{Fl}(\mathcal{M}), \qquad \Sigma \supset \text{ensemble des flèches identiques},
\]
et pour toute catégorie $X$ on définit
\[
\Sigma_X \subset \mathrm{Fl}\,\mathcal{M}(X)
\]
ainsi: si $F, G \in \mathrm{Ob}\,\mathcal{M}(X)$ ($F, G : X \rightrightarrows \mathcal{M}$), une flèche $u : F \longrightarrow G$ (homomorphisme de foncteurs) est dans $\Sigma_X$ si et seulement si $\forall x \in X$, $u(x) : F(x) \longrightarrow G(x)$ est dans $\Sigma$ (`critère de quasi-isomorphie fibre par fibre', pour l'homomophisme entre les copréfaisceaux $F$, $G$ envisagés). % [sic] homomophisme -> homomorphisme, letter r missing

\dgrblock{dgrb176}{The parameter is a diagram shape}
\begin{dcomment}{exposition}{p28:On désignera par contre}{The parameter is a diagram shape}
Grothendieck now asks the path construction to do more than compute a single localized Hom-set. He wants the same calculation after replacing the coefficient category by every category of diagrams. This is the point at which a category of paths begins to behave like an external Hom: its test objects are the indexing categories themselves. The change of letters is deliberate. In this section $\mathcal M$ supplies coefficients, whereas $X$ is allowed to vary as a parameter.

There is no model-category structure in the hypotheses. We have a category $\mathcal M$ and a class $\Sigma$ containing the identities; for the localization calculation even closure of $\Sigma$ under composition is unnecessary. One may enlarge $\Sigma$ to the subcategory it generates without changing the localization. We work in fixed universes, so that the localizations and the Hom-sets under discussion exist in the chosen larger universe. The designation ``quasi-isomorphism'' means membership in this specified class, rather than a statement about chain complexes.

Write $\mathcal M^X=\operatorname{Fun}(X,\mathcal M)$ and let $\Sigma_X$ consist of natural transformations whose components belong to $\Sigma$. Restriction along $f:X\to Y$ preserves these transformations and therefore induces
\[
 f^*: \mathcal M^Y[\Sigma_Y^{-1}]
       \longrightarrow \mathcal M^X[\Sigma_X^{-1}].
\]
The two objects of special interest are the constant diagrams $a_X,b_X$. Thus the functor to be understood is
\[
 H_{a,b}(X)=\operatorname{Hom}_{\mathcal M^X[\Sigma_X^{-1}]}(a_X,b_X),
 \qquad H_{a,b}:\Cat^{\mathrm{op}}\to\Ens.
\]
These are constant diagrams before localization. We are not replacing an arbitrary diagram by its values in $\mathcal M[\Sigma^{-1}]$; that would discard precisely the coherent information retained in Chapters~I and~II. See the commentary on Chapter~I, manuscript p.~13 (``Localizing diagrams and deriving their Kan extensions'') for that distinction.

For a fixed finite path shape $\Phi_n$, the elementary exponential law
\[
 \operatorname{Fun}(\Phi_n,\operatorname{Fun}(X,\mathcal M))
 \cong\operatorname{Fun}(X,\operatorname{Fun}(\Phi_n,\mathcal M))
\]
already contains the calculation. Its objects agree because both sides are functors $X\times\Phi_n\to\mathcal M$; its arrows agree because both sides are natural transformations of such functors. The restrictions on backward arrows are checked at each object of $X$, and the restrictions on the endpoints are also objectwise. The difficulty is therefore not the exponential law. It is the passage from one finite path shape to all path lengths, where one must retain a uniform bound.
\end{dcomment}


Soient
\[
a, b \in \mathrm{Ob}\,\mathcal{M}
\]
et considérons les foncteurs constants correspondants sur $X$
\[
a_X, b_X \in \mathrm{Ob}\,\mathcal{M}(X).
\]

\msp{page 29}On se propose d'appliquer la théorie du \S2 au `calcul' de
\[
\Hom_{\Sigma_X}(a_X, b_X),
\]
on a donc une bijection canonique
\[
\boxed{\Hom_{\Sigma_X}(a_X, b_X) \simeq \pi_0\uCh_{\Sigma_X}(\mathcal{M}(X); a_x, b_x).} % [sic] a_x, b_x -> a_X, b_X, lowercase index against a_X, b_X elsewhere
\]
On se propose d'expliciter la catégorie de chemins du second membre. On a, abstraction faite des données de $\Sigma$, $a$, $b$ dans $\mathcal{M}$, l'expression pour la catégorie de \textit{tous} les chemins
% foot: Version 9 septembre 2001

% ---- from Der-ChV-234_p021-025.tex
\printedpage{19}
% head: A. Grothendieck : Dérivateurs, Ch. V, Édition des Universités de Montpellier II et Paris VII

dans $\cM(X)$
\[
\begin{aligned}
\Ch(\cM(X)) \quad &\simeq \quad \uHom(\ \underbrace{\boldsymbol{\Phi}}_{\substack{\text{pro-catégorie}\\ (\Phi_n)_{n\in\mathbf{N}}}}\ ,\ \underbrace{\cM(X)}_{\uHom(X,M)}\ ) \\ % reading resolved: italic M is printed; category-valued Hom underlines restored.
&\dfeq \quad \varinjlim_n \uHom(\Phi_n,\cM(X)) \\
&\simeq \quad \varinjlim_n \uHom(\Phi_n\times X,\cM) \\
&\simeq \quad \varinjlim_n \uHom(X,\uHom(\Phi_n,\cM)) \\
&\dfeq \quad \varinjlim_n \uHom(X,\Ch_n(\cM)) \\
&\simeq \quad \uHom^{!}(X,\Ch(\cM)),
\end{aligned}
\]
où le signe ! signifie qu'on prend la sous-catégorie pleine des foncteurs
\[
X\longrightarrow \Ch(\cM),
\]
\msp{page 30}
qui se factorisent par une des sous-catégories pleines $\Ch_n(\cM)$. On peut dire aussi que si on considère le système inductif des catégories $\Ch_n(\cM)$ comme un ind-objet de Cat, qu'on va noter
\[
\Chh(\cM) \qquad\qquad\qquad\qquad \text{ind-objet de Cat,}
\]
on a
\[
\Ch(\cM(X)) \xrightarrow{\ \sim\ } \uHom(X,\Chh(\cM)).
\]
Si on en vient à $\Sigma$, $\Sigma_X$, on trouve aussitôt qu'un chemin dans $\cM(X)$ est un $\Sigma_X$-chemin, si et seulement si le foncteur correspondant $X\longrightarrow\Ch(\cM)$ se factorise par $\Ch_{\Sigma}(\cM)$. Ainsi on trouve
\[
\Ch_{\Sigma_X}(\cM(X))\simeq\uHom^{!}(X,\Ch_{\Sigma}(\cM)),
\]
où le signe ! a le même sens que plus haut. Si on introduit le ind-objet $\Chh_{\Sigma}(\cM)$ défini par les $\Ch_{\Sigma,n}(\cM)$, on trouve donc aussi
\[
\Ch_{\Sigma_X}(\cM(X))\xrightarrow{\ \sim\ }\uHom(X,\Chh_{\Sigma}(\cM)).
\]

\dgrblock{dgrb177}{Why the exclamation mark is necessary}
\begin{dcomment}{exposition}{p30:qui se factorisent par une}{Why the exclamation mark is necessary}
Grothendieck introduces the exclamation mark to remember that a path of diagrams has one length for the entire diagram. An object of $\operatorname{Fun}(X,\operatorname{Ch}_\Sigma(\mathcal M))$ may assign a finite path to every $x\in X$ without these lengths being bounded as $x$ varies. Such a functor need not come from a path in $\mathcal M^X$. This is a genuine distinction even before taking connected components or weak equivalences.

Let $K_n$ denote the fixed-endpoint path category at stage $n$, using the fully faithful identity-padding inclusions, and put $K_\infty=\bigcup_nK_n$. Then
\[
 \operatorname{Fun}^{!}(X,K_\infty)
   =\bigcup_n\operatorname{Fun}(X,K_n)
\]
is the full subcategory on bounded functors. Fullness deserves a verification. If $F$ and $G$ factor through $K_n$, every component of a natural transformation $F\Rightarrow G$ lies in $K_n$, because $K_n\subseteq K_\infty$ is full. Its naturality equations are the same equations in either category. Thus no additional bound is needed for arrows between already bounded objects.

The failure for an infinite source can be seen without any homotopy theory. Take $K_n$ to be the discrete category on $\{0,\ldots,n\}$ and $X$ the discrete category on $\mathbf N$. The functor $x\mapsto x$ from $X$ to $K_\infty$ lies in no $\operatorname{Fun}(X,K_n)$. It cannot even be isomorphic to a bounded functor, since the target is discrete. The ind-object $(K_n)$ records a uniform finite stage which the ordinary union has forgotten.

Connected components do commute with the filtered union used here. A component of the union has a representative object at some stage. If two objects become connected, a finite zigzag of arrows witnesses this, and all its objects and arrows occur at one sufficiently large stage. This proves directly that
\[
 \varinjlim_n\pi_0\operatorname{Fun}(X,K_n)
 \xrightarrow{\sim}
 \pi_0\!\left(\varinjlim_n\operatorname{Fun}(X,K_n)\right).
\]
Nothing in this argument permits one to replace the category inside the final parentheses by the unrestricted functor category. The exclamation mark and the double underline serve the same purpose: the first describes bounded functors, while the second retains the system of bounds as an ind-object.
\end{dcomment}


\msp{page 31}
Soit $c$ un chemin dans $\cM(X)$, et
\[
X\longrightarrow\Ch(\cM)
\]
le foncteur qu'il définit, composons-le avec les foncteurs source-but
\[
\Ch(\cM)\longrightarrow\cM\times\cM,
\]
on trouve un foncteur
\[
X\longrightarrow\cM\times\cM,
\]
i.e. un objet $(F_0,F_1)$ de $\cM(X)\times\cM(X)$. Je dis que $F_0$ et $F_1$ sont respectivement la source et le but du chemin $c$. De façon un peu plus précise, on a commutativité (stricte) dans le diagramme de foncteurs
% foot: Version 9 septembre 2001
\printedpage{20}
% head: A. Grothendieck : Dérivateurs, Ch. V, Édition des Universités de Montpellier II et Paris VII
\[
\begin{tikzcd}[column sep=large, row sep=large]
\Ch_{\Sigma_X}(\cM(X)) \arrow[r, "\sim"] \arrow[d, "{\sigma_{\cM(X)},\,\beta_{\cM(X)}}"] & \uHom^{!}(X,\Ch_{\Sigma}(\cM)) \arrow[d, "{\text{déduit de }\sigma_{\cM},\,\beta_{\cM}}"] \\
\cM(X)\times\cM(X) \arrow[r, "\sim"] & \uHom(X,\cM\times\cM),
\end{tikzcd}
\]
où les foncteurs $\sigma$ sont les foncteurs source, les foncteurs $\beta$ les foncteurs but.

Si on se donne arbitrairement les objets $F_0$, $F_1$ dans $\cM(X)$, donc un foncteur
\[
X\longrightarrow\cM\times\cM,
\]
$\Ch_{\Sigma_X}(F_0,F_1)$ s'identifie donc à la catégorie des foncteurs ‘admissibles’ (i.e. se factorisant par un $\Ch_{\Sigma,n}(\cM)$) $X\xrightarrow{\ C\ }\Ch_{\Sigma}(\cM)$ qui \textit{relèvent}
\msp{page 32 vide}
\msp{page 33}
le foncteur $(\sigma_{\cM},\beta_{\cM})$ \texttt{[plutôt} $(F_0,F_1)$\texttt{]} précédent, i.e. qui donnent lieu à un triangle commutatif (strict)
\[
\begin{tikzcd}[row sep=large]
 & \Ch_{\Sigma}(\cM) \arrow[d, "{(\sigma_{\cM},\,\beta_{\cM})}"] \\
X \arrow[ur, "C"] \arrow[r, "{(F_0,F_1)}"] & \cM\times\cM.
\end{tikzcd}
\]
Mais si les foncteurs $F_0$, $F_1$ sont constants, de valeur $a$, $b$ comme tantôt, i.e. $(F_0,F_1)$ constant de valeur $(a,b)$, alors la catégorie des foncteurs $C$ en question est tautologiquement isomorphe à celle des foncteurs de $X$ dans la fibre de $(\sigma_{\cM},\beta_{\cM})$ en $(a,b)$, d'où l'isomorphisme de catégories fondamental
\[
\boxed{\Ch_{\Sigma_X}(X;a_X,b_X)\simeq\uHom^{!}(X,\Ch_{\Sigma}(\cM;a,b)).}
\]

On va expliciter le contenu essentiel de cette formule en une

\dgrblock{dgrb178}{The precise meaning of ind-representation}
\begin{dcomment}{exposition}{p33:On va expliciter le contenu essentiel}{The precise meaning of ind-representation}
The scholium isolates what survives every arbitrary choice in the path calculation. If $\Cathot$ denotes the category with small categories as objects and
\[
 \operatorname{Hom}_{\Cathot}(X,Y)
       =\pi_0\operatorname{Fun}(X,Y),
\]
then the result says that $H_{a,b}$ is a filtered colimit of representable presheaves on $\Cathot$. The homotopy relation in this formula is generated by natural transformations, with arrows in either direction. Composition is well defined because whiskering preserves natural transformations.

For clarity, the fully faithful embedding of ind-objects in presheaves can be proved in the notation needed here. To an ind-object $K=(K_i)$ of a locally small category $C$ associate
\[
 h_K(X)=\varinjlim_i\operatorname{Hom}_C(X,K_i).
\]
If $L=(L_j)$ is another ind-object, Yoneda and the universal property of a colimit give
\[
 \operatorname{Nat}(h_K,h_L)
 \cong\varprojlim_i\operatorname{Nat}(h_{K_i},h_L)
 \cong\varprojlim_i\varinjlim_j\operatorname{Hom}_C(K_i,L_j)
 =\operatorname{Hom}_{\operatorname{Ind}(C)}(K,L).
\]
The limits and colimits occur in that order. No interchange theorem has been used. A natural isomorphism of the presheaves consequently gives a unique isomorphism of the corresponding ind-objects, compatible with the specified representations. This is the exact sense of the uniqueness statement on manuscript p.~34.

To obtain the scholium, apply the preceding path-localization theorem inside $\mathcal M^X$, then apply the exponential law with fixed endpoints, and finally commute $\pi_0$ with the filtered union. This yields
\[
 H_{a,b}(X)
  \cong\varinjlim_n\operatorname{Hom}_{\Cathot}(X,K_n)
  =h_{(K_n)}(X).
\]
The notation in the boxed formula just before the scholium has $X$ where the ambient path category should be $\mathcal M(X)$; the edition prints this slip. The preceding endpoint square determines the intended meaning unambiguously.

The result has a useful limit. It gives uniqueness in $\operatorname{Ind}(\Cathot)$, not uniqueness of the ordinary union $K_\infty$ in $\Cathot$. A colimit of a chosen sequence of actual categories contains information about the representatives and their compatibility which an ind-object in a homotopy category need not retain. The questions at the end of the chapter begin exactly at this boundary.
\end{dcomment}


\textbf{Scholie.} Soient $\cM$ une catégorie avec un ensemble $\Sigma$ de flèches de localisation, $a,b\in\mathrm{Ob}\,\cM$. Alors on peut trouver un objet $\KK=K(a,b)$, $K=(K_n)$, de Ind\,Cat (indexé par $\mathbf{N}$ mais peu importe ce fait), et un isomorphisme canonique pour toute catégorie $X$ (isomorphisme fonctoriel en $X$)
\[
\begin{aligned}
\mathrm{Hom}_{\Sigma_X}(a_X,b_X) \quad &\simeq \quad \pi_0(\uHom(X,\KK)) \\
&\dfeq \quad \pi_0(\ \varinjlim \uHom(X,K_n)) \\
&\simeq \quad \varinjlim \pi_0(\uHom(X,K_n)).
\end{aligned}
\]
% foot: Version 9 septembre 2001

\printedpage{21}
% head: A. Grothendieck : Dérivateurs, Ch. V, Édition des Universités de Montpellier II et Paris VII

\msp{page 34}
Je vais épiloguer sur ce résultat, sans me soucier pour l'instant de la façon particulière dont la ind-catégorie $K$ \texttt{[plutôt} $\KK$\texttt{]} a été obtenue en termes de $\cM$, $\Sigma$, $a$, $b$.

Comme les $\pi_0\uHom(X,K_n)=\mathrm{Hom}_{\Cathot}(X,K_n)$ ne dépendent que de $X$ et de $K_n$ en tant qu'objets de Cat\,hot, je vais interpréter $\KK$ comme un ind-objet de Cat\,hot, et écrire la formule fondamentale sous la forme
\[
\mathrm{Hom}_{\Sigma_X}(a_X,b_X)\simeq\mathrm{Hom}_{\mathrm{Ind\,Cat\,hot}}(X,\KK).
\]

Ainsi, l'énoncé du scholie implique deux choses :

\begin{enumerate}
\item[a)] Le contrafoncteur $X\longmapsto\mathrm{Hom}_{\Sigma_X}(a_X,b_X)$ sur Cat se factorise canoniquement par Cat\,hot, via un foncteur canonique
\[
\Cathot^{o}\ \xrightarrow{\ H=H^{\cM,\Sigma}_{a,b}\ }\mathrm{Ens}\ ;
\]
\item[b)] Ce dernier contrafoncteur sur Cat\,hot est \textit{ind-représentable} par un ind-objet $\KK$ de Cat\,hot.
\end{enumerate}

C'est cet ind-objet de Cat\,hot qui est défini à isomorphisme unique près par $\cM$, $\Sigma$, $a$, $b$, indépendamment de tout artifice de calcul pour expliciter les $\mathrm{Hom}_{\Sigma_X}(a_X,b_X)$.

Notons aussi la conclusion suivante, nullement évidente à priori :

\msp{page 35}
Notons d'abord (ceci étant implicite dans le scholie) que toute flèche dans Cat
\[
X\xrightarrow{\ f\ }Y
\]
définit
\[
\cM(Y)\xrightarrow{\cM(f)}\cM(X) \qquad\qquad\qquad \text{appliquant } \Sigma_Y \text{ dans } \Sigma_X,
\]
d'où
\[
\cM(Y)\Sigma_Y^{-1}\dfeq\mathbf{D}_{\cM,\Sigma}(Y)\longrightarrow\cM_X\Sigma_X^{-1}=\mathbf{D}_{\cM,\Sigma}(X),
\]
\texttt{[plutôt} $\cM(X)$\texttt{]} d'où, pour tout $a,b\in\mathrm{Ob}\,\cM$, vu que $\cM(f)$ transforme $a_Y$, $b_Y$ en $a_X$, $b_X$ :
\[
\boxed{\mathrm{Hom}_{\Sigma_Y}(a_Y,b_Y)\longrightarrow\mathrm{Hom}_{\Sigma_X}(a_X,b_X).}
\]

Ceci posé, on a :

\textbf{Corollaire :} \textit{Si $f$ est un homotopisme, alors quelque soient $a,b\in\cM$, la flèche précédente est un isomorphisme.} % [sic] quelque soient -> quels que soient, grammatical form of the locution

La question naturelle, c'est de savoir si la même conclusion est encore valable, en supposant seulement que $f$ soit un ‘\textit{quasi-isomorphisme}’ (ou ‘équivalence faible’), ici que $f$ définisse un isomorphisme dans la catégorie de fractions Hot de Cat. (À ne pas confondre avec Cat\,hot,
% foot: Version 9 septembre 2001
\printedpage{22}
% head: A. Grothendieck : Dérivateurs, Ch. V, Édition des Universités de Montpellier II et Paris VII
\msp{page 36}
nettement plus fine). S'il en était ainsi, le contrafoncteur
\[
X\longmapsto\mathrm{Hom}_{\Sigma_X}(a_X,b_X)
\]
se factoriserait non seulement par Cat\,hot, mais même par la catégorie de fractions plus grossière Hot. Et dès lors se poserait la question si ce contrafoncteur hypothétique
\[
\mathrm{Hot}^{o}\longrightarrow\mathrm{Ens}
\]
est encore ind-représentable; et de préférence (\footnote[5]{ce sera automatique si on a ind-représentabilité comme foncteur sur Hot, cf. ci-dessous.}) ind-représentable par le ind-objet image de $\KK$ dans Ind\,Hot par
\[
\mathrm{Ind}(\Cathot)\longrightarrow\mathrm{Ind}(\mathrm{Hot}),
\]
déduit du foncteur de localisation $\Cathot\longrightarrow\mathrm{Hot}$. (Car cet ind-objet est un candidat tout trouvé.)

\dgrblock{dgrb179}{Two homotopy invariance questions}
\begin{dcomment}{exposition}{p35:La question naturelle, c'est de savoir}{Two homotopy invariance questions}
Grothendieck has proved invariance under categorical homotopy equivalence, and now asks for invariance under weak equivalence of nerves. The first assertion is immediate from the ind-representation. A natural transformation $f\Rightarrow g$ induces, by precomposition, a natural transformation between the two functors on each $\operatorname{Fun}(Y,K_n)$. These functors therefore induce the same map on $\pi_0$. A finite zigzag gives the same conclusion, and a homotopy inverse to $f$ gives inverse maps on $H_{a,b}$.

The second assertion has additional content. A weak equivalence $Nf$ need not admit an inverse represented by a functor and a finite zigzag of natural transformations. Nor does the construction of $H_{a,b}$, for an arbitrary class $\Sigma$, supply a derived mapping space. Thus the corollary proves factorization through $\Cathot$; it does not prove factorization through $\Hot$. This is the same difference that appeared in the commentary on Chapter~I, manuscript p.~102 (``What the completion requires''), now applied to a particular family of coefficients.

Once that factorization is established, however, the subsequent question about ind-representability has a formal affirmative answer. It does not require the proposed right adjoint to $\Cathot\to\Hot$. More generally let $\varphi:C\to D$ be a localization, let $H=\varinjlim_i h_{K_i}$ on $C$, and suppose $H=\varphi^*H'$ for a presheaf $H'$ on $D$. Restriction $\varphi^*$ is fully faithful by the universal property of localization. Put $L=\varinjlim_i h_{\varphi K_i}$ on $D$. For every presheaf $F$ on $D$, Yoneda gives natural bijections
\[
 \operatorname{Nat}(L,F)
 \cong\varprojlim_i F(\varphi K_i)
 \cong\operatorname{Nat}(H,\varphi^*F)
 \cong\operatorname{Nat}(H',F).
\]
Yoneda, now in the presheaf category, proves $L\cong H'$. Equivalently, left Kan extension along $\varphi^{\mathrm{op}}$ sends $h_K$ to $h_{\varphi K}$, preserves the filtered colimit, and recovers a descended presheaf; this is the usual Kan-extension formalism \citep[\S\S1.1--1.4]{Riehl2014}. Thus, conditional only on weak-equivalence invariance, the image of $(K_i)$ in $\operatorname{Ind}(\Hot)$ represents the descended functor. What remains conditional on the later right adjoint is the characterization by a Thomason--Kan envelope, not this descent of ind-representability.

A useful sufficient hypothesis is that localization of the strict diagram categories presents the expected infinity-categories of diagrams. More explicitly, suppose a homotopy theory $\mathcal C$ satisfies, naturally in small $X$,
\[
 \mathcal M^X[\Sigma_X^{-1}]
   \simeq h\operatorname{Fun}(NX,\mathcal C).
\]
If $Z=\operatorname{Map}_{\mathcal C}(a,b)$, the space of natural transformations between constant diagrams is
\[
 \operatorname{Map}_{\operatorname{Fun}(NX,\mathcal C)}(a_X,b_X)
   \simeq\operatorname{Map}_{\mathcal S}(|NX|,Z).
\]
Indeed, a coherent transformation chooses a point of $Z$ at each object, a path for each arrow, and the compatible higher homotopies for composable strings; these are exactly the simplices of a map from $NX$ into $Z$. Taking $\pi_0$ proves $H_{a,b}(X)\cong[|NX|,Z]$. Under this hypothesis the functor is represented on $\Hot$, rather than merely ind-represented there.

Combinatorial simplicial model categories provide a familiar setting for this statement. The mapping space must be derived, using a cofibrant source and a fibrant target. The model-independent function-complex construction is due to Dwyer and Kan; their comparison theorem explains why the resulting homotopy type is independent of auxiliary model structures \citep[\S4]{DwyerKan1980Function}. The displayed hypothesis is stated explicitly here because it is not a consequence of an arbitrary choice of $\Sigma$.
\end{dcomment}


\dgrblock{dgrb180}{An external Hom that is an ordinary homotopy type}
\begin{dcomment}{examples}{p36:est encore ind-représentable}{An external Hom that is an ordinary homotopy type}
Grothendieck's candidate can be made completely concrete for simplicial sets with their usual weak equivalences. Take $a=\Delta^0$ and let $b=Z$ be a Kan complex. For a small category $X$, a projectively cofibrant replacement of the constant point diagram is supplied by the usual bar construction. At $x\in X$ its value is $N(X/x)$; this is contractible because $X/x$ has the terminal object $\operatorname{id}_x$. The diagram is assembled from the covariant representables by attaching simplices along their boundaries, so it is projectively cofibrant. The constant diagram $Z_X$ is projectively fibrant.

A natural transformation from this resolution to $Z_X$ is equivalent to a simplicial map from its colimit to $Z$. The colimit of the diagram $x\mapsto N(X/x)$ is $NX$: a simplex represented by a string ending in an arrow to $x$ is identified with the same string at its last vertex. Consequently the derived simplicial mapping space between the two constant diagrams is $Z^{NX}$. Since the source simplicial set is cofibrant and $Z$ is Kan,
\[
 \operatorname{Hom}_{\operatorname{Ho}(\mathbf{sSet}^{X})}
       (*_X,Z_X)=\pi_0 Z^{NX}=[NX,Z].
\]
This is the bar-construction calculation underlying the abstract assertion in the preceding block; see \citet[\S\S4.4--4.5 and 5.1]{Riehl2014} for the construction and its homotopical justification.

For $Z=K(A,m)$ with $m\geq1$ and $A$ abelian, it gives $H_{*,Z}(X)=H^m(NX;A)$, with constant coefficients. If $X$ is the one-object category attached to the cyclic group $C_r$, the familiar periodic resolution has alternating differentials $g-1$ and $1+g+\cdots+g^{r-1}$. After applying $\operatorname{Hom}_{\mathbf Z[C_r]}(-,A)$ for the trivial action, these differentials become $0$ and multiplication by $r$. Hence
\[
 H^{2q}(BC_r;A)=A/rA\quad(q\geq1),\qquad
 H^{2q+1}(BC_r;A)=A[r]\quad(q\geq0).
\]
The category $BC_r$ has finitely many objects and arrows, but its nerve has nondegenerate simplices in arbitrarily high dimensions. Thus ``finite category'' in the closing question must not be silently replaced by ``finite CW complex''. The same example also shows that successful homotopy-theoretic coefficients often produce an actually representable functor on $\Hot$. This is distinct from essential constancy in $\operatorname{Ind}(\Cathot)$, the subject of the pessimistic parenthesis on manuscript p.~46. Passing from the first assertion to the second would require precisely the right adjoint discussed next.
\end{dcomment}


\msp{page 37}
Je vais faire quelques rappels catégoriques, généralisant la situation du foncteur canonique de localisation
\[
\Cathot\xrightarrow{\ \varphi\ }\mathrm{Hot},
\]
dont on sait qu'il admet un adjoint à droite \texttt{[si un lecteur a une démonstration de ce fait, prière de contacter G. Maltsiniotis : maltsin@math.jussieu.fr]}
\[
\Cathot\xleftarrow{\ \psi\ }\mathrm{Hot}, \qquad\qquad\qquad \mathrm{Hom}_{\mathrm{Hot}}(\varphi X,Y)\simeq\mathrm{Hom}_{\Cathot}(X,\psi Y),
\]
où $\psi$ est le foncteur ‘enveloppe de Kan’ ou ‘représentant de Kan’. On va donc considérer la situation plus générale
\[
\begin{tikzcd}[column sep=large]
C \arrow[r, shift right=1, "\varphi"'] & C' \arrow[l, shift right=1, "\psi"']
\end{tikzcd}
\qquad\qquad \mathrm{Hom}_{C'}(\varphi X,X')\simeq\mathrm{Hom}_{C}(X,\psi X')
\]
i.e.
\[
\varphi,\psi \text{ couple de foncteurs adjoints.}
\]
On suppose de plus
\[
\left\{
\begin{aligned}
&\varphi \text{ un foncteur localisation, i.e. } \psi \text{ pleinement fidèle,}\\
&\text{i.e. } \varphi\psi\xrightarrow[\sim]{\text{adj}}\mathrm{id}_{C'} \quad .
\end{aligned}
\right.
\]
Soit $W=W_{\varphi}\subset\mathrm{Fl}\,C$, défini par $\varphi$ : $W_{\varphi}=\{u\in\mathrm{Fl}\,C\mid\varphi(u)\text{ isomorphisme}\}$. Voici la proposition que j'avais oubliée, et que j'ai eu du mal à reconstituer, hélas!

\dgrblock{dgrb181}{The reflection does not exist}
\begin{dcomment}{development}{p37:dont on sait qu'il admet un adjoint}{The reflection does not exist}
Grothendieck wants to convert homotopy invariance of a represented functor into a condition on its representing object. For a reflective localization this is an elementary and powerful argument, and Propositions~1 and~2 below are correct formal statements about any reflective localization. The edition questions the right adjoint to $\varphi:\mathrm{Cat}\,\mathrm{hot}\to\mathrm{Hot}$ on which their application rests. It does not exist.

\textbf{Theorem.}\ours
\textit{The functor $\varphi:\mathrm{Cat}\,\mathrm{hot}\to\mathrm{Hot}$ has no right adjoint. More precisely, there is no small category $K$ with bijections $\pi_0\operatorname{Fun}(X,K)\cong[BX,K(\mathbf Z/2,2)]$ natural in the small category $X$.}

\textit{Proof.}
A right adjoint $\psi$ would give such a $K$, namely $\psi$ of the homotopy type $Y=K(\mathbf Z/2,2)$, since morphisms $X\to K$ in $\mathrm{Cat}\,\mathrm{hot}$ are the connected components of $\operatorname{Fun}(X,K)$. Suppose $K$ is given. Taking $X=e$ shows that $\pi_0K=[\mathrm{pt},Y]$ is one point.

Let $\Omega$ be an infinite set of cardinality greater than that of the set of arrows of $K$, and $G$ the group of even permutations of $\Omega$ with finite support. It is the union of the alternating groups of the finite subsets of $\Omega$, so it is simple, and its cardinality is that of $\Omega$. A functor from the one-object category $\mathsf BG$ to $K$ is an object $k$ with a homomorphism $G\to\operatorname{Aut}(k)$. Its kernel is a normal subgroup; it is not trivial because $\operatorname{Aut}(k)$ has fewer elements than $G$; so the homomorphism is trivial. A natural transformation between two such constant functors is a single arrow of $K$. Hence $\operatorname{Fun}(\mathsf BG,K)\cong K$ and $\pi_0\operatorname{Fun}(\mathsf BG,K)$ is one point.

On the other side $[BG,K(\mathbf Z/2,2)]=H^2(G;\mathbf Z/2)$ classifies the central extensions of $G$ by $\mathbf Z/2$, and one of them does not split. Let $V$ be the real vector space with orthonormal basis $(e_\omega)_{\omega\in\Omega}$ and $\mathrm{Cl}(V)$ its Clifford algebra, with $v^2=-|v|^2$. For $\omega\neq\omega'$ the unit vector $\tau=(e_\omega-e_{\omega'})/\sqrt2$ is invertible in $\mathrm{Cl}(V)$ and $v\mapsto-\tau v\tau^{-1}$ is the reflection of $V$ that exchanges $e_\omega$ and $e_{\omega'}$. Let $\widehat G$ be the group of products of an even number of such vectors. Sending such a product to the corresponding product of transpositions is a surjective homomorphism $\widehat G\to G$; its kernel is contained in $\{\pm1\}$, the kernel of the map from the Pin group to the orthogonal group, and $-1=(\tau\tau')^2\in\widehat G$ for two orthogonal such vectors $\tau,\tau'$. So $\widehat G$ is a central extension of $G$ by $\mathbf Z/2$. The element $g=(\omega_1\,\omega_2)(\omega_3\,\omega_4)$ has order two, and its two lifts $\pm\tau\tau'$ have square $-1$; a splitting would lift $g$ to an element of order two. So $H^2(G;\mathbf Z/2)\neq0$, and the two sets cannot be in bijection.\hfill$\square$

The argument shows more than the statement: a category $K$ represents on $\mathrm{Cat}\,\mathrm{hot}$ only what homomorphisms into its automorphism groups and its arrows can see, and a group that is large and simple is invisible to it. The functors $X\mapsto[BX,Y]$ that are representable on all of $\mathrm{Cat}\,\mathrm{hot}$ are those with $Y$ a $1$-type, represented by a groupoid (the proof for a general $Y$, which replaces $G$ by large simple groups with the homology of spheres, is not reproduced here); the fundamental groupoid $X\mapsto\Pi_1X$ is a reflection of $\mathrm{Cat}\,\mathrm{hot}$ onto groupoids modulo natural isomorphism, and it inverts the functors that are bijective on $\pi_0$ and on all $\pi_1$, a larger class than the weak equivalences. Propositions~1 and~2 apply to that reflection.

Thomason's model structure gives a fibrant replacement $K\to RK$ and an equivalence between its homotopy category and $\mathrm{Hot}$ \citep[Theorem~4.9 and Corollary~5.5]{Thomason1980}. It does not identify $\pi_0\operatorname{Fun}(X,RK)$ with $[NX,NK]$ for every $X$: mapping formulas in a model category need a cofibrant source, and the cofibrant objects of that structure are posets. Below, references to an envelope or to a reflective localization of $\mathrm{Cat}\,\mathrm{hot}$ onto homotopy types are to be read as hypotheses that cannot be met; references to Thomason fibrant replacements mean the model-category construction and carry only its conclusions.
\end{dcomment}


\textbf{Proposition 1.} \textit{Soit $K\in \mathit{Ob}C$. Les conditions suivantes sur $K$ sont équivalentes :}
% foot: Version 9 septembre 2001

\printedpage{23}
% head: A. Grothendieck : Dérivateurs, Ch. V, Édition des Universités de Montpellier II et Paris VII

\begin{enumerate}
\item[(i)] \textit{Le foncteur}
\[
X\longmapsto K(X)\dfeq\mathrm{Hom}_C(X,K) \qquad\quad C^{o}\longrightarrow\mathrm{Ens}
\]
\textit{transforme flèches de $W$ (‘quasi-isomorphismes’) en bijections.}
\item[(i$'$)] \textit{Même condition pour les flèches d'adjonction}
\[
X\longrightarrow\psi\varphi X \qquad\quad (\forall X\in\mathrm{Ob}\ \mathrm{C})
\]
\item[(i$''$)] \textit{Même condition pour la flèche d'adjonction} (\footnote[6]{N.B. Il suffit même que l'élément $\mathrm{id}_K$ de $F_K(K)$ provienne de $F_K(\psi\varphi K)$ \texttt{[où} $F_K=K(-)=\mathrm{Hom}(-,K)$\texttt{]}, i.e.\par (i$'''$) Il existe une rétraction de $\psi\varphi X$ sur $X$ \texttt{[plutôt} $\psi\varphi K$ \texttt{sur} $K$\texttt{]} (NB Celle-ci sera même un isomorphisme).})
\[
K\longrightarrow\psi\varphi K.
\]
\item[(ii)] \textit{$K$ appartient à l'image essentielle de $\psi$}
\item[(ii$'$)] \textit{$K\longrightarrow\psi\varphi K$ (adjonction) est un isomorphisme.}
\item[(ii$''$)] \textit{L'homomorphisme fonctoriel en $X$}
\[
\mathrm{Hom}_C(X,K)\longrightarrow\mathrm{Hom}_{C'}(\varphi X,\varphi K)
\]
\textit{est un isomorphisme.}
\end{enumerate}

\textsc{Démonstration.} On a d'abord comme quasi-tautologies (ii$'$)$\Longrightarrow$ (ii)$\Longrightarrow$ (ii$''$)$\Longrightarrow$ (ii$'$), i.e. l'équivalence des trois dernières conditions. En effet, (ii$'$)$\Longrightarrow$ (ii) tautologique, (ii)$\Longrightarrow$ (ii$''$) car si $K=\psi(K')$ on a % reading resolved: single and double mathematical primes as printed.
\msp{page 38}
par adjonction
\[
\mathrm{Hom}_C(X,K)=\mathrm{Hom}_C(X,\psi(K'))\overset{\text{adj}}{\simeq}\mathrm{Hom}_{C'}(\varphi X,K'),
\]
or $\varphi(K)=\varphi\psi(K')\xrightarrow{\ \sim\ }K'$ (adjonction, et $\psi$ pleinement fidèle).
\[
\begin{tikzcd}[column sep=small, row sep=large]
\mathrm{Hom}_C(X,K) \arrow[rr, "\sim"] \arrow[dr, "\varphi_{X,K}"] & & \mathrm{Hom}_C(\varphi X,K') & \text{\footnotesize qui se complète en} \\ % [sic] Hom_C(\varphi X,K') -> Hom_{C'}(\varphi X,K'), \varphi X and K' are objects of C' (cf. the display above)
& \mathrm{Hom}_{C'}(\varphi X,\ \underbrace{\varphi K}_{=\varphi\psi(K')}\ ) \arrow[ur, "\sim", "\text{adj}"'] & &
\end{tikzcd}
\]
ce qui prouve que $\varphi_{X,K}$ est un isomorphisme.

(ii$''$) $\Longrightarrow$ (ii$'$). Car la flèche
% foot: Version 9 septembre 2001

% ---- from Der-ChV-234_p026-030.tex
\printedpage{24}
% head: A. Grothendieck : Dérivateurs, Ch. V, Édition des Universités de Montpellier II et Paris VII
% foot: Version 9 septembre 2001

\[
\begin{tikzcd}
\Hom_C(X,K) \arrow[r,"\alpha_X"] \arrow[dr,dashed] & \Hom_{C'}(\varphi X,\varphi K) \arrow[d,dash,"\wr"',"\mathrm{adj}"] \\
 & \Hom_C(X,\psi\varphi K)
\end{tikzcd}
\]
est déduite de la flèche d'adjonction
\[
K \longrightarrow \psi\varphi K
\]
en prenant $\Hom_C(X,-)$, donc $\alpha_X$ est un isomorphisme pour tout $X$ si et seulement si cette flèche d'adjonction est un isomorphisme. Donc
\[
\text{(ii)} \Longleftrightarrow \text{(ii$'$)} \Longleftrightarrow \text{(ii$''$)}.
\]
D'autre part,
\[
\text{(ii$''$)} \Longrightarrow \text{(i)} \Longrightarrow \text{(i$'$)} \Longrightarrow \text{(i$''$)}
\]
sont tautologiques, en tenant compte pour (i) $\Longrightarrow$ (i) \texttt{[probablement (i) $\Longrightarrow$ (i$'$)]} que les flèches d'adjonction
\[
X \xrightarrow{u} \psi\varphi X \qquad\qquad\qquad \text{dans } C
\]
sont bien dans $W$, car on a un triangle commutatif
\[
\begin{tikzcd}[column sep=large]
\varphi(X) \arrow[r,"\varphi(u)"] \arrow[dr,"\id_{\varphi(X)}"] & \varphi(\psi\varphi X) = (\varphi\psi)(\varphi X) \arrow[d,"{\begin{array}{l}\wr\ \mathrm{adj}\ \varphi\psi \longrightarrow 1,\ \text{qui est un}\\ \text{isomorphisme car } \psi \text{ est pleinement fidèle}\end{array}}"] \\
 & \varphi(X),
\end{tikzcd}
\]
donc $\varphi(u)$ est bien un isomorphisme.

Reste à prouver (i$''$) $\Longrightarrow$ (ii$'$), i.e. on suppose

\noindent\msp{page 39}%
que le foncteur $X \longrightarrow \Hom(X,K)$ transforme le morphisme d'adjonction particulier
\[
K \xrightarrow{i} \psi\varphi K
\]
en
\[
\begin{aligned}
u \quad &\longrightarrow \quad u\circ i\\
\Hom(\psi\varphi(K),K) \quad &\longrightarrow \quad \Hom(K,K),
\end{aligned}
\]
qui est une bijection. Donc $\id_K$ est dans l'image, soit $p$ un élément du premier membre qui lui donne naissance

\printedpage{25}
% head: A. Grothendieck : Dérivateurs, Ch. V, Édition des Universités de Montpellier II et Paris VII
% foot: Version 9 septembre 2001

\[
\begin{tikzcd}
K \arrow[r,hook,shift left,"i"] & \psi\,\underbrace{\varphi(K)}_{K'} \arrow[l,shift left,"p"]
\end{tikzcd}
\qquad\qquad {\scriptstyle \text{donc } p\circ i = \id_K}
\]
Considérons
\[
ip = q : \psi(K') \longrightarrow \psi(K') \qquad\qquad\qquad (\text{NB } q^2 = q),
\]
comme $\psi$ est pleinement fidèle, il existe un unique
\[
q_0 : K' \longrightarrow K' \qquad\qquad\qquad (K' \overset{\text{déf}}{=} \varphi(K))
\]
tel que l'on ait
\[
q = \psi(q_0).
\]
La relation $pi = \id_K$ donne, en appliquant $\varphi$,
\[
\varphi(p)\varphi(i) = \id,
\]
donc comme $\varphi(i)$ est un isomorphisme (c'est
\[
\begin{tikzcd}[column sep=large]
\varphi(K) \arrow[r,"\sim"] \arrow[dr,"\sim"{sloped,pos=0.45},"\id"{pos=0.6}] & \varphi(\psi\varphi)(K) = (\varphi\psi)(\varphi K) \arrow[d,"\wr"',"\mathrm{adj}"] \\
 & \varphi(K),
\end{tikzcd}
\]
rendant commutatif le triangle ci-contre), $\varphi(p)$ l'est. Donc $\varphi(q) = \varphi(i)\varphi(p)$ l'est, puisque $\varphi(i)$ et $\varphi(i)$ \texttt{[plutôt $\varphi(p)$]} le sont. Mais
\[
\varphi(q) = \varphi\psi(q_0) \quad \underbrace{\xrightarrow{\ \sim\ }}_{\mathrm{adj}\ \varphi\psi\ \xrightarrow{\sim}\ \id} \quad q_0,
\]
donc $q_0$ est un isomorphisme, donc aussi $q = \psi q_0$. Comme $q^2 = q$, cela implique $q = \id$. Donc $i$ et $p$ sont des isomorphismes inverses l'un de l'autre, ce qui établit (ii$''$) \texttt{[plutôt (ii$'$)]} et achève la démonstration.

\dgrblock{dgrb182}{A retract of a local object is local}
\begin{dcomment}{exposition}{p39:et achève la démonstration.}{A retract of a local object is local}
Grothendieck's proof of Proposition~1 identifies the precise amount of surjectivity needed: the identity of $K$ must lift across its unit. This is stronger than a convenient trick. It says that the full subcategory of local objects is closed under retracts, and explains why checking every weak equivalence can be reduced to one particular map.

Let $\varphi:C\rightleftarrows C':\psi$ be an adjunction with fully faithful right adjoint, and write $T=\psi\varphi$ and $\eta:1\to T$ for its idempotent localization monad. An object $K$ is local exactly when $\eta_K$ is invertible. Local objects invert $W_\varphi$ under $\operatorname{Hom}_C(-,K)$ by adjunction. Conversely, if precomposition with $\eta_K$ is surjective at the identity of $K$, choose $p:TK\to K$ with $p\eta_K=1_K$.

Here is a short verification of the remaining step. Since $\varphi\eta_K$ is invertible by the triangle identity, $\varphi p$ is its inverse. Thus $\varphi(\eta_Kp)=1_{\varphi TK}$. The functor $\varphi$ is faithful on the full subcategory of local objects: on that subcategory it is quasi-inverse to $\psi$. Both source and target of $\eta_Kp$ are the local object $TK$, so $\eta_Kp=1_{TK}$. Hence $\eta_K$ and $p$ are inverse isomorphisms. This proves the implication from the weakest condition in the footnote and therefore all the equivalences in the proposition.

The representability corollary is then exact. If $H=H'\varphi$, every $W_\varphi$ is sent to a bijection. When $H\cong\operatorname{Hom}_C(-,K)$, the proposition makes $K$ local, and adjunction identifies $H'$ with $\operatorname{Hom}_{C'}(-,\varphi K)$. In the other direction, a representative $K'$ for $H'$ pulls back to the representative $\psi K'$. Essential surjectivity of $\varphi$ onto $C'$ ensures that the resulting natural isomorphisms are determined by the original factorization. The repeated $X$ in the footnote's first formulation is a printed slip; the object at this step is $K$, as the edition itself corrects.
\end{dcomment}


\textbf{Corollaire :} \textit{Considérons des foncteurs $C^o \xrightarrow{H} \Ens$, et $C'^o \xrightarrow{H'} \Ens$, tels que $H \simeq H'\circ\varphi$. Pour que $H$ soit représentable, il faut et il suffit que $H'$ le soit, et dans ce cas, l'objet $K$ de $C$}

\noindent\msp{page 40}%
\textit{qui représente $H$ satisfait aux conditions équivalentes de la proposition. Si $K$, $K'$ les représentent, on a des isomorphismes canoniques $K' \simeq \varphi K$, $K \simeq \psi K'$.}

Si $H'$ est représentable par $K' \in \Ob C'$, on aura donc
\[
H'\circ\varphi(X) \simeq \Hom_{C'}(\varphi X,K') \overset{\mathrm{adj}}{\simeq} \Hom_C(X,\psi(K'))
\]
\printedpage{26}
% head: A. Grothendieck : Dérivateurs, Ch. V, Édition des Universités de Montpellier II et Paris VII
% foot: Version 9 septembre 2001
donc $H$ est représentable par $K = \psi(K')$, satisfaisant la condition (ii) de la proposition. Inversement, si $H$ est représentable par un objet $K$, alors $K$ satisfait à la condition (i), donc par (ii$''$) on a
\[
\Hom_C(X,K) \xrightarrow{\sim} \Hom_{C'}(\varphi X,\varphi K),
\]
donc le foncteur $H' : C^o \longrightarrow \Ens$ qui factorise $H$ s'explicite comme $X' \longmapsto \Hom_{C'}(X',\varphi K)$, donc est représentable par $\varphi K$. % [sic] C^o -> C'^o, H' is a functor on C'^o (cf. the Corollaire)

Passons aux ind-objets. Les foncteurs $\varphi$, $\psi$ passent aux ind-objets, d'où des foncteurs
\[
\begin{tikzcd}[column sep=large]
\Ind C \arrow[r,shift left,"\Psi = \Ind\psi"] \arrow[r,shift right,"\Phi = \Ind\varphi"'] & \Ind C'.
\end{tikzcd} % [sic] both arrows printed pointing to Ind C' -> the arrow Psi = Ind psi apparently intended from Ind C' to Ind C, psi goes from C' to C
\]
On aura encore la formule d'adjonction pour $X = (X_i)_{i\in I}$ dans $\Ind C$, $X' = (X'_{i'})_{i'\in I'}$ dans $\Ind C'$, car
\[
\begin{aligned}
\Hom(X,\Psi X') &= \varprojlim_i \varinjlim_{i'} \Hom(X_i,\psi(X'_{i'}))\\
&= \varprojlim_i \varinjlim_{i'} \Hom(\varphi X_i,X'_{i'})\\
&= \Hom(\Phi X,X').
\end{aligned}
\]

\noindent\msp{page 41}%
On prouve de même que $\Psi$ est pleinement fidèle, car si $X' = (X'_i)_{i\in I}$, $Y' = (Y'_j)_{j\in J}$ sont des objets de $\Ind C'$, on aura
\[
\begin{tikzcd}[column sep=small]
\Hom(X',Y') \arrow[d,"\Psi_{X',Y'}"] & \simeq \varprojlim\limits_i \varinjlim\limits_j & \Hom_{C'}(X'_i,Y'_j) \arrow[d,"\psi_{X'_i,Y'_j}"] \\
\Hom(\Psi X',\Psi Y') & \simeq \varprojlim\limits_i \varinjlim\limits_j & \Hom_C(\psi(X'_i),\psi(Y'_j)),
\end{tikzcd}
\]
donc $\Psi_{X',Y'}$ est un isomorphisme si les $\psi_{X'_i,Y'_j}$ le sont.

Donc la proposition 1 peut s'appliquer à la situation $\Phi$, $\Psi$ entre ind-objets. Je vais l'expliciter sous la forme la plus commode :

\textbf{Proposition 2.} \textit{Soit $K = (K_i)_{i\in I}$ un ind-objet de $C$, et soit $H$ le foncteur}
\[
\left\{
\begin{aligned}
H : X \quad &\longmapsto \quad \underbrace{\Hom_C(X,K)}_{K(X)} \overset{\text{déf}}{=} \varinjlim_i \Hom_C(X,K_i)\\
C^o \quad &\longrightarrow \quad \Ens
\end{aligned}
\right.
\]
\textit{qu'il ind-représente. Conditions équivalentes}

\begin{enumerate}[label={}]
\item[(i)] \textit{Le foncteur $H$ se factorise (à isomorphisme près) par $C'$ via $\varphi$, i.e. transforme les $u \in W_\varphi$ en isomorphismes.}
\printedpage{27}
% head: A. Grothendieck : Dérivateurs, Ch. V, Édition des Universités de Montpellier II et Paris VII
% foot: Version 9 septembre 2001
\item[(i$'$)] \textit{Même condition pour les flèches d'adjonction $X \longrightarrow \psi\varphi(X)$ (lesquelles sont dans $W$).}
\item[(i$''$)] \textit{Même condition pour la flèche d'adjonction $K \longrightarrow \psi\varphi(K)$.}
\item[(ii)] \textit{$K$ appartient à l'image essentielle de} $\Ind C'$ \textit{par $\Psi$.}
\item[(ii$'$)] $K \longrightarrow \Psi\Phi(K)$ \textit{est un} isomorphisme \textit{dans} $\Ind C$.) % [sic] ".)" -> ".", closing parenthesis without opening parenthesis
\item[(ii$''$)] \textit{L'homomorphisme fonctoriel}
\[
K(X) = \Hom_C(X,K) \longrightarrow \Hom_{C'}(\varphi(X),\varphi(K))
\]
\texttt{[plutôt} $\Hom_{\mathrm{Ind}(C)}(X,K) \longrightarrow \Hom_{\mathrm{Ind}(C')}(\Phi(X),\Phi(K))$\texttt{, en identifiant l'objet} $X$ \texttt{de} $C$ \texttt{à son image dans} $\mathrm{Ind}(C)$\texttt{]} \textit{est un isomorphisme.}
\end{enumerate}

\textsc{Démonstration.} L'équivalence de (ii), (ii$'$), (ii$''$) est la proposition 1 appliquée à $\Phi$ et $\Psi$, à cela près qu'il faut noter que (ii$''$) reste évidemment valable si on y remplace $X$ par un ind-objet $(X_j)_{j\in J}$ quelconque. On a

\noindent\msp{page 42}%
\[
\text{(ii$''$)} \Longrightarrow \text{(i)} \Longrightarrow \text{(i$'$)}
\]
tautologiquement, il reste à prouver que (i$'$) implique (i$''$) et (i$''$) (ii$'$). On ne peut pas appliquer directement la proposition 1, car l'hypothèse (i$'$) ici ne porte que sur les objets de $\Ind C$ qui proviennent de $C$, et il n'est pas évident à priori qu'elle passe aux objets quelconques. Mais cela l'est dans le cas d'un homomorphisme de ind-objets $(X_i)$, $(Y_i)$ sur le même ensemble d'indices, donné par un homomorphisme ordinaire (terme à terme) de systèmes inductifs sur $I$, $X_i \longrightarrow Y_i$: si les $X_i \longrightarrow Y_i$ sont dans $W$, l'hypothèse (i) implique que $H(Y) \longrightarrow H(X)$ est un isomorphisme, et de même pour (i$'$) pour $\underbrace{X}_{=(X_i)_i} \;\longrightarrow\; \underbrace{\Psi\Phi(X)}_{=(\psi\varphi(X_i))_i}$\,, où l'homomorphisme d'adjonction est donné terme à terme par les homomorphismes d'adjonction $X_i \longrightarrow \psi\varphi(X_i)$. Donc l'hypothèse (i$'$) ici implique l'hypothèse (i$'$) de la proposition 1 pour $\Phi$ et $\Psi$, donc aussi (ii$'$) (par la proposition 1), et on gagne. % [sic] "et (i'') (ii')" -> "et (i'') => (ii')", the implication sign or verb is missing

\dgrblock{dgrb183}{Why the argument passes to ind-objects}
\begin{dcomment}{exposition}{p42:de la proposition 1 pour}{Why the argument passes to ind-objects}
Grothendieck is careful at the one point where the extension to ind-objects is not a formal substitution. The hypothesis tests units on objects of $C$, whereas Proposition~1, applied to $\operatorname{Ind}(C)$, would test units on all ind-objects. The missing bridge is a limit, and it needs no exactness assumption.

Put $\Phi=\operatorname{Ind}(\varphi)$ and $\Psi=\operatorname{Ind}(\psi)$. For filtered systems $X=(X_i)$ and $Y'=(Y'_j)$,
\[
 \operatorname{Hom}(X,\Psi Y')
  =\varprojlim_i\varinjlim_j\operatorname{Hom}_C(X_i,\psi Y'_j)
  \cong\varprojlim_i\varinjlim_j
        \operatorname{Hom}_{C'}(\varphi X_i,Y'_j).
\]
This proves $\Phi\dashv\Psi$. Applying the same formula to two systems from $C'$ proves that $\Psi$ is fully faithful. The upper arrow in the edition's displayed pair of ind-functors points the wrong way; these types determine its direction.

For a fixed ind-object $K$, write $\widehat H(X)=\operatorname{Hom}_{\operatorname{Ind}(C)}(X,K)$. Its restriction to $C$ is the functor $H$ of the proposition. Suppose each unit $X_i\to\psi\varphi X_i$ induces a bijection under $H$. The unit of an arbitrary system $X$ is represented term by term by these maps. Therefore
\[
 \widehat H(\Psi\Phi X)=\varprojlim_i H(\psi\varphi X_i)
   \xrightarrow{\sim}\varprojlim_i H(X_i)=\widehat H(X).
\]
A limit of a natural diagram of bijections is a bijection: its inverse is obtained by taking the componentwise inverses, which remain compatible. There is no interchange of a filtered colimit with an inverse limit here. Proposition~1 now applies in $\operatorname{Ind}(C)$ and gives $K\cong\Psi\Phi K$.

Condition (i$''$) needs this extension $\widehat H$ to have a meaning, because $K$ is not in general an object of $C$. The text uses $\psi\varphi K$ as shorthand for $\Psi\Phi K$; the corrected statement is a test at the unit of an ind-object. Likewise condition (ii$''$), when tested only on $X\in C$, extends to arbitrary ind-objects by the displayed limit formula. These are not extra hypotheses or repairs to the argument. They make explicit the conventions under which its proof is complete.
\end{dcomment}


Revenons à la situation où $C = \Cathot$, $C' = \Hot$, $\varphi$ le foncteur \texttt{[de]} localisation. Si $K \in \Ob\Cathot$, alors $\psi\varphi(K)$ s'appelle \textit{l'enveloppe de Thomason-Kan} de $K$ - si $K$ est une catégorie, son enveloppe $\bar{K}$ et la flèche $K \longrightarrow \bar{K}$ ne sont définis que dans la catégorie Cat\,hot, i.e. $\bar{K}$ n'est défini qu'à `homotopisme près'. % reading resolved: no underline is printed under psi phi(K).

\noindent\msp{page 43}%
Mais on peut toujours \textit{choisir} un objet $\bar{K}$ de Cat, et une flèche
\[
i : K \longrightarrow \bar{K} \qquad\qquad\qquad \textit{dans } \Cat,
\]
\printedpage{28}
% head: A. Grothendieck : Dérivateurs, Ch. V, Édition des Universités de Montpellier II et Paris VII
% foot: Version 9 septembre 2001
pas seulement dans Cat\,hot, de telle façon de plus que $i$ soit une `\textit{cofibration} de Thomason' (et en particulier une immersion ouverte?). Il en résulte que si on a un système inductif
\[
K = (K_n)_{n\in\mathbf{N}}
\]
indexé par $\mathbf{N}$, et si on choisit pour chaque $n$ une enveloppe de Thomason-Kan
\[
K_n \overset{\alpha_n}{\hookrightarrow} \bar{K}_n
\]
de cette façon, on peut trouver des foncteurs de transition
\[
\bar{K}_n \xrightarrow{\bar{i}_n} \bar{K}_{n+1}
\]
tels que les carrés
\[
\begin{tikzcd}
K_n \arrow[r,hook,"\alpha_n"] \arrow[d,"i_n"'] & \bar{K}_n \arrow[d,"\bar{i}_n"] \\
K_{n+1} \arrow[r,hook,"\alpha_{n+1}"] & \bar{K}_{n+1}
\end{tikzcd}
\]
soient commutatifs, de telle sorte que l'on ait donc un homomorphisme de \textit{systèmes inductifs de} Cat
\[
K \longrightarrow \bar{K}
\]
et a fortiori un homomorphisme de ind-objets de Cat. Ceci posé, les conditions (i), (ii$'$) de la proposition 2 montrent que le foncteur ind-représenté par $K$ \textit{dans} $C = \Cathot$ transforme quasi-isomorphismes

\dgrblock{dgrb184}{What Thomason replacement supplies}
\begin{dcomment}{later}{p43:de cette façon, on peut trouver des foncteurs de transition}{What Thomason replacement supplies}
Grothendieck needs replacements which can be chosen together with the transition maps, so that he still has a sequence of actual categories. This part of the argument has a direct model-category realization. Given a sequence $K_n\to K_{n+1}$, choose trivial cofibrations $\alpha_n:K_n\to R_n$ with $R_n$ Thomason fibrant. The square with top map $\alpha_n$, bottom object $R_{n+1}$ and right map $R_{n+1}\to *$ admits a lift. It gives $R_n\to R_{n+1}$ extending $\alpha_{n+1}i_n$. Since the indexing category is the freely generated chain, no relations beyond these adjacent squares need to be imposed.

The parenthetical question about open immersions also has a precise answer. Thomason cofibrations are sieves, hence full inclusions with no arrows from outside to inside. A technical correction matters here: Thomason originally asserted that retracts of Dwyer maps remain Dwyer maps. \citet{Cisinski1999Dwyer} gives a counterexample and proves that every cofibration is instead a pseudo-Dwyer map; the corrected assertion retains the sieve property and the proper model structure. This repairs that specific later proof, not a gap in Grothendieck's formal Propositions~1--2.

One can also replace any sequence of fibrant objects by a sequence of fibrant objects with cofibration transitions, up to a levelwise weak equivalence. Start with $T_0=R_0$. If $g_n:T_n\to R_n$ has been constructed, factor the composite $T_n\to R_n\to R_{n+1}$ as
\[
 T_n\lhook\joinrel\longrightarrow T_{n+1}
       \xrightarrow[\sim]{g_{n+1}}R_{n+1},
\]
where the first map is a cofibration and the second a trivial fibration. Then $T_{n+1}$ is fibrant and the $g_n$ form a strict map of sequences. The transition inclusions are fully faithful. This proves a useful model-theoretic version of the rearrangement proposed on manuscript p.~46. It gives a levelwise weak equivalence and therefore an isomorphism of ind-objects in $\Hot$. Calling it an isomorphism in $\operatorname{Ind}(\Cathot)$ requires a further comparison; that conclusion is not supplied by the factorization axiom.
\end{dcomment}


\noindent\msp{page 44}%
ordinaires (équivalences faibles) en bijections, \textit{si et seulement si} l'homomorphisme précédent est un isomorphisme de ind-objets de Ind\,Cat\,hot.

Ce sera donc cela la condition naturelle à imposer aux ind-catégories de chemins
\[
\Chh_{\Sigma}(\mathcal{M};a,b) \qquad\qquad\qquad (a,b \in \mathrm{ob}\,\mathcal{M})
\]
provenant d'une `catégorie de modèles' $\mathcal{M}$, $\Sigma$.

Bien sûr, disposant d'un système inductif explicite de catégories (pas seulement d'objets de Cat\,hot)
\[
K_n = \Ch_{\Sigma,n}(\mathcal{M};a,b),
\]
il est tentant de passer à la limite inductive
\[
K_\infty = \varinjlim_n \Ch_{\Sigma,n}(\mathcal{M};a,b) \overset{\text{déf}}{=} \Ch_{\Sigma}(\mathcal{M};a,b),
\]
et ceci d'autant plus que les foncteurs de transition sont des monomorphismes et pleinement fidèles, de sorte que les $K_n$ s'identifient à une suite de sous-catégories pleines de

% ---- from Der-ChV-234_p031-034.tex
\printedpage{29}
% head: A. Grothendieck : Dérivateurs, Ch. V, Édition des Universités de Montpellier II et Paris VII
% foot: Version 9 septembre 2001
$K_\infty$, dont la réunion est $K_\infty$. Mais d'autre part je ne sais donner aucune interprétation convaincante de $K_\infty$, comme de celle du ind-objet de Cat\,hot defini par $\uuK = (K_n)_{n\in\mathbf{N}}$, qui ind-représente le foncteur remarquable % [sic] defini -> défini, missing accent
\[
H = H_{\mathcal{M};x,y} : \Cathot^{o} \longrightarrow \Ens
\]
factorisant le foncteur
\msp{page 45}
\[
X \longmapsto H_{\mathcal{M};x,y}(X) \defeq \Hom_{\mathcal{M}(X)\Sigma_X^{-1}}(x_X, y_X),
\]
par la formule donc
\[
H(X) = \varinjlim_{n} \underbrace{\pi_0\uHom(X, K_n)}_{\Hom_{\Cathot}(X,K_n)} .
\]
Je ne sais donner d'interprétation de ce type ni de la catégorie $K_\infty$ elle-même, ni de son image dans Cat\,hot, ni de son image dans Hot. Il est vrai que si je restreins le foncteur $H$ à la sous-catégorie pleine de (Cat) formée des catégories \textit{finies} (Ob\,$X$ et Fl\,$X$ finis), alors on aura
\[
\uHom^{!}(X, K_\infty) = \uHom(X, K_\infty),
\]
et l'isomorphisme précédant s'interprète comme % [sic] précédant -> précédent, adjective intended
\[
\Hom_{\mathcal{M}(X)\Sigma_X^{-1}}(a_X, b_Y) \defeq H(X) \simeq \pi_0\uHom(X, K_\infty) \defeq \Hom_{\Cathot}(X, K_\infty).
\]
\texttt{[plutôt} $(a_X, b_X)$\texttt{]} Cela prouve du moins que la connaissance de $K_\infty$ est utile. Mais cette formule, pour $X$ restreint aux catégories \textit{finies,} ne définira pas $K_\infty$ comme objet de Cat\,hot, à isomorphisme unique près dans cette catégorie - \textit{sauf} dans le cas où $K_\infty$ est elle-même homotope à une catégorie
\msp{page 46}
finie. (Je ne me rends pas bien compte si ça arrive très souvent \ldots\ Mais on ne peut fonder une théorie générale des dérivateurs et des catégories de modèles `dérivables' sur une hypothèse de ce genre.)

\dgrblock{dgrb185}{Finitely many objects already give a uniform stage}
\begin{dcomment}{development}{p45:formée des catégories}{Finitely many objects already give a uniform stage}
Grothendieck uses finite categories to remove the exclamation mark. The proof is controlled by the number of objects, because fullness has already supplied every arrow between objects that have appeared. We record the resulting refinement to make this source of finiteness explicit.

\textbf{Finite-object factorization.}\ours\quad
Let $K_0\subseteq K_1\subseteq\cdots$ be full subcategories, let $K_\infty$ be their union, and let $X$ have finitely many objects, with no restriction on its sets of arrows. Then the canonical functor
\[
 \varinjlim_n\operatorname{Fun}(X,K_n)
       \longrightarrow\operatorname{Fun}(X,K_\infty)
\]
is an isomorphism of categories. In particular, the path-localization formula on manuscript p.~45 holds for such $X$.

\textit{Proof.} Given $F:X\to K_\infty$, choose a stage containing each of its finitely many object-values, and take the largest of these stages. Every arrow of $X$ has endpoints whose images lie in that stage, so fullness supplies its image arrow there. The composition and identity equations hold because inclusion is faithful. Thus $F$ factors through this one stage. If $F,G$ are two functors represented at possibly different stages, put them in a common stage. All components of every transformation $F\Rightarrow G$ then belong to that stage by fullness; its naturality equations are unchanged. Conversely, inclusions of stages introduce no identifications between different transformations. This proves both bijectivity on objects and bijectivity on every Hom-set.\hfill$\square$

There is a version invariant under equivalence of the source. If $X$ has only finitely many isomorphism classes of objects, choose a finite-object skeleton $S$ and an equivalence $S\to X$. Precomposition gives equivalences of functor categories into every $K_n$ and into $K_\infty$. Passing to $\pi_0$ therefore proves
\[
 \varinjlim_n\pi_0\operatorname{Fun}(X,K_n)
       \cong\pi_0\operatorname{Fun}(X,K_\infty).
\]
The strict equality of bounded and unrestricted functor categories need not hold for this larger class: an actual functor may send infinitely many isomorphic source objects to different target objects appearing at unbounded stages. The statement on components is the appropriate invariant one.

The qualification at the end of manuscript p.~45 also needs its domain of comparison. Ordinary Yoneda gives uniqueness among objects of $\Cathot$ homotopy equivalent to finite categories. It does not give uniqueness among all objects merely because one candidate is finite. The same elementary fence which will make the later length obstruction visible already separates these assertions.

\textbf{Failure of unrestricted uniqueness from finite tests.}\ours\quad
Let $B$ be the poset $0\to1\leftarrow2\to3\leftarrow\cdots$, and let $B_n$ be its full subcategory on $\{0,\ldots,n\}$. The restrictions of $\operatorname{Hom}_{\Cathot}(-,B)$ and $\operatorname{Hom}_{\Cathot}(-,*)$ to finite categories are naturally isomorphic, although $B$ and $*$ are not isomorphic in $\Cathot$.

\textit{Proof.} The maps $r_j(i)=\min(i,j)$, for $0\leq j\leq n$, are endofunctors of $B_n$: they preserve every fence arrow below $j$ and collapse every arrow above it to an identity. For even $j$ there is a transformation $r_j\Rightarrow r_{j+1}$; for odd $j$ there is a transformation in the reverse direction. Thus $r_0$, the constant functor at $0$, and $r_n=1_{B_n}$ are joined by a zigzag. Postcomposition contracts $\operatorname{Fun}(X,B_n)$ for every $X$. When $X$ is finite, the preceding factorization lemma places each functor $X\to B$ in some $B_n$. It is therefore homotopic to the constant functor at $0$, and $\pi_0\operatorname{Fun}(X,B)$ is a singleton. These singleton identifications are natural in $X$.

If $1_B$ were joined to the constant functor at $0$ by $L$ transformations, evaluation at $L+1$ would give a zigzag of $L$ arrows connecting $L+1$ to $0$. Each fence arrow changes the index by at most one, which is impossible. Hence $B$ is not categorically contractible.\hfill$\square$

The nerve of $B$ realizes to an infinite ray and is contractible. This counterexample concerns $\Cathot$; the last-page obstruction in $\Hot$ will require a different example. It gives a precise reason to restrict the finite-object exception to comparisons within the finite subcategory, if that is the intended reading of the parenthesis.

We checked the general filtered-colimit and functor-category framework in \citet[Chapters~1 and 7]{Riehl2014} and the specific argument in the chapter. We did not find this exact refinement stated there. It is an elementary continuation of the printed proof, with no claim of priority. Its relevance is practical: a one-object category with infinitely many endomorphisms still satisfies the lemma, even though it is not a finite category in the author's stated sense.
\end{dcomment}


Mais supposons que le foncteur $H$ envisagé se factorise même par Hot, i.e. qu'il transforme quasi-isomorphismes en bijections. On a vu que cela équivaut à la condition qu'il puisse se représenter par un système inductif sur $\mathbf{N}$ formé de catégories de Thomason-Kan. Il doit être facile de plus, si on y tient, de s'arranger pour que les foncteurs de transition soient des monomorphismes et pleinement fidèles, pour avoir une $\varinjlim$ sympathique. Ce sera bien sûr une catégorie de Thomason-Kan. Cette catégorie est-elle déterminée par le foncteur $H : \Cathot^{o} \longrightarrow \Ens$ à isomorphisme unique près dans Cat\,hot? Bien sûr, il n'est pas question qu'elle représente le foncteur $H$, sauf \texttt{[si]} celui-ci est bel et bien représentable, i.e. si le ind-objet $\uuK$ de Cat\,hot est essentiellement constant (ce qui ne doit pratiquement jamais être le cas). Notons déjà
\printedpage{30}
% head: A. Grothendieck : Dérivateurs, Ch. V, Édition des Universités de Montpellier II et Paris VII
% foot: Version 9 septembre 2001
\msp{page 47}
que si $u$ est un homomorphimse de ind-objets \textit{dans} Cat, ou, ce qui revient maintenant au même, un homomorphisme dans Cat % [sic] homomorphimse -> homomorphisme, transposed letters
\[
K_\infty \xrightarrow{u_\infty} K'_\infty
\]
tel que $\forall n$, $u_\infty(K_n) \subseteq K'_{n'}$ pour $n' = \tau(n)$ convenable, et si on suppose que l'homomorphisme de ind-objets induit un isomorphisme pour les foncteurs ind-représentables \textit{sur} Cat\,hot (par exemple, si $\uuK$ et $\uuK'$ ind-représentent le même foncteur $H$, et l'homomorphisme de ind-objets est celui qui correspond à $\mathrm{id}_H$), alors $u_\infty$ est forcément un isomorphisme \textit{dans} Hot, donc aussi (sauf erreur) dans Cat\,hot puisque $K_\infty$, $K'_\infty$ sont des catégories de Thomason-Kan. Mais ce qui n'est pas évident, c'est l'existence d'un homomorphisme $u_\infty$ comme \texttt{[ci-]}dessus, et s'il existe, qu'il soit unique (ou du moins canonique?) comme flèche de Cat\,hot.

\dgrblock{dgrb186}{Fibrancy passes to the union}
\begin{dcomment}{exposition}{p46:Ce sera bien sûr une catégorie de Thomason-Kan.}{Fibrancy passes to the union}
Grothendieck expects the union of his replacement categories to remain a good representative of a homotopy type. For Thomason fibrancy this expectation is correct and has a finite combinatorial proof. A category $K$ is Thomason fibrant when $\operatorname{Ex}^2NK$ is a Kan complex. For a filtered system $(K_i)$, the nerve commutes with its colimit: an $n$-simplex is a finite composable string, and both the string and the finitely many equations expressing a morphism of such strings occur at one stage. Moreover
\[
 (\operatorname{Ex}^2 S)_n
   =\operatorname{Hom}_{\mathbf{sSet}}(\operatorname{Sd}^2\Delta^n,S)
\]
commutes with filtered colimits, since $\operatorname{Sd}^2\Delta^n$ is a finite simplicial set. Finally, a horn in a filtered colimit of Kan complexes lies in one stage, and has a filler at that stage. Consequently
\[
 \operatorname{Ex}^2N(\varinjlim_i K_i)
       \cong\varinjlim_i\operatorname{Ex}^2NK_i
\]
is Kan whenever all the terms are Kan. This proof even dispenses with fully faithful transition maps.

Fibrancy does not answer the next uniqueness question. There are three different conclusions available from three different kinds of information. A presheaf determines its representing ind-object in $\operatorname{Ind}(\Cathot)$. A coherent map of strict diagrams which is a weak equivalence at every stage gives a weak equivalence on the appropriate homotopy colimits. A specified categorical homotopy between maps of unions requires actual natural transformations, with a finite common length. These conclusions are not interchangeable.

In particular, an isomorphism in the ind-category of a homotopy category is only eventual equality of homotopy classes at each stage. It does not come with compatible choices of the homotopies. The claim on manuscript p.~47 that a corresponding union map must be a weak equivalence is false, even for full inclusions of finite fibrant stages with fibrant union; the example is in the commentary at ms.~49. The verified model-category replacement construction in the preceding block avoids this difficulty by giving actual coherent maps and levelwise weak equivalences. The raw categorical-homotopy strictification questions remain separately recorded in the audit.
\end{dcomment}


Pour l'existence d'un $u_\infty$, on est ramené à prouver que si on a un diagramme à flèches pleines
\[
\begin{tikzcd}
K_n \arrow[r, "i_n", "\text{cofib.}"'] \arrow[d, "u_n"] & K_{n+1} \arrow[dr, dashed] & \\
K'_{\tau(n)} \arrow[rr, "\text{transition}"'] & & K'_N ,
\end{tikzcd}
\]
et si une flèche $u_{n+1}$ est donnée `en pointillés' \textit{dans} Cat\,hot, de façon à rendre le carré commutatif,
\msp{page 48}
alors on peut la réaliser par une flèche \textit{dans} Cat qui rende le diagramme commutatif pour de bon (dans Cat). On est ramené au même problème dans une situation
\[
\begin{tikzcd}
X \arrow[r, hook, shift right, "i"'] & Y , \arrow[l, dashed, shift right, "p"']
\end{tikzcd}
\]
où $X$ est une catégorie de Thomason, $i$ une cofibration de Thomason, $p$ une rétraction de $Y$ dans $X$ \textit{dans} Cat\,hot, on veut la réaliser par une vraie rétraction (\footnote[7]{\texttt{[un ?\ \ dans la marge]}}).

En admettant que ça soit le cas, donc que tout homomorphisme des ind-objets $\uuK$, $\uuK'$ \textit{dans} Cat\,hot se représente par un homomorphisme entre les ind-objets \textit{dans} Cat, donc aussi par un foncteur $K_\infty \longrightarrow K'_\infty$, reste la question si l'homomorphisme correspondant dans Cat\,hot est unique. Donc si on a deux foncteurs $f$, $g$ qui induisent le même homomorphisme sur les foncteurs ind-représentés sur Cat\,hot, si $f$ et $g$ sont homotopes. Donc

\textbf{Question}. Deux foncteurs $f, g : K_\infty \longrightarrow K'_\infty$ tels que pour tout $n \in \mathbf{N}$, $f_n$, $g_n$ qui s'en déduisent par restriction à $K_n$ sont homotopes (et même homotopes en tant que foncteurs de $K_n$ dans un $K'_N$ convenable), sont-ils homotopes? (\footnote[8]{$K$, $K'$ catégories de Thomason, les $K_n$, $K'_n$ aussi.})

\printedpage{31}
% head: A. Grothendieck : Dérivateurs, Ch. V, Édition des Universités de Montpellier II et Paris VII
% foot: Version 9 septembre 2001
\msp{page 49}
Je présume que non - qu'il peut arriver que la longueur minimum requise $l(n)$ pour un chemin dans $\uHom(K_n, K'_{\tau(n)})$ reliant $f_n$ et $g_n$, tende vers l'infini si $n \longrightarrow \infty$. L'hypothèse que les $K'_n$ sont des catégories de Thomason y change-t-il quelque chose (deux objets homotopes dans $\uHom(X, K')$, si $K'$ est de Thomason, sont-ils reliés par une flèche?)

\dgrblock{dgrb187}{A complete instance of the length obstruction}
\begin{dcomment}{development}{p49:Je présume que non - qu'il peut arriver}{A complete instance of the length obstruction}
Grothendieck identifies a concrete mechanism by which stagewise homotopies could fail to assemble: their lengths may become unbounded. The following example realizes that mechanism with full inclusions at every stage. It deliberately leaves the additional target-fibrancy condition aside, so that the obstruction itself can be seen without a model-category comparison.

\textbf{Unbounded categorical homotopies.}\ours\quad
There are full filtrations $A_0\subset A_1\subset\cdots$ and $B_0\subset B_1\subset\cdots$, and functors $f,g:A_\infty\to B_\infty$, such that $f|_{A_n}$ and $g|_{A_n}$ are homotopic in $\operatorname{Fun}(A_n,B_n)$ for every $n$, but $f$ and $g$ are not homotopic in $\operatorname{Fun}(A_\infty,B_\infty)$.

\textit{Proof.} Take $A_n$ to be discrete on $\{0,\ldots,n\}$. Take $B_n$ to be the full subcategory on the same vertices of the infinite alternating fence
\[
 0\longrightarrow1\longleftarrow2\longrightarrow3
       \longleftarrow4\longrightarrow\cdots.
\]
There are no nonidentity composites, so this is a poset. Define $f(i)=0$ and $g(i)=i$. For $0\leq j\leq n$, put $h_j(i)=\min(i,j)$ on $A_n$. Since the source is discrete, these assignments are functors. If $j$ is even there is a transformation $h_j\Rightarrow h_{j+1}$, and if $j$ is odd there is a transformation $h_{j+1}\Rightarrow h_j$: each component is either the identity or the indicated fence arrow. Hence
\[
 h_0=f|_{A_n},\quad h_1,\quad\ldots,\quad h_n=g|_{A_n}
\]
is a categorical homotopy of length $n$, contained in $B_n$.

Suppose a zigzag of $L$ natural transformations connected $f$ and $g$ globally. Evaluate it at the object $L+1$ of $A_\infty$. Each nonidentity arrow in $B_\infty$ changes the vertex by exactly one, so a zigzag of $L$ arrows cannot connect $0$ to $L+1$. This contradiction proves the assertion; it also proves that the displayed stagewise length $n$ is minimal.\hfill$\square$

Both filtrations consist of full subcategories, and the source stages are discrete, hence Thomason fibrant. We have not shown that the target stages or their union are Thomason fibrant; this construction therefore does not answer the stronger question with that additional hypothesis. The next commentary does. It proves exactly the length mechanism suggested in the first sentence of manuscript p.~49. The argument was developed directly from that sentence. We checked the surrounding categorical-homotopy and model-category references cited in this section and did not find this exact fence formulation; it is presented as an elementary illustration, not a priority claim.
\end{dcomment}


\dgrblock{dgrb188}{The same ind-morphism, two different maps of the unions}
\begin{dcomment}{development}{p49:Je présume que non - qu'il peut arriver}{The same ind-morphism, two different maps of the unions}
With Thomason fibrancy imposed on every stage and on the unions, the questions of ms.~47--49 still have negative answers. The example is a strict form of the symmetric-group shift used at ms.~52.

\textbf{Proposition.}\ours
\textit{There are a sequence $T_0\subset T_1\subset\cdots$ of full inclusions of finite categories, each $T_n$ Thomason fibrant and isomorphic in $\mathrm{Cat}\,\mathrm{hot}$ to a group, with Thomason-fibrant union $T$, and a functor $s:T\to T$ with $s(T_n)\subseteq T_{n+1}$, such that:}
\begin{enumerate}[label=(\alph*)]
\item \textit{for every $n$, the functor $s|_{T_n}$ and the inclusion $T_n\to T_{n+1}$ are joined by a zigzag of three natural transformations, so $s$ and $\mathrm{id}_T$ induce the same endomorphism of the ind-object $(T_n)$ of $\mathrm{Cat}\,\mathrm{hot}$;}
\item \textit{$Bs$ is not a weak equivalence; in particular $s$ and $\mathrm{id}_T$ are not homotopic;}
\item \textit{there is a weak equivalence $r:T\to\mathsf BG$ onto a group, between Thomason-fibrant categories, that is not an isomorphism in $\mathrm{Cat}\,\mathrm{hot}$.}
\end{enumerate}

\textit{Proof.}
Let $G_j$ be the group of permutations of $\{0,\dots,j\}$, with $G_j\subset G_{j+1}$ by fixing $j+1$, and $G=\bigcup G_j$. Let $T_n$ have the objects $0,\dots,n$, with $\operatorname{Hom}(i,j)=G_j$ if $i\leq j$ and empty otherwise, and composition by multiplication in the group of the target. The inclusions $T_n\subset T_{n+1}$ are full. The functor $r_n:T_n\to\mathsf BG_n$ that forgets the objects and the functor $j_n:\mathsf BG_n\to T_n$ onto the object $n$ satisfy $r_nj_n=\mathrm{id}$, and the arrows $i\to n$ labelled by the unit form a natural transformation $\mathrm{id}\Rightarrow j_nr_n$. So $T_n\cong\mathsf BG_n$ in $\mathrm{Cat}\,\mathrm{hot}$.

Fibrancy. A category is Thomason fibrant when it has the right lifting property for the functors $c\,\mathrm{Sd}^2\Lambda\to c\,\mathrm{Sd}^2\Delta$ associated with the horn inclusions, and these are sieves \citep[Proposition~4.2]{Thomason1980}. Let $F:A\to T_n$ be given on such a sieve $A\subset B$. The functor $r_nF$ extends to $B$, because the nerve of a groupoid is a Kan complex. Extend $F$ by sending every object of $B$ outside $A$ to $n$ and every arrow that meets the complement of $A$ to the arrow with the label given by the extension of $r_nF$. No arrow leads from the complement into $A$, and $\operatorname{Hom}(i,n)=G_n$ for all $i$, so this is defined, and it is a functor because composition in $T_n$ multiplies labels. So $T_n$ is fibrant, and the union $T$ is fibrant by the preceding commentary on ms.~46.

The shift. For $g\in G_j$ let $\sigma(g)\in G_{j+1}$ fix $0$ and act on $\{1,\dots,j+1\}$ as $g$ does on $\{0,\dots,j\}$; let $s:T\to T$ send $i$ to $i+1$ and the arrow labelled $g$ to the arrow labelled $\sigma(g)$. On $T_n$ both $s$ and the inclusion into $T_{n+1}$ map, by the arrows labelled by the unit, to the functors with the same labels and constant value $n+1$ on objects, and these two are isomorphic by the cycle $c=(0\,1\,\cdots\,n+1)$, since $cgc^{-1}=\sigma(g)$ for $g\in G_n$. This is (a).

The functors $r_n$ are compatible and give $r:T\to\mathsf BG$; the nerve and homotopy groups commute with the union, so $Br$ is a weak equivalence and $BT$ is a $K(G,1)$. Under this identification $s$ induces $\sigma:G\to G$ on $\pi_1$, which is not surjective (its image fixes $0$). This is (b). For (c): if $r$ were invertible in $\mathrm{Cat}\,\mathrm{hot}$, there would be a functor $j:\mathsf BG\to T$ with $rj$ joined to the identity by natural transformations, that is, with $rj:G\to G$ an inner automorphism. But $j$ sends the object to some $i$ and $G$ into $\operatorname{Hom}(i,i)=G_i$, so the image of $rj$ lies in $G_i\neq G$.\hfill$\square$

Consequences for these pages. The assertion of ms.~47, that a morphism of such systems inducing an isomorphism of ind-objects induces a weak equivalence of the unions, is false: $s$ induces the identity of the ind-object. The question of ms.~48, whether two maps of unions that are homotopic on every stage inside a later stage are homotopic, has the answer no, with homotopies of length three at every stage; what fails is their compatibility, not the growth of their length as in the fence example above. The union of fibrant stages is fibrant, but the union $T$ of stages each isomorphic in $\mathrm{Cat}\,\mathrm{hot}$ to a groupoid is not isomorphic to one, by (c). And a weak equivalence between Thomason-fibrant categories need not be an isomorphism in $\mathrm{Cat}\,\mathrm{hot}$ (ms.~47, ``sauf erreur''): fibrancy is not the envelope property.
\end{dcomment}


\dgrblock{dgrb189}{Fibrancy does not reduce every homotopy to one arrow}
\begin{dcomment}{examples}{p49:y change-t-il quelque chose}{Fibrancy does not reduce every homotopy to one arrow}
The parenthetical question asks a sharper thing than a uniform bound: for a Thomason fibrant target, must two homotopic functors already be joined by a single natural transformation? A three-object category gives a negative answer, and the fibrancy can be checked directly against the generating horns.

Every category $C$ with a terminal object $t$ is Thomason fibrant. Indeed, \citet[Propositions~2.5 and 4.2]{Thomason1980} identify the lifting test with inclusions
\[
 c\operatorname{Sd}^2\Lambda^k[n]
     \lhook\joinrel\longrightarrow c\operatorname{Sd}^2\Delta[n],
\]
which are sieves. More generally, if $A\subseteq D$ is a sieve and $F:A\to C$ is a functor, extend $F$ by assigning $t$ to every object outside $A$. There are no arrows from an outside object to an object of $A$. Send arrows between outside objects to $1_t$, and each arrow from inside to outside to the unique map into $t$. Composition is respected by uniqueness of maps to the terminal object. This constructs the required lift against every displayed horn, so $C\to *$ is a Thomason fibration.

Now take $C$ to be $0\to1\leftarrow2$, whose terminal object is $1$, and take $X=*$. The functors selecting $0$ and $2$ lie in the same component of $\operatorname{Fun}(*,C)=C$, connected by the two-arrow zigzag through $1$. There is no arrow $0\to2$ or $2\to0$, so there is no single transformation in either direction between them. The target is fibrant by the preceding proof.

This answers the exact one-arrow question. It does not say that minimal lengths can grow without bound for a fixed fibrant target, and it does not establish a counterexample to the preceding question with fibrant full filtrations. A claim of bounded shortening, or of compatibility of stagewise homotopies, has different quantifiers and needs a separate argument.
\end{dcomment}


Bien sûr, on peut prétendre que les catégories $\uCh_{\Sigma,n}(\mathcal{M}; a, b)$ ont été définies `canoniquement', et leur limite inductive $\uCh_{\Sigma}(\mathcal{M}; a, b)$ aussi. Mais la définition était assez arbitraire et abracadabrante - je suis convaincu que dix mathématiciens, se proposant de développer cette même idée des catégories de chemins, dans le même but que moi, arriveraient à dix définitions différentes! Pourtant, il n'est pas exclu que l'on finisse par trouver un système transitif d'isomorphismes dans Cat\,hot, entre les catégories de chemins `infinis' (i.e. de longueur non précisée) $K_\infty = K_\infty(\mathcal{M}; a, b)$ que nous aurons construites, dans le but
\msp{page 50}
principal (disons) de représenter le foncteur $H$ \textit{sur les catégories finies.} Et que les isomorphismes obtenus seront, bien entendu, des homomorphismes dans Ind(Cat\,hot). J'ai l'impression de ne pas avoir encore axiomatisé à fond la notion même de chemins, qui après tout a un sens géométrique assez évident et important par lui-même. Donc par ce qu'il faut entendre par `\textit{une} théorie des chemins dans une catégorie'. Une telle théorie $\mathcal{T}$ associerait donc à $\mathcal{M}$, $\Sigma$, $a$, $b$ un système inductif de catégories $\Ch^{\mathcal{T}}_{n,\Sigma}(\mathcal{M}; a, b)$. Et les contraintes géométriques naturelles à imposer à de telles théories, en vue aussi de leur utilisation pour expliciter
\[
\Hom_{\mathcal{M}(X)\Sigma_X^{-1}}(a_X, b_X) \xrightarrow{\ \sim\ } \pi_0\uHom(X, \Chh^{\mathcal{T}}(\mathcal{M}; a, b)) % reading resolved: two underlines, the ind-path object.
\]
donnerait un système transitif d'isomorphismes non seulement entre les ind-objets de Cat\,hot, mais peut-être aussi entre les $\Ch^{\mathcal{T}}_{\infty,\Sigma}(\mathcal{M}; a, b)$, pour des théories differentes. Si on dispose d'un tel énoncé, la catégorie $\Ch_{\infty,\Sigma}(\mathcal{M}; a, b)$, ou plus exactement ses images dans Cat\,hot et Hot, % [sic] differentes -> différentes, missing accent
\msp{page 51}
prendraient un sens convaincant.

\dgrblock{dgrb190}{What a canonical mapping space requires}
\begin{dcomment}{later}{p51:prendraient un sens convaincant}{What a canonical mapping space requires}
Grothendieck wants a theory of paths whose result is independent of its combinatorial presentation. The preceding scholium already supplies such independence at the level it actually computes: different constructions representing the same $H_{a,b}$ give uniquely isomorphic ind-objects in $\Cathot$. To obtain a canonical homotopy type of paths, one needs information beyond the connected-component calculation.

Dwyer and Kan's simplicial localization provides the standard stronger construction. For a relative category $(\mathcal M,W)$ it gives simplicial mapping objects with the localized Hom-sets as their connected components \citep[\S2 and Proposition~3.1]{DwyerKan1980Calculating}. Their hammock construction uses diagrams of zigzags and coherent higher comparisons; in a model category its mapping objects recover the usual derived function complexes \citep[Proposition~4.4 and Corollary~4.7]{DwyerKan1980Function}. Thus composition and its higher compatibility are retained while allowing different presentations of the same homotopy theory.

There is a precise caution in comparing it with these pages. Grothendieck's $K_\infty$ is the ordinary union of full categories of fixed-endpoint paths, and the morphisms of those categories are the transformations specified in sections~1--2. The theorem proved here identifies $\pi_0$ of bounded diagram categories with an ordinary localized Hom-set. That theorem alone is not a weak-equivalence comparison between $NK_\infty$ and a Dwyer--Kan function complex. Replacing vertical morphisms, changing allowed zigzags, or using a homotopy colimit may change the higher homotopy type while leaving this $\pi_0$ calculation intact.

Accordingly, the modern theory answers the demand for an invariant mapping-space object, but does not automatically verify each proposed comparison between the particular unbounded path categories in the manuscript. In the favorable model case discussed above, the desired functor $H_{a,b}$ is represented on $\Hot$ by the derived mapping space. A direct identification of a chosen $K_\infty$ with that object is a separate theorem to be proved under stated assumptions. This distinction explains both the force of the scholium and the legitimacy of the questions which follow it.
\end{dcomment}


Il y a d'autre part un autre point de vue, qui peut-être pourrait `canonifier' les catégories $\uCh_\infty$ en tant qu'objet de Cat\,hot et surtout de Hot (on passerait à Cat\,hot en appliquant le foncteur $\psi$, enveloppe de Thomason-Kan \ldots). En effet, on est intéressé à construire, en termes de $\mathcal{M}$, $\Sigma$, un `dérivateur' (ou un prédérivateur) $\mathbf{D}_{\mathcal{M},\Sigma}$ surtout sur le `domaine' des catégories finies. L'étendre à un domaine plus large est peut-être important dans un `deuxième souffle', mais peut être considéré comme du luxe dans le premier mouvement de fondation d'une théorie des dérivateurs. Ceci dit, on peut se poser des questions du type suivant.

Soit
\[
\Hot_0 \subset \Hot
\]
\printedpage{32}
% head: A. Grothendieck : Dérivateurs, Ch. V, Édition des Universités de Montpellier II et Paris VII
% foot: Version 9 septembre 2001
la sous-catégorie pleine de Hot, définie par les catégories finies (\footnote[9]{\textit{NB} Hot$_0$ est équivalente à une petite catégorie.}) - appelons les types d'homotopie `de type fini' (ou de \textit{présentation} finie?). On a un foncteur évident
\msp{page 52}
\[
\begin{aligned}
\Hot \quad &\longrightarrow \quad \Hot_0^{\wedge} = \uHom(\Hot_0^{o}, (\Ens)) \, ,\\
X \quad &\longmapsto \quad \tilde{X} = (\xi \longmapsto \Hom_{\Hot}(\xi, X)).
\end{aligned}
\]
Ce foncteur est \textit{conservatif} \texttt{[cela semble faux]}, puisque un homomorphisme $X \longrightarrow X'$ dans Hot est un isomorphisme s'il induit des isomorphismes pour les foncteurs $\tilde{X} \longrightarrow \tilde{X}'$, et même seulement pour leurs valeurs pour les sphères $S^n$ ($n \geq 0$). (Sauf erreur. Y a-t-il des canulars, parce que le critère d'équivalence faible par les $\pi_i$ ne s'applique pas directement, en l'absence de points bases. Ça ne doit pas être très sérieux \ldots\ \texttt{[cela semble en fait sérieux]}). Mais est-il peut-être même \textit{pleinement fidèle} (\footnote[10]{\texttt{[un ?\ \ dans la marge]}})? De sorte que Hot se réaliserait comme sous-catégorie pleine de $\Hot_0^{\wedge}$, contenant les foncteurs représentables?


\dgrblock{dgrb191}{Finite unbased tests miss a symmetric-group shift}
\begin{dcomment}{examples}{p52:puisque un homomorphisme}{Finite unbased tests miss a symmetric-group shift}
Grothendieck's final question concerns a specific restricted Yoneda functor, and here a concrete counterexample settles both the conservativity assertion and the stronger full-faithfulness question. The mechanism is that every finite test sees a conjugating element, but no one element works globally. The finite-support symmetric-group example is the classical case recorded in \citet[Proposition~4.1]{MatumotoMinamiSugawara1984}; the attribution and its relation to the later general result are confirmed by \citet[Remark~2.3]{ArlinChristensen2023}. We give the elementary verification in the exact form needed for these pages.

Let $G$ be the group of permutations of $\mathbf N=\{0,1,\ldots\}$ with finite support. Define a homomorphism $s:G\to G$ by
\[
 s(\sigma)(0)=0,\qquad s(\sigma)(n+1)=\sigma(n)+1.
\]
It is injective but is not surjective, since every image permutation fixes $0$, whereas the transposition $(0\ 1)$ belongs to $G$. Therefore $Bs:BG\to BG$ is not a weak homotopy equivalence: on fundamental groups it is $s$.

Now let $P$ be any finitely generated group and $\rho:P\to G$ a homomorphism. Choose a finite generating set. The union of the supports of the images of those generators is finite, so every element of $\rho(P)$ fixes all integers greater than some $m$. Let $c=(0\ 1\ \cdots\ m+1)$. Conjugation by $c$ carries the action on $\{0,\ldots,m\}$ to the shifted action on $\{1,\ldots,m+1\}$ and fixes $0$. Directly, for every $p\in P$,
\[
 c\rho(p)c^{-1}=s(\rho(p)).
\]
Hence $s\rho$ and $\rho$ are conjugate homomorphisms.

For a connected CW complex $T$, homotopy classes of maps to $BG=K(G,1)$ correspond to homomorphisms $\pi_1T\to G$ modulo conjugation. One way to verify this description is to choose a vertex, extend a chosen group homomorphism across the $2$-skeleton, and then use the vanishing higher homotopy groups of $BG$ for existence and uniqueness of further extensions. Forgetting the chosen basepoint changes the homomorphism by conjugation. If $T$ is the realization of the nerve of a finite category, each component has finitely presented fundamental group: its finite $2$-skeleton gives generators and relations. Applying the preceding conjugacy calculation on each component therefore shows
\[
 (Bs)_*:[T,BG]\longrightarrow[T,BG]
       \quad\hbox{is the identity}
\]
for every $T\in\Hot_0$. The same holds for every sphere. This proves that the restricted Yoneda functor is not conservative.

It is also not faithful. The maps $1_{BG}$ and $Bs$ are distinct in $\Hot$, since the latter is not an equivalence and the former is. Nevertheless they induce the same natural endomorphism of the presheaf $T\mapsto[T,BG]$ on $\Hot_0$: both are the identity at each test object. Thus the proposed full faithfulness fails even at injectivity on morphisms. The two editorial warnings at the end of the chapter are justified; the missing basepoint information cannot be recovered merely by adding more finite unbased tests.
\end{dcomment}


\dgrblock{dgrb192}{The modern scope of the final obstruction}
\begin{dcomment}{status}{p52}{The modern scope of the final obstruction}
The counterexample above addresses exactly the finite-category tests in the manuscript. A much stronger theorem is now available: \citet[Theorem~2.1]{ArlinChristensen2023} prove that no set of objects of the unpointed homotopy category jointly reflects isomorphisms. Their proof enlarges the symmetric-group construction to a sufficiently large regular cardinal. Consequently the obstruction is not removed by choosing a different small family of unbased tests. The same paper distinguishes this theorem from an earlier argument of Heller whose asserted boundedness property had a gap.

The loss occurs at the passage to sets of homotopy classes. Retaining more structure changes the answer: \citet[Theorem~4.3]{ArlinChristensen2023} show that spheres jointly detect equivalences in the homotopy $2$-category of spaces, whose Hom-categories are fundamental groupoids of mapping spaces. Already the automorphisms in the mapping groupoid out of a point retain the fundamental group which a set of unbased components forgets. At the infinity-categorical level, mapping spaces provide the full Yoneda data. These results give a precise modern reason for the insistence, throughout these chapters, on keeping coherent diagrams and their transformations together.
\end{dcomment}
\begin{otherlanguage}{english}
\dgroverview{III-V}{Chapter V \textperiodcentered\ \textup{III-V} \textperiodcentered\ Overview}
These chapters study how a homotopy theory can retain operations that
ordinary homotopy classes conceal. They approach the problem from three directions.
Chapter III asks for an action of homotopy types on coefficients and for an
external Hom that represents the resulting families of morphisms. Chapter IV
constructs coherent diagrams by extending coefficients along carefully chosen
inclusions. Chapter V returns to the elementary presentation of a localization
by zigzags, then asks how much of its homotopy theory those zigzags remember.
The order matters. The external Hom is a desired object; coherent diagrams
provide a way to construct operations on it; the categories of paths are a
candidate model. What is established at each stage must be distinguished from
what the next stage would require.

These chapters have three independent manuscript numberings. Chapter III
contains manuscript pages 1--12, Chapter IV contains 73--83 (the edition
explicitly says that pages 1--72 are absent), and Chapter V contains 1--52,
with page 32 marked blank. We therefore cite manuscript locations with a
chapter numeral. In the edition, the mathematical text occupies forty-seven pages, besides
titles and editorial notices.
The inserted comments preserve the edition's corrections in square brackets.
The corrections to the transcription restore what that edition actually prints;
mathematical repairs are stated in the commentary instead of being silently
substituted for the historical text.

Chapter III begins with a concrete requirement. If $p:X\to e$ and $M$ is a
coefficient, a tensor by the homotopy type of $X$ should be $p_!p^*M$, and a
cotensor should be $p_*p^*M$. Their adjunction follows directly from the two
adjunctions involving $p^*$:
\[
 \Hom(p_!p^*M,N)\cong\Hom_{\D(X)}(p^*M,p^*N)
 \cong\Hom(M,p_*p^*N).
\]
The issue is descent in $X$: a weak homotopy equivalence of indexing categories
must induce an isomorphism on constant-coefficient homology and cohomology.
There is also a size issue. A derivator defined on a restricted class of
shapes cannot be presumed to see every homotopy type. These are separate
questions, and neither is settled by the formal adjunction alone.

The modern action of spaces on derivators gives a precise setting for the
first construction, while model presentations provide derived mapping spaces.
The distinction between an action and a closed enrichment remains essential:
the former is not, by itself, a representability theorem for the external
Hom. We relate the discussion to the universal property of the derivator of
spaces and to the later account of homotopy-theoretic module structures
\citep{Cisinski2008,GrothShulman2017}. For familiar derived categories one can
see all the operations explicitly. With cohomological grading,
\[
 K\otimes M=C_*(K;\mathbf Z)\otimes^{\mathbf L}_{\mathbf Z}M,
 \qquad
 M^K=\mathbf R\Hom_{\mathbf Z}(C_*(K;\mathbf Z),M).
\]
The mapping space from $M$ to $N$ is obtained from the good truncation
$\tau_{\leq0}\mathbf R\Hom_k(M,N)$ by the Dold--Kan construction and
forgetting the abelian structure. Thus its $n$th homotopy group at zero is
$\Hom_{D(k)}(M,N[-n])$. A positive Ext group appears as a higher homotopy
group after the corresponding shift of the target; it is not a positive
homotopy group of this unshifted mapping space merely because it is an Ext
group.

The circle calculation explains why pointedness enters the manuscript.
Unreduced chains on $S^1$ give $\mathbf Z\oplus\mathbf Z[1]$, hence
$S^1\otimes M\simeq M\oplus M[1]$ in the derived category. The suspension
is the reduced summand. Grothendieck's first identification of the tensor
with suspension is therefore corrected by his own closing observation; it
should not be reported as an unresolved defect of his final construction.
His suggestion of a bounded-above enhancement, however, requires a boundedness
hypothesis on derived Hom. An elementary example over the dual numbers shows
why arbitrary coefficients do not supply one.

Chapter IV explains why these operations must be constructed on coherent
diagrams. The underlying-diagram functor is conservative, but it discards
information about coherent compositions and homotopies. The commentary on Chapter~I
already examines this failure in detail. Here the new point is constructive:
Kan extension along a suitable sequence of inclusions supplies the coherence
that a diagram drawn in the fundamental category does not specify.
The loop object is obtained by extending a single coefficient by zero to
a punctured square and then taking its right Kan extension to the whole
square. Full faithfulness of these extensions explains both the existence
and the correctly qualified uniqueness of the resulting cartesian square.
This is the pattern developed systematically for pointed and stable
derivators in \citet{Groth2013}.

The same method makes the adjunction maps visible. For $f:I\to J$,
Grothendieck forms the collage $K(f)$ with mixed morphisms
$\Hom_{K(f)}(i,j)=\Hom_J(f(i),j)$. A presheaf on $K(f)$ is exactly a
pair of presheaves together with a map $f^*G\to F$. The elementary
description identifies the required units and counits before the derived
case is attempted. We complete the construction by giving the indexing
functors and checking the pointwise Kan-extension comparison. The benefit
is a canonical coherent realization of these particular adjunction maps.
It does not assert a canonical lift of every arrow drawn in a homotopy
category.

Chapter V begins with alternating finite paths, extended by identities at
their two ends. A backward segment is required to belong to $\Sigma$;
the connected components of the resulting category with fixed endpoints
then give the morphisms of $\mathcal M[\Sigma^{-1}]$. Here ``fixed'' includes
the morphisms: transformations in the fibre have identity components at
the two endpoints. Replacing this fibre by the full subcategory on paths
with those endpoint objects would change the construction. We give the
quotient and composition argument with this convention throughout.

The manuscript's parity correction is also substantive. Concatenating two
of the initially chosen odd-indexed standard paths reverses the standard
orientation. Passing to the cofinal subsystem of even indices repairs the
problem. Grothendieck's analysis of admissible contractions establishes the
necessary compatibility with padding. His more general interval lemma is
left as an assumption, with contractibility proposed as a strengthening.
The development supplied here proves that strengthening: for a finite
oriented path as target, every nonempty category of functors with a prescribed
restriction is categorically contractible. Numerical minimum and maximum
give explicit endofunctors and a zigzag between the identity and a constant
functor. The fixed-restriction version also covers the cancellation diagrams
whose two endpoints have the same image. We give the full argument and state
the extent of the literature search; the label \emph{ours} records responsibility
for this proof, not a claim of priority.

Passing to diagrams exposes a uniformity condition that a bare union of
path categories forgets. A path in $\mathcal M^X$ has one finite length
bound valid at every object of $X$. Consequently the natural representing
object is an ind-category $(K_n)$, with
\[
 \Hom_{\mathcal M^X[\Sigma_X^{-1}]}(a_X,b_X)
   \cong\mathop{\rm colim}_n\pi_0\operatorname{Fun}(X,K_n).
\]
This is an unconditional statement about ordinary localizations. It is
not yet an identification of $NK_\infty$ with a derived mapping space.
The later simplicial-localization theory supplies canonical mapping spaces
and comparison theorems under their stated hypotheses; the commentary keeps
those theorems separate from what this particular path model proves.

The formal propositions on reflective localizations and their ind-categories
can be preceded by a simpler affirmative result. If $H_{a,b}$ inverts nerve
weak equivalences, the image of its representing ind-object already
ind-represents the descended functor on $\Hot$. Indeed, left Kan extension
along the opposite of a localization carries representable presheaves to
representable presheaves and preserves filtered colimits; restriction along
the localization is fully faithful. This conclusion needs no right adjoint
between the original categories.

The stronger formal propositions on reflective localizations and their ind-categories
are sound when the fully faithful right adjoint is supplied. Their application
to $\mathrm{Cat\,hot}\to\Hot$ needs that additional premise. Here
$\mathrm{Cat\,hot}$ means categories and functors modulo finite zigzags of
natural transformations, whereas $\Hot$ inverts weak equivalences of nerves.
The editorial request for a proof of the right adjoint is consequently
material. Thomason's model structure does not alone establish the asserted
adjunction with this particular source category. We leave that premise
unresolved in this audit and explain the valid replacement-based conclusions
without promoting an unsuccessful search to a claim that the problem is open.

The last page admits a definite answer. Testing an unpointed homotopy type
by every finite category does not give a conservative or a faithful restricted
Yoneda functor. Let $G$ be the group of permutations of $\mathbf N$ with finite
support, and let $s:G\to G$ shift a permutation one place to the right while
fixing zero. The induced $Bs$ is not an equivalence. Every map from a finite
category's nerve into $BG$, however, uses a finitely generated fundamental
group and hence only finitely many permutation supports. On those supports
the shift is conjugation by one finite permutation. Thus $Bs$ induces the
identity on all the proposed tests. The same example separates $Bs$ from
the identity despite their identical restricted Yoneda transformations.
The modern general theorem says that no set of unpointed spaces jointly
reflects isomorphisms in $\Hot$ \citep[Theorem~2.1 and Remark~2.3]{ArlinChristensen2023}.
Retaining mapping groupoids changes this conclusion: the sphere tests do
reflect equivalences in the homotopy 2-category, as their Theorem~4.3 proves.
The obstruction on the closing page is precisely information discarded by
passing too early to sets of unpointed homotopy classes.
\end{otherlanguage}
\dgrchunkend
